\documentclass[11pt,a4paper]{article}
\usepackage[T1]{fontenc}
\usepackage[utf8]{inputenc}
\usepackage{lmodern,amsmath,amssymb,amsthm,mathtools,microtype,mathrsfs}
\usepackage[margin=25mm,top=23mm,bottom=25mm]{geometry}
\usepackage{booktabs,array,longtable,enumitem,xcolor,fancyhdr}
\usepackage[colorlinks=true,linkcolor=blue!45!black,urlcolor=blue!45!black,citecolor=blue!45!black]{hyperref}
\hypersetup{pdftitle={Reading the Alpoge--Claude RH and Dirichlet-GRH proof: working version 0.50},pdfauthor={ChatGPT},pdfsubject={Partial human-readable reconstruction with explicit arithmetic inputs},pdfkeywords={Riemann hypothesis, Lean, theta lifts, prime sums, Arakelov slopes}}
\setlength{\parindent}{0pt}
\setlength{\parskip}{5pt plus 1pt}
\setlist{topsep=3pt,itemsep=3pt,leftmargin=*}
\setlength{\headheight}{14pt}
\pagestyle{fancy}\fancyhf{}
\fancyhead[L]{\small RH and Dirichlet GRH: a reconstruction}
\fancyhead[R]{\small Working version 0.50}
\fancyfoot[C]{\thepage}
\renewcommand{\headrulewidth}{0.3pt}
\numberwithin{equation}{section}
\newtheorem{theorem}{Theorem}[section]
\newtheorem{proposition}[theorem]{Proposition}
\newtheorem{lemma}[theorem]{Lemma}
\newtheorem{corollary}[theorem]{Corollary}
\theoremstyle{remark}\newtheorem{remark}[theorem]{Remark}
\newcommand{\Rea}{\operatorname{Re}}
\newcommand{\HS}{\mathrm{HS}}
\newcommand{\tr}{\operatorname{tr}}
\newcommand{\codim}{\operatorname{codim}}
\newcommand{\covol}{\operatorname{covol}}
\newcommand{\OK}{\mathcal O_K}
\newcommand{\src}[2]{\href{https://github.com/CocoSgt/riemann-hypothesis-lean/blob/7cbf7a0c32fca1badf05931879a37c3f8a2b03b4/src/#1}{\texttt{#2}}}
\newcommand{\status}[1]{\textbf{Status:} #1}
\begin{document}
\begin{center}
{\LARGE\bfseries RH and Dirichlet GRH: new human proofs}\\[7pt]
{\large Version 0.50 --- 23 November 2026}\\[5pt]
{\large Frames, Bessel periods and denominator bounds}
\end{center}
\textbf{Scope.} This edition contains the current audit and all new chapters.
It is a partial human reconstruction of the accepted formal proof.
Earlier named claims are in the cumulative manuscript and preserved components.
Each result below states its original inputs and exact remaining producers.
\setcounter{tocdepth}{1}
\tableofcontents
\clearpage
\begingroup\emergencystretch=3em
\section{Version 0.50: the current proof frontier}
\label{sec:dependency-audit-v50}

\subsection{The objective and the exact advance}

We accept the supplied Alpoge--Claude Lean proof as correct and reconstruct
its mathematics in ordinary prose. The task here is to write the human
argument, with simplifications whenever they preserve the needed conclusion.
The pinned source commit remains
\texttt{7cbf7a0c32fca1badf05931879a37c3f8a2b03b4}.
The version-0.49 manuscript and all 3014 earlier claim records are
preserved. Seven new chapters add 98 independently reviewed claims,
bringing the written inventory to 3112.

The new results make substantial progress inside the existing obligations.
They close no additional \emph{whole} obligation in the selected sufficient
route. The graph still has 121 nodes and 161 edges, with 24 open leaves:
17 class C and 7 class B. Exactly 17 are shared by ordinary RH and
unrestricted Dirichlet GRH; 7 are additional Dirichlet-GRH inputs.
These counts measure the inventory of the chosen decomposition. They do
not measure a percentage of the mathematical work or assert that the
decomposition is minimal.

\subsection{The new chapters and their mathematical content}

\paragraph{TDR1--TDR13: actual framed points descend to finite data.}
The Lie eigenbasis and theta-section generator are represented by their
actual frame bundles over the given moduli scheme. Finite presentation
then descends every original framed point to a finitely generated
coefficient subalgebra, without assuming that the target algebra is
reduced, Noetherian or flat. For a reduced target, taking a radical
quotient of a finite polynomial ring gives the literal test ring and
the original \texttt{Fr\allowbreak Map}. This proves the whole
\texttt{all\allowbreak Test\_\allowbreak of\_\allowbreak sec\allowbreak Base\allowbreak Change} implication and its
\texttt{From\allowbreak Test} return. The exact section-base-change premise and
the supplied moduli scheme are retained.

\paragraph{NFR1--NFR13: normalized theta points satisfy every single-frame condition.}
The proof supplies the original embedding and column dictionary,
holomorphic periods, connected frame domain and symplectic basis with
its exact modulus-eight adaptation. The Fourier recurrence gives a
direct one-dimensionality proof for the actual complex theta-section
line. It is enough here without an additional section-twisting premise.
The actual Hermitian form, unit-minor determinant, theta mean and
quadratic gauge have their precise constants. The source's square-root
normalization makes the Gaussian mean exactly one. These calculations
prove the complete \texttt{Theta\allowbreak Data}, \texttt{Is\allowbreak Norm\allowbreak Un} and
\texttt{Frame\allowbreak OK} results. Simultaneous choice gives the first
\texttt{Frames\allowbreak OK} field; its other three fields, \texttt{Cover\allowbreak Tri},
\texttt{Fib\allowbreak Tri} and full \texttt{Is\allowbreak Unf} remain to be supplied.

\paragraph{CAF1--CAF11: one character seed before every later field and prime.}
The exact local norm, Teichmuller and conductor argument constructs the
original seed character uniformly. An elementary cyclotomic argument
supplies the auxiliary primes, and the actual fixed field has the
required CM and maximal-real subgroup identities. The original
\texttt{Dat} modulus and the independent \texttt{Count\allowbreak V Host\allowbreak Data}
modulus are kept distinct. In particular, the latter retains the
allowed case $\ell=3$. These are character seeds; the rank-two and
rank-three automorphic representations remain a later construction.

\paragraph{MFC1--MFC14: the common field and marked frame are actual objects.}
An elementary CRT argument constructs the real quadratic auxiliary field
with the required local behavior. The proof supplies every original
\texttt{UFacts} field, the class-rank bound chosen before the tower
parameter, and the complete literal \texttt{link\_\allowbreak ex}. Both original
characters satisfy strictness and belong to their prescribed cone
classes. The resulting common $L,g_1,\Phi,C_0,w$ satisfies
\texttt{Big} and 21 of the 28 \texttt{Ok} fields. When $O=c.O$ is
the original coefficient DVR, the CAF seed joins this same frame,
with its modulus chosen first.
The seven fields \texttt{phi/\allowbreak gen/\allowbreak cl/\allowbreak ec/\allowbreak b1p/\allowbreak b1/\allowbreak h1} and the full entry
implications \texttt{Gen\_\allowbreak gd}, \texttt{Gen\_\allowbreak host} remain. Thus
the complete common automorphy endpoint \texttt{Gd\allowbreak End} is still a
producer, with the same marked $g_1$ and coefficient embedding required
throughout.

\paragraph{BNV1--BNV23: a nonzero base integral and its exact universal return.}
The original integration exponent and full finite parameter walk are
proved. An explicit isotropic frame and Cayley transform match the
source character, including its factor $-2it$. The Haar exponential
chart, a global Cayley majorant and Casimir estimates give a uniform
tangent Gaussian and a nonzero complex Gaussian integral. The
dominating bound is chosen before all integration variables; the
original almost-everywhere strong measurability is sufficient. Schur
averaging has exactly the factor $1/\dim V_\pi$ and retains the full
allowed ambient second-measure scope. Data transport and compact
cyclicity return nonvanishing to the original supplied endomorphism
integral. The whole \texttt{Step\allowbreak Statement} is still needed to
propagate from the large-parameter base to every target.

\paragraph{BSP1--BSP10: the analytic end of weight propagation.}
The complete original \texttt{Gap\allowbreak Need}, \texttt{Psi\allowbreak Datum},
\texttt{Cap\allowbreak Val}, \texttt{Mult}, \texttt{Per} and
\texttt{Fin} statements are supplied at their stated premises.
The Capelli polynomial acts by its exact scalar on the entire generated
module; differentiation commutes with the original vector period.
Coefficient orthogonality yields pointwise reproduction on the actual
right-regular space. For $\rho$ the Frobenius norm, $\beta>0$, and
group elements $y,g$, the useful simpler majorant is
\[
 \rho(y)^{-\beta}\rho(y^{-1}g)^{-\beta}
 \le \rho(g)^{-\beta/4}\rho(y)^{-\beta/2}.
\]
The exact source exponents make its factors separately integrable.
This proves the final Fubini and nonvanishing extraction. The standard
and dual tensor coefficient maps, complementary-minor eigenspaces and
evaluations, and shifted quotient constructions remain the suppliers
of the original two-sided step.

\paragraph{DNE1--DNE14: the whole conditional global denominator estimate.}
The original patched table, including its prime-three equal-share
branch, supplies all seventeen admissibility fields. Explicit
cyclotomic, small-order, band, exceptional and ordinary-prime estimates
control all the original sums. Uniform exponential comparisons replace
a numerical cutoff. The final constant is
\[
                    K=2000(|C_a|+1)^3,
\]
chosen before the family, with the original eventual-threshold order.
The exact implication is
\[
 \texttt{Concl\allowbreak F}\ \Longrightarrow\ \texttt{Adm\allowbreak F}
                  \ \Longrightarrow\ \texttt{Link\allowbreak N}.
\]
The original \texttt{Sha\allowbreak Fin\allowbreak Odd} and summability guards remain.
Constructing an actual common \texttt{Concl\allowbreak F} family is still required
for the whole selected denominator obligation.

\subsection{How the completed pieces connect}

The two geometry chapters feed the simultaneous principal-frame
construction. The two arithmetic chapters feed the common marked-frame
and automorphy construction. The Bessel base, finite walk and analytic
return reduce the remaining nonvanishing work to the explicit
two-sided step suppliers. Completing that step yields the original
\texttt{Bessel\allowbreak NVGap}; the inherited wave-packet construction then
supplies its setting conclusion at the original \texttt{Forms0} scope.
The denominator chapter consumes the local-bound family when that
actual family is available. The already completed v0.49 local-gain
estimate is preserved and does not need to be reproved.

These are deductions on the same original objects. Existence of a
compatible-looking field, unrelated theta model, abstract representation
or separate row does not replace the original common field, frame,
section, quotient or family. Their full producer premises are printed
where they are used.

\subsection{Current whole-obligation register}

Each whole contract below is preserved verbatim from the previous
selected graph. Historical qualifications embedded in those contracts
remain for traceability. The following version-0.50 progress paragraph,
where present, gives the current completed subresults and remaining
scope. A completed conditional implication does not erase its premises.
The ledger records the human dependency graph; it is not a mechanically
extracted transitive Lean dependency graph. Every new proof has a
separate independent mathematical review and exact source/proof hashes.
Draft TeX validation checks syntax and references; no PDF visual
certification or new formal replay is asserted.

\paragraph{1. GLOBAL\_\allowbreak DENOMINATOR — C — Actual LinkN local table and binding producer}

Closed(Des $\mathcal C$ 3) LinkN: the actual etaSum is bounded with exact dZ0 and same-host vvTerm, subject only to the literal ShaFinOdd and summability guards.

\emph{Shared objects and choice order.} one FamP; one class $\mathcal A$ and its target-compatible subclass; one coefficient field per family datum; one x and i at each application; original selected rows, words and levels; thresholds chosen before the common family; actual D.nd; actual etaJ and Aj; sha\_\allowbreak eq

\emph{Completed inputs (version 0.50).} DNE1-DNE14 (whole original ConclF -> AdmF -> LinkN consumer), Inherited original local-bound and length interfaces at their supplied-data scope

\emph{Remaining producer scope.} Construct one actual common family satisfying the original ConclF, including all its local table bounds, actual rows, arithmetic data and the original ShaFinOdd and vvTerm-summability guards. DNE now proves the full ConclF-to-AdmF-to-LinkN consumer and every original export quantifier, with K=2000(|Ca|+1)\textasciicircum{}3 chosen before the family and the eventual threshold in the original order. The prime estimates and analytic scale comparison need not be reopened. The existential common-family and ConclF producer remains, so GLOBAL\_\allowbreak DENOMINATOR stays open.

\emph{Source.} \href{https://github.com/CocoSgt/riemann-hypothesis-lean/blob/7cbf7a0c32fca1badf05931879a37c3f8a2b03b4/src/Formalism/SharpRestriction/GlobalTheoremInputs.lean#L45-L48}{LinkN}.

\emph{Source.} \href{https://github.com/CocoSgt/riemann-hypothesis-lean/blob/7cbf7a0c32fca1badf05931879a37c3f8a2b03b4/src/Formalism/SharpRestriction/GlobalTheoremInputs.lean#L40-L44}{LinkN\_\allowbreak at}.

\emph{Source.} \href{https://github.com/CocoSgt/riemann-hypothesis-lean/blob/7cbf7a0c32fca1badf05931879a37c3f8a2b03b4/src/Development/AutomorphicArithmetic/CharacterBoxes/DivinNOfHypV.lean#L117-L134}{divinN\_\allowbreak of\_\allowbreak hypV}.



\paragraph{2. GLOBAL\_\allowbreak RETAINED\_\allowbreak DEFECT — C — Original retained-row defect bound LinkIRR}

Rows.DelpNonneg and delp$\le$exp(-x\textasciicircum{}(ez/7)) uniformly for every late actual row datum. CURRENT VERSION0.47 QUALIFICATION: Finish the same-D/common-family adelic setting after ASG: normalized nonzero Bessel W and decay, measured-field block, original Siegel-Weil rescaling and measure/character ties, Bessel uniqueness/local ties, then the original AutTie/RowDictTie/LabTie. RowsFinDim is now proved universally by RFD. Supply the actual CountIdempotentsClass theorem or its three sufficient arithmetic cylinder counts, with one simultaneous family threshold; reuse the completed ARO structure. Completed new original producers: ARO1-ARO30, RFD1-RFD9, ASG1-ASG18. Any earlier narrative naming these completed producers as outstanding is historical; this qualification is the current scope authority. CURRENT VERSION0.48 QUALIFICATION: Finish the same-D/common-family adelic setting after ASG and AMD: normalized nonzero Bessel W and decay, original orbital/fiber and Siegel-Weil measure/character ties, Bessel uniqueness/local ties, and the full original transport theorem, then AutTie/RowDictTie/LabTie. The complete MExists measured-data block and unipotent covolume normalization are now supplied, and universal RowsFinDim remains available. Supply the actual CountIdempotentsClass theorem or its three sufficient arithmetic cylinder counts with one simultaneous family threshold. Completed current producers or conditional consumers: ARO1-ARO30, RFD1-RFD9, ASG1-ASG18, AMD1-AMD17. Earlier frontier descriptions are historical; this qualification is the current scope authority. Original whole-contract premises and quantifiers are retained. CURRENT VERSION0.49 QUALIFICATION: Prove the exact original BesselNVGap nonzero integral on its supplied discrete-series, compact-representation and character data. CoefWeightedLoc, WlocExists with its original nonvanishing antecedent, WFromLoc, the normalized global product and all-later-measured-block decay are complete. On the inherited Forms0 scope this leaves BesselNVGap for ConcreteSettingExists. Continue original orbital/fiber and Siegel-Weil measure/character ties, the required Bessel uniqueness/local ties and the full original transport theorem, then AutTie/RowDictTie/LabTie. Supply the actual CountIdempotentsClass theorem or its three sufficient arithmetic cylinder counts with one simultaneous family threshold. Completed current producers or conditional consumers: ARO1-ARO30, RFD1-RFD9, ASG1-ASG18, AMD1-AMD17, BWP1-BWP24 at their exact BesselNVGap/Forms0 scope. Earlier frontier descriptions are historical; this qualification is the current scope authority. Original whole-contract premises and quantifiers are retained.

\emph{Shared objects and choice order.} one FamP; one class $\mathcal A$ and its target-compatible subclass; one coefficient field per family datum; one x and i at each application; original selected rows, words and levels; thresholds chosen before the common family; AllRows A; LabelCountGoodInjective; AllJoinA; CountAllRows

\emph{Completed inputs (version 0.50).} DNE1-DNE14 separate denominator consumer, WRS1-WRS10 full selected-class local gain, Inherited first DivDef3 and row-pair integrability

\emph{Remaining producer scope.} Complete the original F3 local coefficient bridges and F4 common-overfield readings, with valuation/norm/GoodCoef/iq and the entire coefficient field, and actual complex Hecke standard/no-Jordan plus arithmetic character data on the same form datum. Construct the actual base rows, host/frame lengths, nonzero chains and group ties, MultEq support, eigensystem/carrier/parts, all-embedding logLam, HD/d3b and NEb. Keep the stronger second DivRows3 scope distinct from the completed first DivDef3 and row-pair integrability. The partial common marked frame and CountV seed do not by themselves install HOST\_\allowbreak NUMBER\_\allowbreak S, HOST\_\allowbreak FRAME\_\allowbreak LENGTH, TASK\_\allowbreak HOST\_\allowbreak VAN\_\allowbreak E\_\allowbreak INSTALL, TASK\_\allowbreak HOST\_\allowbreak XPP\_\allowbreak INSTALL or TASK\_\allowbreak COUNT\_\allowbreak V\_\allowbreak INSTALL. Continue the distinct retained-defect, period-height and normalized spectral producers, and the original compatible second-layer selection on one common arithmetic/geometric family. GLOBAL\_\allowbreak LOCAL\_\allowbreak GAIN remains closed at its complete v0.49 scope; the current denominator work only completes its separate consumer. Common design existence, NBUILT compatibility, original true-column records, shared fields/embeddings, Fits and eta agreement and eventual-family thresholds remain explicit. A stronger historical WPt/TSel construction may still serve other producers but is not needed to reopen the proved selected-class local-gain estimate.

\emph{Source.} \href{https://github.com/CocoSgt/riemann-hypothesis-lean/blob/7cbf7a0c32fca1badf05931879a37c3f8a2b03b4/src/Formalism/SharpRestriction/GlobalTheoremInputs.lean#L55-L58}{LinkIRR}.

\emph{Source.} \href{https://github.com/CocoSgt/riemann-hypothesis-lean/blob/7cbf7a0c32fca1badf05931879a37c3f8a2b03b4/src/Development/AbelianVarietyTheory/Proofs/GlobalInterpolation/LinkIRROfJoinEx.lean#L17-L38}{linkIRR\_\allowbreak of\_\allowbreak join\_\allowbreak ex}.

\emph{Source.} \href{https://github.com/CocoSgt/riemann-hypothesis-lean/blob/7cbf7a0c32fca1badf05931879a37c3f8a2b03b4/src/Development/AbelianVarietyTheory/Proofs/GlobalInterpolation/LinkIRROfJoinEx.lean#L143-L166}{link\_\allowbreak irr}.



\paragraph{3. GLOBAL\_\allowbreak PERIOD\_\allowbreak HEIGHT — C — Full actual period and height comparison LinkDIVt}

Closed conditional TuSize -> DGamSize -> (LinkEta -> DIVinN -> LinkDIV), with every original LinkDIV\_\allowbreak at clause.

\emph{Shared objects and choice order.} one FamP; one class $\mathcal A$ and its target-compatible subclass; one coefficient field per family datum; one x and i at each application; original selected rows, words and levels; thresholds chosen before the common family; same RR target; same Bu/besS; same Nu and Bk

\emph{Source.} \href{https://github.com/CocoSgt/riemann-hypothesis-lean/blob/7cbf7a0c32fca1badf05931879a37c3f8a2b03b4/src/Formalism/SharpRestriction/GlobalTheoremInputs.lean#L77-L90}{LinkDIV\_\allowbreak at}.

\emph{Source.} \href{https://github.com/CocoSgt/riemann-hypothesis-lean/blob/7cbf7a0c32fca1badf05931879a37c3f8a2b03b4/src/Formalism/SharpRestriction/GlobalTheoremInputs.lean#L94-L96}{LinkDIVt}.

\emph{Source.} \href{https://github.com/CocoSgt/riemann-hypothesis-lean/blob/7cbf7a0c32fca1badf05931879a37c3f8a2b03b4/src/Development/AbelianVarietyTheory/Proofs/GlobalInterpolation/LinkBrest.lean#L7-L25}{GlobalTheory.Top.link\_\allowbreak brest}.

\emph{Source.} \href{https://github.com/CocoSgt/riemann-hypothesis-lean/blob/7cbf7a0c32fca1badf05931879a37c3f8a2b03b4/src/Development/AbelianVarietyTheory/Proofs/GlobalInterpolation/WireSel.lean#L11-L41}{GlobalTheory.Top.wire\_\allowbreak sel}.



\paragraph{4. GLOBAL\_\allowbreak SPECTRAL — C — Original normalized kernel/Gram estimates LinkS}

Closed TuSize -> LinkS with exact Gamma normalization, off-diagonal Hilbert-Schmidt bound, diagonal error and Dall/Dstar estimate, errors exp(-x\textasciicircum{}(epsilon/8)). CURRENT VERSION0.48 QUALIFICATION: Complete the original same-family spectral/Bessel and Siegel-Weil inputs with all measure, character, transport and common-object ties. The complete MExists measured block and prescribed-measure covolume normalization are supplied; these are not the full normalized spectral theorem. Completed current producers or conditional consumers: ASG1-ASG18, AMD1-AMD17, GCF1-GCF18 at their stated scopes. Earlier frontier descriptions are historical; this qualification is the current scope authority. Original whole-contract premises and quantifiers are retained. CURRENT VERSION0.49 QUALIFICATION: Complete the original same-family spectral/Bessel and Siegel-Weil inputs with every measure, character, transport and common-object tie. The normalized packet and decay are now constructed conditional on the exact BesselNVGap, and the analytic setting uses the original supplied Forms0 scope. These do not imply the full normalized spectral theorem. Completed current producers or conditional consumers: ASG1-ASG18, AMD1-AMD17, GCF1-GCF18 at their stated scopes, BWP1-BWP24 at their exact BesselNVGap/Forms0 scope, WRS1-WRS10 original local gain. Earlier frontier descriptions are historical; this qualification is the current scope authority. Original whole-contract premises and quantifiers are retained.

\emph{Shared objects and choice order.} one FamP; one class $\mathcal A$ and its target-compatible subclass; one coefficient field per family datum; one x and i at each application; original selected rows, words and levels; thresholds chosen before the common family; actual KD; same S.Cu spaces; GamOf; Dstar; Dall

\emph{Completed inputs (version 0.50).} BNV1-BNV22 (IntS, base nonvanishing, walk and exact universal return), BNV23 conditional on original StepStatement, BSP1-BSP10 (full original analytic propagation suppliers), BWP1-BWP24 at original nonvanishing/Forms0 scope, AMD1-AMD24 at stated scope

\emph{Remaining producer scope.} Complete the original two-sided StepStatement. Standard StepUp at lambda+e\_\allowbreak s still needs TMapStatement, complementary-minor EigStatement, EvalPStatement and QuotPStatement; dual StepDown at lambda-e\_\allowbreak s needs TMapDStatement, EigStatement, EvalMStatement and QuotMStatement. BNV now supplies IntS, the large-parameter StepIBase, the full finite walk, data transport and universal return; BSP supplies GapNeed, PsiDatum, CapVal, Mult, Per and Fin at their full original premises. Once the remaining suppliers are proved on their actual modules, periods and quotients, BNV23 yields BesselNVGap and the inherited BWP setting result with Forms0 retained. The actual orbital/fiber and Siegel-Weil character/measure comparisons, Bessel uniqueness/local ties, full transport, AutTie/RowDictTie/LabTie, and the simultaneous arithmetic cylinder counts for CountIdempotentsClass remain separate.

\emph{Source.} \href{https://github.com/CocoSgt/riemann-hypothesis-lean/blob/7cbf7a0c32fca1badf05931879a37c3f8a2b03b4/src/Formalism/SharpRestriction/GlobalTheoremInputs.lean#L145-L149}{LinkS\_\allowbreak at}.

\emph{Source.} \href{https://github.com/CocoSgt/riemann-hypothesis-lean/blob/7cbf7a0c32fca1badf05931879a37c3f8a2b03b4/src/Formalism/SharpRestriction/GlobalTheoremInputs.lean#L152-L157}{LinkS}.

\emph{Source.} \href{https://github.com/CocoSgt/riemann-hypothesis-lean/blob/7cbf7a0c32fca1badf05931879a37c3f8a2b03b4/src/Development/AbelianVarietyTheory/Proofs/GlobalInterpolation/LinkAlt.lean#L16-L50}{GlobalTheory.Top.link\_\allowbreak alt}.



\paragraph{5. AUT\_\allowbreak UNRAM\_\allowbreak COEF — C — Actual unramified-coefficient automorphy}

For each actual Dat with coefficient ramification one and fixed iota, choose one marked g1 supplying FinA, AutB and LemW together. CURRENT VERSION0.48 QUALIFICATION: Construct the actual AutomorphyUnramifiedCoefficients package: for every original Dat with coefficient ramification one and fixed iota, choose one marked g1 that supplies FinA, AutB and LemW together. The original InpZ record, full RigII, CNR3, finite PropGK consumer and now the entire OneChar implication are supplied. Continue RH Cases I/V and the independent additional Dirichlet-GRH CaseIII/promodularity, SeedIV, tensor-factor and rank-four-chain tasks. Completed current producers or conditional consumers: OCH1-OCH30 (the full consumer implication), KIZ original InpZ, KTR1-KTR13, RNC16, PGK1-PGK6. Earlier frontier descriptions are historical; this qualification is the current scope authority. Original whole-contract premises and quantifiers are retained. CURRENT VERSION0.49 QUALIFICATION: Construct one common GdEnd for each original datum and fixed iota: the same marked g1 must supply FinEnd, AutEnd3 and AutEnd2. UCE now proves all ordinary local/curve/coefficient consumers and the exact implication to FinA/AutB/LemW, including arbitrary entry maps and every original outer-corner extension. The producer still needs the actual common abelian-CM GdFrameEx and LinkedEndE with both rank-two and rank-three runs, full ModUAlt/PfromAut and the two actual seeds, followed by the original local/automorphic descent. At this pin the alternative upstream route gives GdEndAlt; its additional local/descent fields must be matched to the actual original points, so no name-only GdEndAlt-to-GdEnd identification is allowed. Keep the supplied InpZ, RigII, CNR3, finite PropGK and whole OneChar implication. Continue RH Cases I/V and the independent additional Dirichlet-GRH CaseIII/promodularity, SeedIV, tensor-factor and rank-four-chain tasks. Completed current producers or conditional consumers: UCE1-UCE19 (complete GdEnd consumer), OCH1-OCH30, KIZ original InpZ, KTR1-KTR13, RNC16, PGK1-PGK6. Earlier frontier descriptions are historical; this qualification is the current scope authority. Original whole-contract premises and quantifiers are retained.

\emph{Completed inputs (version 0.50).} CAF1-CAF11 (uniform original character seed), MFC1-MFC14 (actual linked frame, 21 of 28 Ok fields), UCE1-UCE19 (full original GdEnd consumer), OCH1-OCH30, KIZ original InpZ, KTR1-KTR13, RNC16, PGK1-PGK6

\emph{Remaining producer scope.} Complete the seven original Ok fields phi/gen/cl/ec/b1p/b1/h1 in the common GC construction, and separately the full entry implications Gen\_\allowbreak gd and Gen\_\allowbreak host. CAF supplies the character SeedClause uniformly before all later compatible fields and primes; MFC supplies the actual linked common field, all twelve UFacts, the uniform b0 class-rank bound, both strict cone classes, Big and 21 of 28 Ok fields. In the original coefficient-DVR case it also joins the CAF seed, with its modulus chosen first. The same actual marked g1 and fixed coefficient embedding must then receive FinEnd, AutEnd3 and AutEnd2, the actual rank-two and rank-three automorphic seeds, full ModUAlt/PfromAut, and original local/automorphic descent to GdEnd. UCE supplies the full GdEnd consumer, not the producer. The alternative GdEndAlt route still requires its exact additional local/descent fields. Do not exclude ell=3 from the independent CountV HostData seed scope merely because the final GD.gen guard excludes it.

\emph{Source.} \href{https://github.com/CocoSgt/riemann-hypothesis-lean/blob/7cbf7a0c32fca1badf05931879a37c3f8a2b03b4/src/Development/AutomorphicArithmetic/AutomorphicVarieties/AutomorphicCases/AutomorphyUnramifiedCoefficients.lean#L24-L27}{AutomorphyUnramifiedCoefficients}.

\emph{Source.} \href{https://github.com/CocoSgt/riemann-hypothesis-lean/blob/7cbf7a0c32fca1badf05931879a37c3f8a2b03b4/src/Development/AutomorphicArithmetic/AutomorphicVarieties/AutomorphicCases/AutomorphyUnramifiedCoefficients.lean#L12-L13}{FinA}.

\emph{Source.} \href{https://github.com/CocoSgt/riemann-hypothesis-lean/blob/7cbf7a0c32fca1badf05931879a37c3f8a2b03b4/src/Development/AutomorphicArithmetic/AutomorphicVarieties/AutomorphicCases/AutomorphyUnramifiedCoefficients.lean#L15-L17}{AutB}.

\emph{Source.} \href{https://github.com/CocoSgt/riemann-hypothesis-lean/blob/7cbf7a0c32fca1badf05931879a37c3f8a2b03b4/src/Development/AutomorphicArithmetic/AutomorphicVarieties/AutomorphicCases/AutomorphyUnramifiedCoefficients.lean#L19-L22}{LemW}.



\paragraph{6. TASK\_\allowbreak CASE3\_\allowbreak AUT\_\allowbreak JOIN — B — Complete the original Case-III automorphy/image join}

Install the actual Case-III ProModAlt, SeedIVAlt and shared PMExAlt with the inherited CohomologyHyp/Ferrero-Washington/HFit/SeedII facts into the original whole Case-III automorphy/image conclusion.

\emph{Source.} \href{https://github.com/CocoSgt/riemann-hypothesis-lean/blob/7cbf7a0c32fca1badf05931879a37c3f8a2b03b4/src/Development/AutomorphicArithmetic/AutomorphicVarieties/ModularityArgument/ThresholdArguments/DenseThresholdCases/PromodularityCaseThree.lean#L43-L44}{AutomorphyRankFourDense}.

\emph{Source.} \href{https://github.com/CocoSgt/riemann-hypothesis-lean/blob/7cbf7a0c32fca1badf05931879a37c3f8a2b03b4/src/Development/AutomorphicArithmetic/AutomorphicVarieties/ModularityArgument/ThresholdArguments/DenseThresholdCases/CuFinal.lean#L44-L44}{main4\_\allowbreak alt}.



\paragraph{7. TASK\_\allowbreak CASE3\_\allowbreak PROMOD\_\allowbreak JOIN — C — Reconstruct the actual Case-III chain-to-promodularity proof}

Use the original tensor-factor tower, ChainRankFour, ChainDesc4, seeded rank-four chain reduction, actual PMExAlt and SeedIVAlt to construct ProModAlt on every original Ctx3 star point at the same fixed iota, preserving all finite-coefficient, inertia and support maps.

\emph{Source.} \href{https://github.com/CocoSgt/riemann-hypothesis-lean/blob/7cbf7a0c32fca1badf05931879a37c3f8a2b03b4/src/Development/AutomorphicArithmetic/AutomorphicVarieties/ModularityArgument/ThresholdArguments/DenseThresholdCases/PromodularityCaseThreeOfLemmas.lean#L131-L153}{UpperTheory.V.Cu.promodularity\_\allowbreak case\_\allowbreak three\_\allowbreak of\_\allowbreak chain\_\allowbreak reduction}.

\emph{Source.} \href{https://github.com/CocoSgt/riemann-hypothesis-lean/blob/7cbf7a0c32fca1badf05931879a37c3f8a2b03b4/src/Development/AutomorphicArithmetic/AutomorphicVarieties/ModularityArgument/ThresholdArguments/DenseThresholdCases/PromodularityCaseThree.lean#L40-L42}{PromodularityCaseThreeDense}.



\paragraph{8. SEED\_\allowbreak IV\_\allowbreak ALT — C — Actual Case-III seed with prechosen residual datum}

Base Ctx3/Four tower, iota and Amb -> w,s before every later Dies tower -> compatible finite-coefficient seed with exact residual support/root entries.

\emph{Source.} \href{https://github.com/CocoSgt/riemann-hypothesis-lean/blob/7cbf7a0c32fca1badf05931879a37c3f8a2b03b4/src/Development/AutomorphicArithmetic/AutomorphicVarieties/ModularityArgument/ThresholdArguments/DenseThresholdCases/PromodularityCaseThree.lean#L19-L25}{SeedIVAlt}.



\paragraph{9. TENSOR\_\allowbreak FACTOR\_\allowbreak TOWER — B — Actual tensor-factor tower and characters}

Original finite-field, tame, factor and character constraints; source-level construction is sizeable and has not been reconstructed here.

\emph{Source.} \href{https://github.com/CocoSgt/riemann-hypothesis-lean/blob/7cbf7a0c32fca1badf05931879a37c3f8a2b03b4/src/Development/AutomorphicArithmetic/AutomorphicVarieties/ModularityArgument/ThresholdArguments/DenseThresholdCases/PromodularityCaseThreeOfLemmas.lean#L131-L153}{UpperTheory.V.Cu.promodularity\_\allowbreak case\_\allowbreak three\_\allowbreak of\_\allowbreak chain\_\allowbreak reduction}.



\paragraph{10. CHAIN\_\allowbreak RANK\_\allowbreak FOUR — C — Actual SomeIn-to-AllIn propagation}

Same source Space and PM ideal, residual point Good4, original bound 30(C+1)\textasciicircum{}2 <= d; geometry must preserve all identities.

\emph{Source.} \href{https://github.com/CocoSgt/riemann-hypothesis-lean/blob/7cbf7a0c32fca1badf05931879a37c3f8a2b03b4/src/Development/AutomorphicArithmetic/AutomorphicVarieties/ModularityArgument/ThresholdArguments/DenseThresholdCases/PromodularityCaseThree.lean#L26-L32}{ChainRankFour}.



\paragraph{11. CHAIN\_\allowbreak DESC4 — B — Original residual chain descent}

Finite residual descent retaining source Chain4 compatibility.

\emph{Source.} \href{https://github.com/CocoSgt/riemann-hypothesis-lean/blob/7cbf7a0c32fca1badf05931879a37c3f8a2b03b4/src/Development/AutomorphicArithmetic/AutomorphicVarieties/ModularityArgument/ThresholdArguments/DenseThresholdCases/PromodularityCaseThreeOfLemmas.lean#L131-L153}{UpperTheory.V.Cu.promodularity\_\allowbreak case\_\allowbreak three\_\allowbreak of\_\allowbreak chain\_\allowbreak reduction}.



\paragraph{12. RANK4\_\allowbreak CHAIN\_\allowbreak REDUCTION — C — Seeded reduction to the original residual chain}

Select T1 before every later Deep tower; retain one compatible seed and Chain4 for every original residual lift over Q.

\emph{Source.} \href{https://github.com/CocoSgt/riemann-hypothesis-lean/blob/7cbf7a0c32fca1badf05931879a37c3f8a2b03b4/src/Development/AutomorphicArithmetic/AutomorphicVarieties/ModularityArgument/ThresholdArguments/DenseThresholdCases/PromodularityCaseThree.lean#L33-L39}{RankFourChainReductionDense}.

\emph{Source.} \href{https://github.com/CocoSgt/riemann-hypothesis-lean/blob/7cbf7a0c32fca1badf05931879a37c3f8a2b03b4/src/Development/AutomorphicArithmetic/AutomorphicVarieties/ModularityArgument/ThresholdArguments/DenseThresholdCases/CuPChainIV.lean#L16-L82}{UpperTheory.V.Cu.chainIV\_\allowbreak alt}.



\paragraph{13. TASK\_\allowbreak COUNT\_\allowbreak V\_\allowbreak INSTALL — B — Complete the original CountV finite assembly}

From exactly the listed source host/cohomology inputs, complete countV\_\allowbreak of\_\allowbreak finiteness\_\allowbreak duality\_\allowbreak euler\_\allowbreak host, including original host/frame formation, normalized length, the CaseV guard and R1<=2\textasciicircum{}c2.

\emph{Completed inputs (version 0.50).} CAF9-CAF11 (full original CountV HostData character seed, including permitted ell=3), MFC1-MFC14 on original Gen/GC hypotheses, Inherited CountV consumers at their supplied-frame scope

\emph{Remaining producer scope.} Complete the seven original Ok fields phi/gen/cl/ec/b1p/b1/h1 in the common GC construction, and separately the full entry implications Gen\_\allowbreak gd and Gen\_\allowbreak host. CAF supplies the character SeedClause uniformly before all later compatible fields and primes; MFC supplies the actual linked common field, all twelve UFacts, the uniform b0 class-rank bound, both strict cone classes, Big and 21 of 28 Ok fields. In the original coefficient-DVR case it also joins the CAF seed, with its modulus chosen first. The same actual marked g1 and fixed coefficient embedding must then receive FinEnd, AutEnd3 and AutEnd2, the actual rank-two and rank-three automorphic seeds, full ModUAlt/PfromAut, and original local/automorphic descent to GdEnd. UCE supplies the full GdEnd consumer, not the producer. The alternative GdEndAlt route still requires its exact additional local/descent fields. Do not exclude ell=3 from the independent CountV HostData seed scope merely because the final GD.gen guard excludes it. Complete the original F3 local coefficient bridges and F4 common-overfield readings, with valuation/norm/GoodCoef/iq and the entire coefficient field, and actual complex Hecke standard/no-Jordan plus arithmetic character data on the same form datum. Construct the actual base rows, host/frame lengths, nonzero chains and group ties, MultEq support, eigensystem/carrier/parts, all-embedding logLam, HD/d3b and NEb. Keep the stronger second DivRows3 scope distinct from the completed first DivDef3 and row-pair integrability. The partial common marked frame and CountV seed do not by themselves install HOST\_\allowbreak NUMBER\_\allowbreak S, HOST\_\allowbreak FRAME\_\allowbreak LENGTH, TASK\_\allowbreak HOST\_\allowbreak VAN\_\allowbreak E\_\allowbreak INSTALL, TASK\_\allowbreak HOST\_\allowbreak XPP\_\allowbreak INSTALL or TASK\_\allowbreak COUNT\_\allowbreak V\_\allowbreak INSTALL.

\emph{Source.} \href{https://github.com/CocoSgt/riemann-hypothesis-lean/blob/7cbf7a0c32fca1badf05931879a37c3f8a2b03b4/src/Development/AutomorphicArithmetic/AutomorphicVarieties/HeightAndVanishing/Counts.lean#L69-L69}{CountV}.

\emph{Source.} \href{https://github.com/CocoSgt/riemann-hypothesis-lean/blob/7cbf7a0c32fca1badf05931879a37c3f8a2b03b4/src/Development/AutomorphicArithmetic/AutomorphicVarieties/ModularityArgument/ThresholdArguments/DenseThresholdCases/CountVOf.lean#L33-L60}{UpperTheory.V.UnitaryThreeInertia.countV\_\allowbreak of\_\allowbreak finiteness\_\allowbreak duality\_\allowbreak euler\_\allowbreak host}.



\paragraph{14. THM\_\allowbreak S6\_\allowbreak S — C — Original unitary-three ordinary support and automorphy}

Original host/frame and fixed iota -> one marked g1 -> every universal R,u -> weight-zero finiteness and every Type-I BPrime\_\allowbreak s point.

\emph{Source.} \href{https://github.com/CocoSgt/riemann-hypothesis-lean/blob/7cbf7a0c32fca1badf05931879a37c3f8a2b03b4/src/Development/AutomorphicArithmetic/AutomorphicVarieties/ModularityArgument/ThresholdArguments/DenseThresholdCases/ThmS6CoreS.lean#L17-L26}{ThmS6\_\allowbreak s}.

\emph{Source.} \href{https://github.com/CocoSgt/riemann-hypothesis-lean/blob/7cbf7a0c32fca1badf05931879a37c3f8a2b03b4/src/Development/AutomorphicArithmetic/AutomorphicVarieties/ModularityArgument/ThresholdArguments/DenseThresholdCases/ThmS6CoreS.lean#L12-L16}{BPrime\_\allowbreak s}.



\paragraph{15. HOST\_\allowbreak FRAME\_\allowbreak LENGTH — B — Original Case-V host/frame and normalized-length installation}

Actual HostData, frame, local-condition equality, abstract count and normalized original Selmer length must all refer to the same objects.

\emph{Completed inputs (version 0.50).} CAF and MFC common-frame subresults at their exact original inputs, DNE weight/host-degree dictionary on its supplied tables

\emph{Remaining producer scope.} Complete the original F3 local coefficient bridges and F4 common-overfield readings, with valuation/norm/GoodCoef/iq and the entire coefficient field, and actual complex Hecke standard/no-Jordan plus arithmetic character data on the same form datum. Construct the actual base rows, host/frame lengths, nonzero chains and group ties, MultEq support, eigensystem/carrier/parts, all-embedding logLam, HD/d3b and NEb. Keep the stronger second DivRows3 scope distinct from the completed first DivDef3 and row-pair integrability. The partial common marked frame and CountV seed do not by themselves install HOST\_\allowbreak NUMBER\_\allowbreak S, HOST\_\allowbreak FRAME\_\allowbreak LENGTH, TASK\_\allowbreak HOST\_\allowbreak VAN\_\allowbreak E\_\allowbreak INSTALL, TASK\_\allowbreak HOST\_\allowbreak XPP\_\allowbreak INSTALL or TASK\_\allowbreak COUNT\_\allowbreak V\_\allowbreak INSTALL.

\emph{Source.} \href{https://github.com/CocoSgt/riemann-hypothesis-lean/blob/7cbf7a0c32fca1badf05931879a37c3f8a2b03b4/src/Development/AutomorphicArithmetic/AutomorphicVarieties/ModularityArgument/ThresholdArguments/DenseThresholdCases/CountVOf.lean#L33-L60}{UpperTheory.V.UnitaryThreeInertia.countV\_\allowbreak of\_\allowbreak finiteness\_\allowbreak duality\_\allowbreak euler\_\allowbreak host}.



\paragraph{16. HOST\_\allowbreak NUMBER\_\allowbreak S — C — Uniform original-host rank-three catalogue}

One exponent before all closures, coefficient rings, levels, characters, frames, embeddings and marked points; one separating trace function controls all later finite pairwise nonisomorphic collections.

\emph{Completed inputs (version 0.50).} CAF and MFC common-frame subresults at their exact original inputs, Inherited host-count consumers

\emph{Remaining producer scope.} Complete the original F3 local coefficient bridges and F4 common-overfield readings, with valuation/norm/GoodCoef/iq and the entire coefficient field, and actual complex Hecke standard/no-Jordan plus arithmetic character data on the same form datum. Construct the actual base rows, host/frame lengths, nonzero chains and group ties, MultEq support, eigensystem/carrier/parts, all-embedding logLam, HD/d3b and NEb. Keep the stronger second DivRows3 scope distinct from the completed first DivDef3 and row-pair integrability. The partial common marked frame and CountV seed do not by themselves install HOST\_\allowbreak NUMBER\_\allowbreak S, HOST\_\allowbreak FRAME\_\allowbreak LENGTH, TASK\_\allowbreak HOST\_\allowbreak VAN\_\allowbreak E\_\allowbreak INSTALL, TASK\_\allowbreak HOST\_\allowbreak XPP\_\allowbreak INSTALL or TASK\_\allowbreak COUNT\_\allowbreak V\_\allowbreak INSTALL.

\emph{Source.} \href{https://github.com/CocoSgt/riemann-hypothesis-lean/blob/7cbf7a0c32fca1badf05931879a37c3f8a2b03b4/src/Development/AutomorphicArithmetic/AutomorphicVarieties/ModularityArgument/ThresholdArguments/DenseThresholdCases/HostNumberS.lean#L18-L28}{HostNumber\_\allowbreak s}.



\paragraph{17. TASK\_\allowbreak HOST\_\allowbreak VAN\_\allowbreak E\_\allowbreak INSTALL — B — Transport K/duality to literal HostVanE}

Use the already supplied LemmaK/local-global duality and finiteness on the actual CountV host/frame and exact original coefficient modules.

\emph{Source.} \href{https://github.com/CocoSgt/riemann-hypothesis-lean/blob/7cbf7a0c32fca1badf05931879a37c3f8a2b03b4/src/Development/AutomorphicArithmetic/AutomorphicVarieties/UnitaryThreeDeformations/HostNumber.lean#L9-L12}{HostVanE}.

\emph{Source.} \href{https://github.com/CocoSgt/riemann-hypothesis-lean/blob/7cbf7a0c32fca1badf05931879a37c3f8a2b03b4/src/Development/AutomorphicArithmetic/AutomorphicVarieties/LevelCohomology/SoHostVanE.lean#L11-L18}{UpperTheory.V.X.So.hostVanE}.



\paragraph{18. TASK\_\allowbreak HOST\_\allowbreak XPP\_\allowbreak INSTALL — B — Transport CorX to literal HostXpp}

Preserve Psi\textasciicircum{}-2 eta, the prime-to-ell quotient, actual local conditions and the original CountV host while installing BR7.

\emph{Source.} \href{https://github.com/CocoSgt/riemann-hypothesis-lean/blob/7cbf7a0c32fca1badf05931879a37c3f8a2b03b4/src/Development/AutomorphicArithmetic/AutomorphicVarieties/UnitaryThreeDeformations/HostNumber.lean#L14-L17}{HostXpp}.

\emph{Source.} \href{https://github.com/CocoSgt/riemann-hypothesis-lean/blob/7cbf7a0c32fca1badf05931879a37c3f8a2b03b4/src/Development/AutomorphicArithmetic/AutomorphicVarieties/LevelCohomology/Xpp.lean#L11-L84}{UpperTheory.V.X.xpp}.



\paragraph{19. DESIGN\_\allowbreak CHOICES\_\allowbreak COMMON — C — Construct the original common design family and bundle}

For the one target HostExW, construct the original compatible DesignChoices on Cdes=WithTw(WithSm Cprime), where Cprime=WithChN(A intersect target class), A is the one Built construction class, and cT(eps,y)=cTmin(y). Retain every original field, tuple, word, metric and common bundle duty. DSZ10 and RMP11 give size/rank properties of class members, not member existence. The original RecExP/RecAtG and common object/row route remains coupled. Construct ONE original common design family R and its recursive/geometric half with original Inputs/Pins/Holds/ExAt/InGeo and required CoreP. PKT now reconstructs full exists\_\allowbreak pkT, rowPk\_\allowbreak nonempty and exists\_\allowbreak moved with Copy/Coh on those original given inputs and a supplied fitting later stage: one catalogue, both ChT, complete LfT, actual integral lattices and all metric/measure/period constants are retained; one numerical record precedes every eta and V/S/N/G. RDR supplies one Side0 before all readers, ReadPk/pins and full ReadL with exact o.eta=q.eta. The initial actual geometry, Fits/division/central conditions, primitive rational tests and eta agreement must still be produced. Inherited CSR supplies finite columns, one ReadC before later twists, period/evaluation and scalar/error transports, RowSideDetermined and same-R Bundle/BTw/BWit/BunU consumers. RPI supplies exact RowPairInt for actual pairs and finite translates with original mu\_\allowbreak b/fd0; first literal DivDef3 at the same updated R,F,Y,Z,C,sEE follows from its retained original readings, Pins/Holds, stored nonzero IsItW chain, five real-group predicates and group tie. Remaining common fields include host/frame, MultEq nonzero-part support, that original nonzero chain/group tie, eigensystem/carrier/parts and all-embedding logLam, HD/d3b, NEb, actual rational tests and base rows. Retain every simultaneous family hypothesis and eventual threshold. Stronger second DivRows3 remains separate; universal quotient-height CuspRapidDecay and actual common member existence are not inferred from the comparison constructors. Version 0.41 supplies the finite phase-correct rational theta route: exact coefficient extension, split representatives, square-law/PhiIs, Pins from CoversAt, class-sum/point-identity transfer to OneMapOverL, actual-good-projector RatG and both-level LiftTestB. Its original higher covers, all-tuple class-sum uniqueness including phase rescaling, one PointIdentCore before O, base ClassSumsSpan and geometric LatFull remain explicit producers. MLat is available from the original LatFull input. This sharpens the rational-test work without claiming actual common geometry, Fits or eta agreement. Version0.46 proves finite generation of the original whole integral cusp lattice by actual integral-frame evaluation. It also proves the exact identity-cover, analytic unfolding, original-field coefficient descent, global clearing-scalar and both-level lattice-transport implications. Full spanning retains the supplied original IsNab connection and the five explicit principal-chart, level-system, raw-Gaussian, short-rule/dictionary and ramified-jet producers. The higher tuple covers, actual common class sums, PointIdentCore and ClassSumsSpan remain separate. CURRENT VERSION0.47 QUALIFICATION: Construct one original common geometric family and the higher CoversAt/all-tuple class sums and uniqueness with original phase rescalings, one PointIdentCore with common lambda,z, and base ClassSumsSpan. CNC now supplies the full original IsNab; PCL supplies the conditional trace-dual principal-lattice and quantitative normalization. Actual uniformization/HasPrinc and the five principal-chart, finite-level-system, Gaussian/reproduction, short-rule/simultaneous-dictionary and ramified-jet producer groups remain. Keep all original both-level, coefficient-field, Fits and eta agreements. Completed new original producers: ICL1-ICL21, CNC1-CNC16, PCL1-PCL12. Any earlier narrative naming these completed producers as outstanding is historical; this qualification is the current scope authority. CURRENT VERSION0.48 QUALIFICATION: Construct the original finite holomorphic FamilyUnif/N1c, normalized theta families and local lifting, then the common principal charts and finite-level theta systems. PUN now supplies the exact fixed-space models, all-family center lattices, A=1 exponent matching and individual primitive points with the full point-equality criterion; CNC supplies the full canonical connection and PCL the quantitative trace-dual normalization. Complete higher CoversAt and all-tuple class sums/uniqueness with the original phase rescalings, one PointIdentCore with common lambda,z, and base ClassSumsSpan. The Gaussian/reproduction, short-rule/simultaneous-dictionary and ramified-jet producers, and all original both-level, coefficient-field, Fits and eta agreements remain explicit. Completed current producers or conditional consumers: ICL1-ICL21, CNC1-CNC16, PCL1-PCL12, PUN1-PUN24, GCF1-GCF18 at their stated scopes. Earlier frontier descriptions are historical; this qualification is the current scope authority. Original whole-contract premises and quantifiers are retained. CURRENT VERSION0.49 QUALIFICATION: Construct the adjoining simultaneous principal framed points and test-ring dictionary. Full finite holomorphic FamilyUnif/N1c, normalized theta families and local lifts are now supplied by HUF. The next bounded target is allTest\_\allowbreak of\_\allowbreak secBaseChange: represent FrPt by the Lie-eigenframe bundles and theta-section frame line, use the full n=1 SecBaseChangeSt on every nonflat coefficient base change, and prove finite-generated subalgebra descent with the exact TRing/FromTest transport. Then supply the full FramesOK/CoverTri/FibTri and IsUnf/hasPrinc package, including adapted mod-eight symplectic basis, normalized theta multiplier/weight and arithmetic invariance. Continue higher CoversAt, all-tuple class sums and uniqueness with original phase rescalings, one PointIdentCore with common lambda,z, base ClassSumsSpan, Gaussian/reproduction, short-rule dictionaries and ramified jets, retaining both-level, coefficient-field, Fits and eta agreements. Common-family nonemptiness remains distinct from compatibility of an already supplied member. Completed current producers or conditional consumers: ICL1-ICL21, CNC1-CNC16, PCL1-PCL12, PUN1-PUN24, HUF1-HUF21 (full original FamilyUnif/N1c), GCF1-GCF18 and WRS1-WRS10 at their stated scopes. Earlier frontier descriptions are historical; this qualification is the current scope authority. Original whole-contract premises and quantifiers are retained.

\emph{Completed inputs (version 0.50).} TDR1-TDR13 (arbitrary finite descent and original-premise AllTest), NFR1-NFR13 (ThetaData, IsNormUn, every FrameOK field and simultaneous first field), HUF1-HUF21, PUN1-PUN24, ICL1-ICL21, CNC1-CNC16, PCL1-PCL12, GCF and WRS at original supplied-member scope

\emph{Remaining producer scope.} Complete the three remaining original FramesOK fields: triangular weight covariance, arithmetic invariance and holomorphic readings. Then construct CoverTri and FibTri and assemble full IsUnf/hasPrinc on the same Moduli, model, lattice, section and tangent frames. TDR now proves arbitrary-algebra finite descent and the whole AllTest implication under the exact SecBaseChangeSt; NFR proves ThetaData, IsNormUn and every FrameOK field, including mean exactly one. These results do not yet supply the other simultaneous frame fields or the full principal charts. Continue the original finite level system, Gaussian/reproduction, short-rule/simultaneous dictionaries and ramified jets; higher CoversAt, all-tuple class sums and uniqueness with their phase rescalings, one PointIdentCore with common lambda and z, and base ClassSumsSpan remain. Preserve both levels, the complete coefficient field, Fits and eta agreements, and original common-family quantifier order. Common design nonemptiness remains distinct from the compatibility of a supplied member.



\paragraph{20. LOCAL\_\allowbreak NBUILT\_\allowbreak COMMON — C — Actual local specialization on the common selected family}

Construct, with the required closed-class property on the same Des(C,3), one eventual family of actual NBuilt certificates for the original Fm.D(x,i), retaining the certificate before forgetting to TSel. HLP now reconstructs full HeckeLift0, all-index commutation, HeckeStd0 and HeckeIntTw from the original MtInput and complex Hecke facts, in the full ambient GenAt range k<=nV. Every maximal ideal is an actual supported residual character class; literal MtSub is the O-span of projections of the original MtLat generators, is finite free, and localizes onto the exact NCov.part. Literal Psi has its exact range and off-space projection error, both original span assertions hold, and ActEx/ActHK are constructed on the same lattice under original semisimplicity. Original MjLat relation quotients, universal coefficient-map relation fields, selected tame generators, coefficient embeddings and all identity transports give the complete nonempty\_\allowbreak spec constructor on its full printed input list. A single compositum combines the four original coefficient requirements with big fixed beforehand. Whole ndSpec\_\allowbreak exists still needs actual common geometric/choice MtInput, complex HeckeStdPlain and the no-Jordan producer, original F3 bridges with valuation/GoodCoef/iq, original F4 common-overfield reading of the entire coefficient field, and exact arithmetic character/reciprocity/admissibility suppliers. Original AutSpec, Galois specialization and the remaining NBuilt fields on that same installation also remain. Do not replace the actual projected lattice or maps by abstract objects with the same dimensions.

\emph{Completed inputs (version 0.50).} TDR and NFR actual-frame subresults, CAF and MFC actual marked-frame subresults, WRS selected built-class witness extraction at existing membership

\emph{Remaining producer scope.} Continue the distinct retained-defect, period-height and normalized spectral producers, and the original compatible second-layer selection on one common arithmetic/geometric family. GLOBAL\_\allowbreak LOCAL\_\allowbreak GAIN remains closed at its complete v0.49 scope; the current denominator work only completes its separate consumer. Common design existence, NBUILT compatibility, original true-column records, shared fields/embeddings, Fits and eta agreement and eventual-family thresholds remain explicit. A stronger historical WPt/TSel construction may still serve other producers but is not needed to reopen the proved selected-class local-gain estimate. Complete the three remaining original FramesOK fields: triangular weight covariance, arithmetic invariance and holomorphic readings. Then construct CoverTri and FibTri and assemble full IsUnf/hasPrinc on the same Moduli, model, lattice, section and tangent frames. TDR now proves arbitrary-algebra finite descent and the whole AllTest implication under the exact SecBaseChangeSt; NFR proves ThetaData, IsNormUn and every FrameOK field, including mean exactly one. These results do not yet supply the other simultaneous frame fields or the full principal charts. Continue the original finite level system, Gaussian/reproduction, short-rule/simultaneous dictionaries and ramified jets; higher CoversAt, all-tuple class sums and uniqueness with their phase rescalings, one PointIdentCore with common lambda and z, and base ClassSumsSpan remain. Preserve both levels, the complete coefficient field, Fits and eta agreements, and original common-family quantifier order.



\paragraph{21. SECOND\_\allowbreak LAYER\_\allowbreak COMPATIBLE\_\allowbreak SELECTION — C — Original second-layer selector and class duties}

For the same A, supply Duty on A and a Closed common original TSel family whose selected records are the original true-column rec7 records (or have the entire dependent tuple transported), with their real witness and original coefficient-field root data. Existence/compatibility remains arithmetic; merely abstractly isomorphic fields or independently selected columns do not suffice. CSR21–22 now construct the same-R Bundle/BTw/BWit/BunU outputs from the exact selected realization/row-side packet and NEb, with the original EvP/ParEv quantifier order. Actual common selection, the five stated realization packets and NEb remain upstream; none is replaced by an independently chosen compatible record. Version 0.38 closes the selected analytic extension and reconstructs the first updated DivDef3 consumer; it leaves the original common selection, geometric readings, rational tests, base rows, host/parts/NEb and exact retained DI premises upstream, with unchanged EvP/ParEv thresholds. Version 0.39 supplies the original given-data reader constructors and exact RowPairInt/first DivDef3 on original premises. It leaves one actual common geometric/comparison/test member, eta agreement, compatible selection, rational tests, base rows, host/parts/NEb and original EvP/ParEv thresholds upstream. Version 0.40 also supplies original comparison/RowPk/fitting-stage Moved algorithms and exact common data transport. The actual initial geometric member, fitting tests, eta agreement and compatible second-layer selector with original eventual thresholds remain upstream. Version 0.41 supplies the finite phase-correct rational theta route: exact coefficient extension, split representatives, square-law/PhiIs, Pins from CoversAt, class-sum/point-identity transfer to OneMapOverL, actual-good-projector RatG and both-level LiftTestB. Its original higher covers, all-tuple class-sum uniqueness including phase rescaling, one PointIdentCore before O, base ClassSumsSpan and geometric LatFull remain explicit producers. MLat is available from the original LatFull input. This sharpens the rational-test work without claiming actual common geometry, Fits or eta agreement. CURRENT VERSION0.49 QUALIFICATION: Continue the distinct global denominator, period/height, normalized spectral and second-layer compatible-selection producers, together with existence of the original common arithmetic/geometric family. GLOBAL\_\allowbreak LOCAL\_\allowbreak GAIN and its actual GroupCompatibilityFamily supplier are now complete on the unchanged selected built class; no additional WPt/TSel realization is needed for that estimate. Do not reopen local first-growth, mean, auxiliary admission, selected D3 tops/Adm2F, the single GIs2 stage or the exact comparison/kernel transport. A full RecT or TSel may remain useful for other tasks, but is not claimed by this local-gain closure. Completed current producers or conditional consumers: GCF1-GCF18 at their stated scopes, WRS1-WRS10 original selected-class compatibility/local gain. Earlier frontier descriptions are historical; this qualification is the current scope authority. Original whole-contract premises and quantifiers are retained.

\emph{Completed inputs (version 0.50).} WRS1-WRS10 original selected-class compatibility/local gain, DNE1-DNE14 original denominator consumer, TDR/NFR and CAF/MFC subresults at stated scope, Inherited exact same-record realization/transport consumers

\emph{Remaining producer scope.} Continue the distinct retained-defect, period-height and normalized spectral producers, and the original compatible second-layer selection on one common arithmetic/geometric family. GLOBAL\_\allowbreak LOCAL\_\allowbreak GAIN remains closed at its complete v0.49 scope; the current denominator work only completes its separate consumer. Common design existence, NBUILT compatibility, original true-column records, shared fields/embeddings, Fits and eta agreement and eventual-family thresholds remain explicit. A stronger historical WPt/TSel construction may still serve other producers but is not needed to reopen the proved selected-class local-gain estimate.



\paragraph{22. BASE\_\allowbreak THETA\_\allowbreak READINGS — C — Actual base reading on the common bundle}

For every family in the original class, eventually on the SAME selected Bundle and for each i<m, supply the actual R,F,Read,ReadL, weak models/Weil data and phase-one pins. Identify their LiftTie.theta, mtRead, mjRead and ThDef3 basis with that Bundle.P.A3, P.Zs, P.S0 and its stored maps. Exact field, embeddings, metrics, measures, row family and weak choices are retained. This remains one part of the coupled BTH12 producer. PKT now supplies the complete original comparison packet, RowPk and fitting-stage Moved with Copy/Coh on original given geometry, Inputs/Pins/Holds/ExAt/InGeo and Fits. Its ChT/LfT/Reps, actual lattice carriers, original Haar and norm constants, period maps, representatives and all-place RrT are constructed together; its one numerical record precedes every eta/V/S/N/G. RDR gives the original level-zero side, moved-Read and ReadPk/pin constructors on those objects, with one Side0 before every reader and all twelve Read ties. Full ReadL retains the exact o.eta(j,v)=q.eta(j,v) guard for every 1<=j<=m and place v; arbitrary Over does not imply it. Inherited CSR fixes one ReadC before all TwArgs, transports both metrics and all thirteen reading maps, and gives actual period/evaluation with kappa=(c*nP)\textasciicircum{}(-1/2). RPI supplies exact row-pair integrability, so first updated DivDef3 follows from the retained original Pins/Holds, nonzero chain, five real-group predicates and group tie. Initial theta extension remains complete. Remaining work is actual initial common geometry and its original predicates, Fits/primitive rational tests, eta agreement and outer family/eventual/same-Bundle duties, reFace compatibility, MultEq support, original chain/group data, actual base rows and any stronger DivRows3 theta square. No new compatible reading, family threshold or Haar normalization is substituted. Version 0.41 supplies the finite phase-correct rational theta route: exact coefficient extension, split representatives, square-law/PhiIs, Pins from CoversAt, class-sum/point-identity transfer to OneMapOverL, actual-good-projector RatG and both-level LiftTestB. Its original higher covers, all-tuple class-sum uniqueness including phase rescaling, one PointIdentCore before O, base ClassSumsSpan and geometric LatFull remain explicit producers. MLat is available from the original LatFull input. This sharpens the rational-test work without claiming actual common geometry, Fits or eta agreement.

\emph{Completed inputs (version 0.50).} TDR1-TDR13, NFR1-NFR13, HUF1-HUF21, PUN1-PUN24, CNC1-CNC16, Inherited actual-reading consumers at their original inputs

\emph{Remaining producer scope.} Complete the three remaining original FramesOK fields: triangular weight covariance, arithmetic invariance and holomorphic readings. Then construct CoverTri and FibTri and assemble full IsUnf/hasPrinc on the same Moduli, model, lattice, section and tangent frames. TDR now proves arbitrary-algebra finite descent and the whole AllTest implication under the exact SecBaseChangeSt; NFR proves ThetaData, IsNormUn and every FrameOK field, including mean exactly one. These results do not yet supply the other simultaneous frame fields or the full principal charts. Continue the original finite level system, Gaussian/reproduction, short-rule/simultaneous dictionaries and ramified jets; higher CoversAt, all-tuple class sums and uniqueness with their phase rescalings, one PointIdentCore with common lambda and z, and base ClassSumsSpan remain. Preserve both levels, the complete coefficient field, Fits and eta agreements, and original common-family quantifier order. Complete the original F3 local coefficient bridges and F4 common-overfield readings, with valuation/norm/GoodCoef/iq and the entire coefficient field, and actual complex Hecke standard/no-Jordan plus arithmetic character data on the same form datum. Construct the actual base rows, host/frame lengths, nonzero chains and group ties, MultEq support, eigensystem/carrier/parts, all-embedding logLam, HD/d3b and NEb. Keep the stronger second DivRows3 scope distinct from the completed first DivDef3 and row-pair integrability. The partial common marked frame and CountV seed do not by themselves install HOST\_\allowbreak NUMBER\_\allowbreak S, HOST\_\allowbreak FRAME\_\allowbreak LENGTH, TASK\_\allowbreak HOST\_\allowbreak VAN\_\allowbreak E\_\allowbreak INSTALL, TASK\_\allowbreak HOST\_\allowbreak XPP\_\allowbreak INSTALL or TASK\_\allowbreak COUNT\_\allowbreak V\_\allowbreak INSTALL.



\paragraph{23. BASE\_\allowbreak THETA\_\allowbreak RATIONAL\_\allowbreak TESTS — C — Simultaneous original base rational tests}

On the SAME original Bundle/reading/theta map, for every i<m supply one finite separating complex-linear test family whose values lie in the fixed embedded coefficient field on theta of EVERY original rational source vector and on EVERY stored target-coordinate basis vector. BTH8/BTH11 use this same LiftTestB for rational adjoint image and matrix coordinates. RTG now reconstructs prescribed-field coefficient extension, original split class representatives/RepsOK, the self-dual split-lattice square law and finite residue PhiIs; every Pins field follows from the actual CoversAt. With the base coefficient record literally fixed, exact class-sum existence/uniqueness including phase-rescaled O, PointIdentCore with common lambda,z before every O, and base ClassSumsSpan imply all-tuple finite lift span and OneMapOverL. The phase is exactly cj(twAt embLL A3.A.v0)(R.phs i). RatG uses the actual good projector eOp(y), and full LiftTestB follows at BOTH original levels. These finite implication proofs are complete. Remaining original producers are the prescribed higher covers, those actual class sums/uniqueness, PointIdentCore, base ClassSumsSpan and the base/common geometric datum; exact MLat is supplied only from original LatFull, whose geometric production remains separate. Finite dimension or an unrelated rational theta model cannot replace these same-object identities. Original outer eventual-family/Bundle quantifiers remain. Version0.46 proves finite generation of the original whole integral cusp lattice by actual integral-frame evaluation. It also proves the exact identity-cover, analytic unfolding, original-field coefficient descent, global clearing-scalar and both-level lattice-transport implications. Full spanning retains the supplied original IsNab connection and the five explicit principal-chart, level-system, raw-Gaussian, short-rule/dictionary and ramified-jet producers. The higher tuple covers, actual common class sums, PointIdentCore and ClassSumsSpan remain separate. CURRENT VERSION0.47 QUALIFICATION: Construct one original common geometric family and the higher CoversAt/all-tuple class sums and uniqueness with original phase rescalings, one PointIdentCore with common lambda,z, and base ClassSumsSpan. CNC now supplies the full original IsNab; PCL supplies the conditional trace-dual principal-lattice and quantitative normalization. Actual uniformization/HasPrinc and the five principal-chart, finite-level-system, Gaussian/reproduction, short-rule/simultaneous-dictionary and ramified-jet producer groups remain. Keep all original both-level, coefficient-field, Fits and eta agreements. Completed new original producers: ICL1-ICL21, CNC1-CNC16, PCL1-PCL12. Any earlier narrative naming these completed producers as outstanding is historical; this qualification is the current scope authority. CURRENT VERSION0.48 QUALIFICATION: Construct the original finite holomorphic FamilyUnif/N1c, normalized theta families and local lifting, then the common principal charts and finite-level theta systems. PUN now supplies the exact fixed-space models, all-family center lattices, A=1 exponent matching and individual primitive points with the full point-equality criterion; CNC supplies the full canonical connection and PCL the quantitative trace-dual normalization. Complete higher CoversAt and all-tuple class sums/uniqueness with the original phase rescalings, one PointIdentCore with common lambda,z, and base ClassSumsSpan. The Gaussian/reproduction, short-rule/simultaneous-dictionary and ramified-jet producers, and all original both-level, coefficient-field, Fits and eta agreements remain explicit. Completed current producers or conditional consumers: ICL1-ICL21, CNC1-CNC16, PCL1-PCL12, PUN1-PUN24. Earlier frontier descriptions are historical; this qualification is the current scope authority. Original whole-contract premises and quantifiers are retained. CURRENT VERSION0.49 QUALIFICATION: Construct the adjoining simultaneous principal framed points and test-ring dictionary. Full finite holomorphic FamilyUnif/N1c, normalized theta families and local lifts are now supplied by HUF. The next bounded target is allTest\_\allowbreak of\_\allowbreak secBaseChange: represent FrPt by the Lie-eigenframe bundles and theta-section frame line, use the full n=1 SecBaseChangeSt on every nonflat coefficient base change, and prove finite-generated subalgebra descent with the exact TRing/FromTest transport. Then supply the full FramesOK/CoverTri/FibTri and IsUnf/hasPrinc package, including adapted mod-eight symplectic basis, normalized theta multiplier/weight and arithmetic invariance. Continue higher CoversAt, all-tuple class sums and uniqueness with original phase rescalings, one PointIdentCore with common lambda,z, base ClassSumsSpan, Gaussian/reproduction, short-rule dictionaries and ramified jets, retaining both-level, coefficient-field, Fits and eta agreements. Completed current producers or conditional consumers: ICL1-ICL21, CNC1-CNC16, PCL1-PCL12, PUN1-PUN24, HUF1-HUF21 (full original FamilyUnif/N1c and theta family). Earlier frontier descriptions are historical; this qualification is the current scope authority. Original whole-contract premises and quantifiers are retained.

\emph{Completed inputs (version 0.50).} TDR1-TDR13, NFR1-NFR13, HUF1-HUF21, PUN1-PUN24, ICL1-ICL21, CNC1-CNC16, PCL1-PCL12, Inherited both-level RatG/LiftTestB consumers at their exact supplied inputs

\emph{Remaining producer scope.} Complete the three remaining original FramesOK fields: triangular weight covariance, arithmetic invariance and holomorphic readings. Then construct CoverTri and FibTri and assemble full IsUnf/hasPrinc on the same Moduli, model, lattice, section and tangent frames. TDR now proves arbitrary-algebra finite descent and the whole AllTest implication under the exact SecBaseChangeSt; NFR proves ThetaData, IsNormUn and every FrameOK field, including mean exactly one. These results do not yet supply the other simultaneous frame fields or the full principal charts. Continue the original finite level system, Gaussian/reproduction, short-rule/simultaneous dictionaries and ramified jets; higher CoversAt, all-tuple class sums and uniqueness with their phase rescalings, one PointIdentCore with common lambda and z, and base ClassSumsSpan remain. Preserve both levels, the complete coefficient field, Fits and eta agreements, and original common-family quantifier order.



\paragraph{24. SLR\_\allowbreak BASE\_\allowbreak ROWS\_\allowbreak PRODUCER — C — Actual base row identities on that same bundle}

For the SAME selected Bundle, common original fields and readings used by the other three BTH12 inputs, supply actual RowsHold0\_\allowbreak 3 uniformly eventually for every family. Its original primitive row realization, spectrum/type and basis identifications must be retained. BTH supplies a theta consumer only; it does not construct the base row identity. No separate witness may be chosen merely to satisfy this row requirement.



\clearpage
\section{All coefficient test rings for the original framed points}

\paragraph{The assertion proved.}
Fix the original \(E,H,D,l,X_o,c\), and put \(K=c.F\) and
\(X=X_o\times_{\operatorname{Spec}K_A(l)}\operatorname{Spec}K\).
The supplied \texttt{Moduli} record, including its universal object,
smoothness and quasiprojectivity, is retained.  The only additional
section-theoretic premise is the original
\(\mathrm{BC}=\texttt{Out.\allowbreak Sec\allowbreak Base\allowbreak Change\allowbreak St}\).  This premise is applied
to a locally Noetherian base such as \(X_o\) or \(X\).  Its conclusion
allows \emph{every} commutative coefficient ring on the other side of
the base change.  We never replace it by a flat-base-change assertion,
and we do not assume the stronger predicate
\texttt{Sec\allowbreak Base\allowbreak Change\allowbreak Any\allowbreak St} with arbitrary original base schemes.

The two principal conclusions of this chapter are
\begin{align}
 &\text{for every commutative }K\text{-algebra }S
   \text{ and }y\in\operatorname{FrPt}(S),\notag\\
 &\qquad\exists\ S_0\subseteq S\text{ a finitely generated }K
       \text{-subalgebra},\notag\\
 &\qquad y_0\in\operatorname{FrPt}(S_0),\quad
       \operatorname{FrMap}(S_0\hookrightarrow S,y_0,y);       \tag{TDR.1}\\
 &\text{if }S\text{ is reduced, there are }n,I,\phi,y'\text{ with}\notag\\
 &\qquad I\subseteq K[t_1,\ldots,t_n]\text{ radical},\quad
       \phi:K[t_1,\ldots,t_n]/I\longrightarrow S,\notag\\
 &\qquad y'\in\operatorname{FrPt}(K[t_1,\ldots,t_n]/I),\notag\\
 &\qquad
       \operatorname{FrMap}(\phi,y',y).                       \tag{TDR.2}
\end{align}
Thus (TDR.1) includes nonreduced, non-Noetherian and non-finite-type
\(S\), and (TDR.2) is the literal \texttt{All\allowbreak Test} conclusion.
The algebra map \(\phi\), tangent vectors and section are part of the
conclusion; the result is not merely equality of isomorphism classes
of abelian varieties or a statement about field-valued points.

\paragraph{Geometric dictionary.}
Let \(\pi:A\to X\) and \(L\) be the pullbacks of the actual universal
abelian scheme and its prescribed principal theta line.  Write
\(m_\tau=D.\mathrm{mult}(\tau)\), for
\(\tau:E\hookrightarrow\mathbb C\), and
\(\tau_K:E\to K\) for the same embedding with codomain restricted
using \(c.\mathrm{conj}\).  The finite set of all these embeddings,
including those with \(m_\tau=0\), is retained.  A framed point has an
underlying map \(x:\operatorname{Spec}S\to X\), ordered tangent
eigenbases \((u_{\tau,i})_{i<m_\tau}\), and a section \(s\) which is
an \(S\)-module generator of \(H^0(A_x,L_x)\).  The source notation
\(\operatorname{PPx}(x)\) is the same line, expressed as its
one-fold \texttt{LPow} on the canonical pullback.  We identify these
lines only by that canonical pullback isomorphism.

The standard algebraic foundations used below are the description of
tangent vectors by dual numbers, Lie base change for a smooth group
scheme, the sheaf properties of regular functions and line-bundle
sections, and affine scheme/ring duality.  Frame schemes and all
finite-descent arguments needed here are constructed below.  No
analytic uniformization, choice of Haar measure, or new result about
automorphic forms is used.

\subsection{The actual frame functor over every coefficient algebra}

% BEGIN TDR1
\paragraph{TDR1. Tangent eigenbases are frames of fixed vector bundles.}
\label{claim:TDR1}
There are finite locally free bundles \(V_\tau\) of rank \(m_\tau\)
on \(X\) such that, for every commutative \(K\)-algebra \(S\) and
\(x:\operatorname{Spec}S\to X\), the source's
\(\tau\)-eigen tangent vectors are exactly
\(\Gamma(\operatorname{Spec}S,x^*V_\tau)\).
Under this identification \texttt{Is\allowbreak Comb\allowbreak At} is ordinary linear
combination.  The source's \texttt{basis} condition is precisely an
isomorphism \(S^{m_\tau}\simeq\Gamma(\operatorname{Spec}S,x^*V_\tau)\).
These identifications commute with arbitrary \(K\)-algebra maps.

\begin{proof}
A morphism from \(\operatorname{Spec}S[\epsilon]/(\epsilon^2)\)
to a smooth group scheme which reduces to its identity is a section
of its relative Lie bundle.  On an affine chart, its pullback on
functions is \(a\mapsto e^*a+\epsilon\delta(a)\), where
\(\delta\) is an \(S\)-linear derivation at the identity.  Multiplying
two such group-valued points adds their derivations: terms involving
two infinitesimal factors vanish because \(\epsilon^2=0\).
Precomposition by \(\epsilon\mapsto b\epsilon\) multiplies the
derivation by \(b\).  Consequently the finite product in
\texttt{Is\allowbreak Comb\allowbreak At} is exactly the sum of the indicated scalar
multiples.  This also proves that the source \texttt{tan} predicate,
including both its structural and identity-section equations, is
the relative tangent condition just described.

The actual \(\mathcal O_E\)-action differentiates to the Lie bundle
\(V=\operatorname{Lie}(A/X)\).  It extends to the split algebra
\[
 E\otimes_{\mathbb Q}K\simeq\prod_{\tau:E\hookrightarrow\mathbb C}K,
 \qquad a\otimes b\longmapsto(\tau_K(a)b)_\tau .             \tag{TDR.3}
\]
Here \(K\) contains every conjugate of \(E\), and characteristic
zero makes \(E/\mathbb Q\) separable.  Explicitly, choose a primitive
element of \(E\); the Lagrange interpolation polynomials in its
distinct \(K\)-valued roots give mutually orthogonal idempotents
\(e_\tau\) with sum one.  All denominators are nonzero elements
of \(K\), hence remain units in every nonzero \(K\)-algebra.
The zero ring causes no exception: all the modules and assertions
over its empty spectrum are trivial.

Set \(V_\tau=e_\tau V\).  These are direct summands of a finite
locally free bundle, hence finite locally free.  The eigencondition
for every \(a\in\mathcal O_E\) is equivalent to membership in this
summand.  Indeed it implies the condition for \(E\), by clearing
integer denominators, and the interpolation idempotents then isolate
the required summand.  Conversely the idempotent summand has precisely
the required character.  The original determinant/multiplicity
condition in the universal object, evaluated over algebraic closures
of residue fields, says that its rank is \(m_\tau\).  A finite
locally free module has locally constant rank detected on these
fibers, so the same rank holds on \(X\).

Lie formation for a smooth group scheme commutes with every base
change: equivalently, pull back its locally free invariant
differentials and dualize.  Each \(e_\tau\) is an idempotent over
\(K\), so its summand also commutes with base change, including
nonflat maps.  Finally, unique expression of every eigenvector in
the ordered \(u_{\tau,i}\) is exactly bijectivity of the indicated
linear map from \(S^{m_\tau}\).  No argument using only generic
fibers is involved.  For \(m_\tau=0\), both modules are zero and
the same statement holds with the empty basis.
\end{proof}
% END TDR1

% BEGIN TDR2
\paragraph{TDR2. The theta-section factor is one invertible sheaf with arbitrary base change.}
\label{claim:TDR2}
Under \(\mathrm{BC}\), there is an invertible sheaf \(\Theta\) on
\(X\), canonically the direct-image sheaf of the actual \(L\), and
for every \(x:\operatorname{Spec}S\to X\) a natural isomorphism
\[
 \Gamma(\operatorname{Spec}S,x^*\Theta)
       \simeq H^0(A_x,L_x).                                 \tag{TDR.4}
\]
An element on the right satisfies the literal source \texttt{gen}
condition if and only if its element on the left trivializes
\(x^*\Theta\).  Such an element remains a generator after every
algebra map \(S\to T\), with no reducedness or flatness assumption.

\begin{proof}
The original moduli scheme is smooth over the field \(K_A(l)\),
so it is locally Noetherian; the same is true after the finite field
extension to \(K\).  Its universal object is the supplied principally
polarized abelian scheme.  Apply \(\mathrm{BC}\) with \(n=1\).
The index set of sections has size \(1^g=1\).  On a neighborhood
of each point it therefore supplies one section whose pullback along
\emph{every} map from an affine scheme is a basis of the pulled-back
section module.  Ordinary pullback of the section supplies the
\texttt{Is\allowbreak Pull\allowbreak Sec} witness in that predicate.  Its two asserted
conditions, linear independence and spanning, say exactly that the
section module is freely generated by this section.

On affine subopens these sections identify \(\pi_*L\) with the
structure sheaf.  On overlaps two such identifications differ by a
unique unit: each generator is a unique scalar multiple of the other,
and composing the two expressions gives the scalar product one.
The overlap units satisfy the cocycle law by the same uniqueness.
They define an invertible sheaf \(\Theta=\pi_*L\), rather than
merely a collection of one-dimensional vector spaces on closed points.

The same local generators and overlap units pull back along any
\(x\).  The defining assertion of \(\mathrm{BC}\) identifies the
pulled-back generator with a basis of \(H^0(A_x,L_x)\) on every
affine subopen mapping to a trivializing neighborhood.  Both sides
of (TDR.4) are sheaves for the Zariski topology on the base: sections
of the pulled-back line on \(A_x\) glue over preimages of base
opens.  The local identifications consequently glue to (TDR.4),
naturally in \(x\), and they are the canonical section pullback
maps.  This also proves the assertion when \(S\) is non-Noetherian.

On an affine scheme an invertible sheaf corresponds to a projective
module of rank one.  The source \texttt{gen} says that
\(a\mapsto a s\) is a bijection from \(S\) to that module,
with scalar multiplication written chartwise as \texttt{scal}.
This is exactly its trivialization.  Tensoring that isomorphism with
an arbitrary \(S\)-algebra \(T\) preserves the isomorphism;
no flatness is necessary.  The section obtained is the actual
pullback section by the preceding naturality.
\end{proof}
% END TDR2

% BEGIN TDR3
\paragraph{TDR3. Exact base change of the original framed-point record.}
\label{claim:TDR3}
For every \(K\)-algebra map \(\phi:S\to T\) and original
\(y\in\operatorname{FrPt}(S)\), there is a unique
\(\phi_*y\in\operatorname{FrPt}(T)\) with
\(\operatorname{FrMap}(\phi,y,\phi_*y)\).
It is functorial in \(\phi\).  Thus this relation is exactly
pullback of the underlying point, each ordered eigenframe and the
chosen section, with no scalar ambiguity.

\begin{proof}
Set \(x_T=x\circ\operatorname{Spec}\phi\), pull back each
basis of \(x^*V_\tau\), and pull back the generator of
\(x^*\Theta\).  TDR1 and TDR2 give all the tangent, eigenbasis
and section-generator fields of the new record.  In dual-number
notation its tangent vectors are literally
\[
 u'_{\tau,i}
   =u_{\tau,i}\circ\operatorname{Spec}(\phi[\epsilon]),
 \qquad
 \phi[\epsilon](a+b\epsilon)=\phi(a)+\phi(b)\epsilon .       \tag{TDR.5}
\]
For the section let \(p_m:A_{x_T}\to A_x\) be the canonical
pullback map.  It commutes with both projections in the definition
of \texttt{Fr\allowbreak Map}.  On every original line-bundle chart the new
section is \(p_m^*s\), which is exactly the universally quantified
\texttt{app\allowbreak LE} equation in that definition.  If a chart is further
restricted, the equation remains true because restriction commutes
with pullback.  This verifies the literal relation, including its
last field, rather than only an isomorphism of section lines.

For uniqueness, the first two fields of \texttt{Fr\allowbreak Map} determine
the point and all tangent vectors.  The two projection equations
uniquely determine \(p_m\) by the fiber-product property, and the
section equations determine every chart component of the section.
All remaining fields are properties of these data.  The identity
and composition laws now follow either directly from (TDR.5) and
section pullback or from this uniqueness.  In particular if
\(z=\phi_*y\), then \((\psi\phi)_*y=w\) if and only if
\(\psi_*z=w\).  This is the cancellation law used in the source
descent proof.
\end{proof}
% END TDR3

% BEGIN TDR4
\paragraph{TDR4. Change of frame, and descent of equality along injections.}
\label{claim:TDR4}
For two original framed points \(y,y'\) over the same underlying
\(x:\operatorname{Spec}S\to X\), there is a unique collection
\[
 g_\tau\in\operatorname{GL}_{m_\tau}(S),\qquad a\in S^\times
 \quad\text{such that}\quad
 u'_{\tau,i}=\sum_j(g_\tau)_{ij}u_{\tau,j},\qquad s'=a s.
                                                               \tag{TDR.6}
\]
These are the original \texttt{Fr\allowbreak Chg} equations.  They commute
with every coefficient algebra map.  If \(S\hookrightarrow T\)
is injective and \(y,y'\) have identical pullbacks over \(T\),
then \(y=y'\).

\begin{proof}
Each ordered basis of a free module gives unique coordinates for
the other.  The coordinate maps in both directions are inverse,
so the matrices in (TDR.6) have unit determinant.  TDR1 identifies
these coordinate equations with \texttt{Is\allowbreak Comb\allowbreak At}; the convention
is rows, exactly as in \texttt{Fr\allowbreak Chg}.  By TDR2, \(s\) and \(s'\)
are bases of the same rank-one module.  Expressing each in the other
gives unique scalars whose product is one.  The section equation
\(s'=a s\) is the source \texttt{Sec\allowbreak BC} for the identity map.
This proves existence and uniqueness, including zero-rank
eigenspaces and arbitrary coefficient rings.

Applying a ring map to (TDR.6) gives the coordinates of the
pulled-back frames, by TDR3.  If those frames coincide, the images
of all matrices and of the scalar are respectively the identity
matrices and one.  Injectivity of the ring map forces the same
equations over \(S\).  Equation (TDR.6) then gives equality of
the original points.  The hypothesis that their underlying maps
\(x\) already agree is retained; we are not using an unjustified
assertion about arbitrary maps of schemes after an injective ring
extension.
\end{proof}
% END TDR4

% BEGIN TDR5
\paragraph{TDR5. Reduced affine charts with actual reference frames.}
\label{claim:TDR5}
Every point of \(X\) has an affine open neighborhood
\(j:\operatorname{Spec}R\hookrightarrow X\), where \(R\) is
a reduced finitely generated \(K\)-algebra, together with an
original framed point \(y_0\in\operatorname{FrPt}(R)\) whose
underlying map to \(X\) is \(j\).

\begin{proof}
The finitely many \(V_\tau\) and the line \(\Theta\) are locally
free, by TDR1 and TDR2.  Intersect trivializing neighborhoods for
all of them and then take an affine open neighborhood of the chosen
point.  The base \(X\) is smooth over \(K\), hence reduced and
locally of finite type.  Its affine open neighborhood is reduced;
its structural map is of finite type because it is affine and
locally of finite type.  Thus \(R\) has the asserted properties.

Choose the ordered bases of each \(V_\tau\) and a generator of
\(\Theta\) on this affine open.  TDR1 translates the bases into
the required dual-number tangent maps with the \(\tau_K\)-action,
and TDR2 translates the line generator into the actual section of
\(\operatorname{PPx}\), with its unique coefficient rule.
They give every field of \(y_0\).  Its underlying point is the
composition of \(j\) with the projection \(X\to X_o\), so the
two projection equations in \texttt{Open\allowbreak Chart} hold literally.
TDR12 below presents \(R\) as the original \texttt{TRing};
transport along that algebra isomorphism and TDR3 then give
\texttt{exists\_\allowbreak open\allowbreak Chart\_\allowbreak test\allowbreak Ring} with the same image in \(X\).
The finite-generation part of this proof does not assume that
the entire moduli scheme is covered by one affine chart.
\end{proof}
% END TDR5

% BEGIN TDR6
\paragraph{TDR6. The explicit universal frame algebra and its exact classifier.}
\label{claim:TDR6}
For a fixed affine chart and reference frame from TDR5, set
\[
 B=R[\,X_{\tau,ij},s\,]
       \left[\frac{1}{s\prod_\tau\det(X_\tau)}\right].       \tag{TDR.7}
\]
It is a reduced finitely generated \(K\)-algebra, and hence is an
original test ring up to a specified algebra isomorphism.  With
\(\iota:R\to B\), the universal matrices and scalar satisfy
\texttt{Frame\allowbreak Ring}.  There is an actual universal frame
\(u\in\operatorname{FrPt}(B)\) such that, for every commutative
\(K\)-algebra \(T\), map \(\psi:R\to T\), and
\(z\in\operatorname{FrPt}(T)\) over \(j\circ\operatorname{Spec}\psi\),
there is a unique \(b:B\to T\) with
\[
 b\iota=\psi,\qquad\operatorname{FrMap}(b,u,z).             \tag{TDR.8}
\]

\begin{proof}
There are finitely many variables, including one variable for the
theta-section scalar.  Localization in (TDR.7) makes \(s\) and
each determinant units: each divides their product, and a factor
of a unit in a commutative ring is a unit.  Conversely, any choice
of invertible matrices and a unit in \(T\) sends this product to
a unit.  The polynomial universal property followed by the
localization universal property gives a unique algebra map
\(B\to T\) with these values and restriction \(\psi\).
This is the entire \texttt{Frame\allowbreak Ring} statement, with its
universal quantifier over all \(T\).

A polynomial ring over a reduced ring is reduced.  One direct
verification is to embed the reduced ring into the product of
its integral-domain quotients by prime ideals and use the
coefficientwise embedding of polynomial rings; the target has no
nonzero nilpotents.  Localization also preserves reducedness:
if a localized element is nilpotent, multiplying the nilpotence
relation by a suitable denominator gives a nilpotent element in
the original ring, whence the localized element is zero.
Localization at one element is a finite-type operation, realized
by adjoining one inverse.  Therefore \(B\) is reduced and
finitely generated over \(K\).

Pull \(y_0\) back to \(B\) by TDR3, then change its bases by
the universal matrices and its section by \(s\), using TDR4.
This defines \(u\), with the source \texttt{Univ\allowbreak Frame} witness.
For (TDR.8), first pull \(y_0\) back by \(\psi\).  TDR4
gives the unique matrices and scalar comparing this reference
frame with \(z\).  Use precisely those coordinates in the
universal property of (TDR.7).  TDR3 and the compatibility in
TDR4 show that the resulting map carries \(u\) to \(z\).
If another map does so, uniqueness of the comparison coordinates
forces its values on every variable and on \(R\); it is therefore
the same map.  If \(m_\tau=0\), its determinant is one and it
contributes no variables, as required by the source.
\end{proof}
% END TDR6

% BEGIN TDR7
\paragraph{TDR7. Representability uses the actual theta and Lie frame bundles.}
\label{claim:TDR7}
The original functor \(S\mapsto\operatorname{FrPt}(S)\), with
\texttt{Fr\allowbreak Map} as base change, is represented by
\[
 \mathcal Y=
 \left(\prod_{\tau,\,X}
   \operatorname{Isom}_X(\mathcal O_X^{m_\tau},V_\tau)\right)
       \times_X\operatorname{Isom}_X(\mathcal O_X,\Theta).
                                                               \tag{TDR.9}
\]
The morphism \(\mathcal Y\to X\) is affine and of finite
presentation.  In particular \(\mathcal Y\) is separated and
locally of finite type over \(K\).

\begin{proof}
On an open where a vector bundle of rank \(m\) is trivial,
its frame scheme is \(\operatorname{Spec}R[X_{ij},1/\det X]\).
On overlaps the invertible transition matrix changes these
coordinates by matrix multiplication, an algebraic isomorphism.
The cocycle law for the bundle is the gluing law for these schemes.
This constructs the indicated \(\operatorname{Isom}\) schemes.
The line factor is constructed in exactly the same way using
\(R[s,s^{-1}]\).  All these morphisms are affine and of finite
presentation, properties which can be checked on an open cover
of their base.  Their finite fiber product has the same properties;
on the common trivializing charts its algebra is precisely (TDR.7).

An \(S\)-point of (TDR.9) consists of a map \(x\) and the
listed trivializations on \(\operatorname{Spec}S\).  On an
affine base these are exactly the global ordered bases and
section generator in TDR1--TDR2.  Thus the correspondence with
the literal source record is bijective.  Its compatibility with
algebra maps is the actual \texttt{Fr\allowbreak Map} relation, by TDR3;
no choice of a scalar or of a new model remains in this comparison.
Finally \(X\) is quasiprojective over \(K\), hence separated,
and \(\mathcal Y\to X\) is affine, hence separated.  Their
composition is separated.  Smoothness makes \(X\) locally of
finite presentation over \(K\), and the constructed morphism
has finite presentation, proving the final local finite-type
assertion.  This representability holds even for nonreduced
test coefficient algebras; the reduced-scheme representing
axiom of \texttt{Moduli} has not been extended or used outside
its original scope.
\end{proof}
% END TDR7

\subsection{Finite coefficients, without flatness assumptions}

% BEGIN TDR8
\paragraph{TDR8. Localization embeddings and simultaneous coefficient capture.}
\label{claim:TDR8}
Let \(S\) be any commutative \(K\)-algebra,
\(S_0\subseteq S\) a finitely generated subalgebra, and
\(h_1,\ldots,h_n\in S\) generate the unit ideal.
For each \(i\), let \(B_i\) be a finitely generated
\(K\)-algebra and \(a_i:B_i\to S_{h_i}\) an algebra map.
There is a finitely generated subalgebra \(S_1\subseteq S\)
containing \(S_0\) and all \(h_i\), in which the \(h_i\)
still generate the unit ideal, and maps
\[
 a_i':B_i\longrightarrow (S_1)_{h_i},\qquad
 \kappa_i:(S_1)_{h_i}\lhook\joinrel\longrightarrow S_{h_i},
 \qquad \kappa_i a_i'=a_i.                                  \tag{TDR.10}
\]
The \(\kappa_i\) are the canonical injective maps, with no
assumption that \(S_1\to S\) is flat or that the \(h_i\)
are non-zero-divisors.

\begin{proof}
First, if \(R\subseteq S\) and \(h\in R\), the canonical
map \(R_h\to S_h\) is injective.  If \(r/h^m\) maps to
zero, then \(h^N r=0\) in \(S\) for some \(N\).
The same equality holds in the subring \(R\), so \(r/h^m=0\)
in \(R_h\).  This argument includes nilpotents and zero divisors.

Choose \(c_i\in S\) with \(\sum_i c_i h_i=1\).
For each \(i\), choose a finite algebra generating set \(G_i\)
of \(B_i\).  Write
\(a_i(b)=s_{i,b}/h_i^{N_{i,b}}\) for \(b\in G_i\).
Adjoin to \(S_0\) the finitely many \(h_i,c_i,s_{i,b}\)
and call the resulting subalgebra \(S_1\).  The very same
equation \(\sum_i c_i h_i=1\) holds there, so its principal
opens still cover.  Every chosen generator image belongs to the
image of \(\kappa_i\), which is a subalgebra.  Hence the entire
image of \(a_i\) lies in that image.  Invert \(\kappa_i\)
on its image to obtain \(a_i'\); this is an algebra homomorphism
because \(\kappa_i\) is an injective algebra homomorphism.

Equivalently one may send the generators to the displayed fractions.
All relations automatically hold: their values map to zero under
the injective \(\kappa_i\), hence are already zero.  This is why
the argument needs only finite generation of the \(B_i\), not
an additional finite list of relations.  For an empty cover the
unit-ideal equation forces \(S\) to be the zero ring; the unique
zero subalgebra and empty family satisfy the assertion.
\end{proof}
% END TDR8

% BEGIN TDR9
\paragraph{TDR9. Principal-cover gluing of the original frames.}
\label{claim:TDR9}
Fix \(x:\operatorname{Spec}S\to X\) and a finite principal
cover \(D(h_i)\) with \((h_i)=S\).  Frames over its members,
with underlying points the restrictions of \(x\), glue uniquely
when their \texttt{Fr\allowbreak Map} pullbacks agree over
\(S_{h_i h_j}\).  Equality of two frames over \(x\) can also
be checked after restriction to these localizations.

\begin{proof}
For every \(\tau,i\), the local tangent vectors are sections
of \(x^*V_\tau\).  Compatibility glues them by the sheaf axiom.
Their local inverse coordinate maps glue as well, or, equivalently,
the resulting map \(\mathcal O^{m_\tau}\to x^*V_\tau\)
is an isomorphism because it is one on each member of the cover.
This gives exactly the required ordered eigenbases by TDR1.
The local theta sections glue as sections of \(x^*\Theta\),
and their local inverses show that the glued section is a frame
of that line.  By TDR2 its section on \(A_x\) satisfies the
source \texttt{gen} condition.  All pullback equations are those
of TDR3, so the result is a frame of the original functor.

For clarity, the coefficient uniqueness implicit here works over
arbitrary \(S\).  If \(b\in S\) becomes zero in every
\(S_{h_i}\), choose \(N_i\) with \(h_i^{N_i}b=0\).
The ideal generated by these powers has radical the unit ideal,
and hence is itself the unit ideal.  A linear combination of the
powers is one, so \(b=0\).  Thus local equality of all matrix
entries and of the section scalar in TDR4 implies global equality.
Existence of glued module sections is the ordinary affine sheaf
equalizer, which follows by clearing finitely many denominators
and the same unit-ideal argument.  The proof does not require
faithful flatness of any individual localization map.
\end{proof}
% END TDR9

% BEGIN TDR10
\paragraph{TDR10. Finite-subalgebra descent of an actual scheme point.}
\label{claim:TDR10}
Let \(Z\) be a separated, locally finite-type \(K\)-scheme.
Every map \(x:\operatorname{Spec}S\to Z\), for an arbitrary
commutative \(K\)-algebra \(S\), is the base change of a map
\(x_0:\operatorname{Spec}S_0\to Z\) for a finitely generated
\(K\)-subalgebra \(S_0\subseteq S\).  A prescribed finite
set of elements of \(S\) may be included in \(S_0\).

\begin{proof}
Choose affine open neighborhoods \(U_i=\operatorname{Spec}B_i\)
of the image points, with \(B_i\) finitely generated over \(K\).
The inverse images cover \(\operatorname{Spec}S\), and its
principal opens form a basis.  Quasicompactness therefore gives
a finite principal cover \(D(h_i)\) such that
\(x(D(h_i))\subseteq U_i\).  We have maps
\(a_i:B_i\to S_{h_i}\).  Choose coefficients of an equation
\(\sum_i c_i h_i=1\), and use TDR8 to capture all the maps
in localizations of a single finitely generated subalgebra \(S_0\).
Start that construction with the algebra generated by the
prescribed finite set.  It gives maps
\(x_i:D_{S_0}(h_i)\to U_i\to Z\).

They agree on overlaps.  To see this without assuming flatness,
put \(R=(S_0)_{h_i h_j}\) and \(T=S_{h_i h_j}\).
The two overlap maps to \(Z\) coincide after base change to
\(T\), since both give the original \(x\).
As \(Z\) is separated over \(K\), their equalizer is the
pullback of the closed diagonal, hence is a closed subscheme of
\(\operatorname{Spec}R\), say defined by an ideal \(J\subseteq R\).
After base change to \(T\) the ideal is zero; thus every element
of \(J\) maps to zero in \(T\).  TDR8 says \(R\to T\)
is injective, so \(J=0\).  The two maps consequently agree
already over \(R\).  The maps \(x_i\) glue to \(x_0\),
and their base changes glue to \(x\).

The zero ring is handled by its empty spectrum.  This proof uses
only finite data, and in particular it does not impose a
Noetherian or reducedness assumption on \(S\).  It applies to
\(Z=X\) because the supplied moduli scheme is quasiprojective,
hence separated, and smooth, hence locally finite type.
It also applies directly to the actual frame scheme \(\mathcal Y\)
of TDR7.  The source's more general locally-finite-presentation
factorization theorem in \texttt{Moduli\allowbreak Point\allowbreak From\allowbreak Fg\allowbreak Subalgebra}
is therefore available here at every instance actually used.
\end{proof}
% END TDR10

% BEGIN TDR11
\paragraph{TDR11. Full finitely generated descent of every original framed point.}
\label{claim:TDR11}
Assuming exactly \(\mathrm{BC}\) and the supplied original data,
(TDR.1) holds.  This is the complete original
\texttt{Test\allowbreak Descent.\allowbreak exists\_\allowbreak fg\_\allowbreak descent} theorem.

\begin{proof}
There is a short proof: regard \(y\) as an \(S\)-point of
the explicitly constructed \(\mathcal Y\), use TDR10 for that
separated locally finite-type scheme, and apply the exact functor
identification TDR7.  It gives the actual frame \(y_1\), and
its base-change equality is precisely \texttt{Fr\allowbreak Map} by TDR3.
We give the source-compatible local construction as well, to
record explicitly which finite data are retained and how the
underlying point and frames are joined.

First descend the underlying map of \(y\) to
\(x_0:\operatorname{Spec}S_0\to X\) by TDR10.
Choose a finite principal cover \(D(g_i)\) of
\(\operatorname{Spec}S_0\) such that each restriction of
\(x_0\) factors through a chart \(U_i=\operatorname{Spec}R_i\)
of TDR5.  Let \(h_i\) be the same elements in \(S\).
They still generate its unit ideal.  Let \(y_i\) be the
pullback of \(y\) to \(S_{h_i}\), and let
\(\psi_i:R_i\to S_{h_i}\) be the map supplied by the chart
factorization.  TDR6 supplies a universal frame \(u_i\) over
an explicit algebra \(B_i\) and a unique map
\[
 a_i:B_i\to S_{h_i},\qquad
 a_i|_{R_i}=\psi_i,\qquad(a_i)_*u_i=y_i.                   \tag{TDR.11}
\]

Apply TDR8 simultaneously to these finitely many \(a_i\),
starting with \(S_0\).  It gives one finitely generated
\(S_1\subseteq S\) and maps \(a_i'\) as in (TDR.10).
Let \(v_i=(a_i')_*u_i\) over \((S_1)_{h_i}\).
Its underlying point is exactly the restriction of
\(x_1=x_0\circ\operatorname{Spec}(S_1\supseteq S_0)\).
Indeed the chart map
\(R_i\to(S_0)_{g_i}\to(S_1)_{h_i}\) and the restriction
of \(a_i'\) to \(R_i\) become \(\psi_i\) in \(S_{h_i}\)
by (TDR.11).  The canonical localization map is injective,
so these two algebra maps already agree.  This proves equality
of the actual underlying maps, before attempting to glue frames.

On a double overlap \((S_1)_{h_i h_j}\), the two restrictions
of \(v_i,v_j\) have that same underlying map and become the
restriction of \(y\) after the canonical injection into
\(S_{h_i h_j}\).  TDR4 therefore makes them equal on the
original overlap.  TDR9 glues them to
\(y_1\in\operatorname{FrPt}(S_1)\).  Its base change to
\(S\) and \(y\) agree on the principal cover \(D(h_i)\),
so TDR9 gives their global equality.  In the source notation this
is the literal witness
\(\operatorname{FrMap}(S_1.\mathrm{val},y_1,y)\).

Every choice was made from the given \(y\).  There is no
independent choice of a new moduli object, no scalar rescaling
after the section is descended, and no assertion that the
inclusion \(S_1\subseteq S\) is flat.  Nilpotents in \(S\)
are retained if they occur among the finite coefficients.
\end{proof}
% END TDR11

% BEGIN TDR12
\paragraph{TDR12. Reduced finite-type algebras are exactly the original test-ring presentations.}
\label{claim:TDR12}
A commutative \(K\)-algebra \(R\) is reduced and finitely
generated if and only if it is isomorphic to
\(p.R=K[t_1,\ldots,t_n]/I\) for some literal
\(p:\texttt{TRing}\ K\), where \(I\) is radical.

\begin{proof}
Choose finitely many algebra generators to obtain a surjective
\(K\)-algebra map \(f:K[t_1,\ldots,t_n]\to R\).
If \(R\) is reduced and \(u^N\in\ker f\), then
\(f(u)^N=0\), whence \(f(u)=0\).  Thus \(I=\ker f\)
is radical, and the first isomorphism theorem gives the specified
algebra isomorphism \(p.R\simeq R\).  Conversely the images
of the variables generate the quotient; if a quotient class has
a nilpotent power, the radical condition makes that class zero.
This proof includes the unit ideal and the zero algebra, and
does not require a field-valued test point or a prime ideal.
\end{proof}
% END TDR12

% BEGIN TDR13
\paragraph{TDR13. The original AllTest theorem and the complex-frame return.}
\label{claim:TDR13}
Under \(\mathrm{BC}\), for every original \(E,H,D,l,X_o,c\),
the original \(\texttt{All\allowbreak Test}\ X_o\ c\) holds.  In
particular, for every embedding \(\sigma:K\to\mathbb C\)
and every actual \(y\in\operatorname{FrPt}(\mathbb C_\sigma)\),
the exact source \texttt{From\allowbreak Test} predicate holds.
The assertion applies simultaneously to every member of any
already chosen family of complex framed points.

\begin{proof}
Let \(S\) be reduced and let \(y\in\operatorname{FrPt}(S)\).
TDR11 gives a finitely generated subalgebra \(S_1\subseteq S\),
an actual frame \(y_1\), and
\(\operatorname{FrMap}(\iota,y_1,y)\), where \(\iota\)
is inclusion.  The subalgebra is reduced: a nilpotent element
would map to a nilpotent in \(S\), hence to zero, and inclusion
is injective.  TDR12 gives \(p:\texttt{TRing}\ K\) and
an algebra isomorphism \(e:p.R\xrightarrow{\sim}S_1\).
Pull \(y_1\) back by \(e^{-1}\) to obtain the literal
\(y'\in\operatorname{FrPt}(p.R)\), using TDR3.  Define
\(\phi=\iota\circ e\).  Functoriality gives
\[
 \phi_*y'=\iota_*e_*(e^{-1})_*y_1
         =\iota_*y_1=y.                                    \tag{TDR.12}
\]
Thus \((p,\phi,y')\) is exactly the triple required by
\texttt{From\allowbreak Test}, proving (TDR.2) and the full universally
quantified \texttt{All\allowbreak Test} statement.

The coefficient algebra \(\mathbb C_\sigma\), denoted
\texttt{CAt}\(\,\sigma\) in the source, is the ordinary
reduced ring \(\mathbb C\) with its \(K\)-structure given
by \(\sigma\).  Apply the result to it.  This does not alter
the embedding, the underlying point, its chosen eigenbases or
its chosen section.  Since the theorem holds for every such
frame, it applies to every \(\operatorname{un}(k,x)\) in
an already chosen family.  The test-ring witness may depend
on \((k,x)\); this is the original pointwise quantifier in
\texttt{exists\_\allowbreak unf\_\allowbreak has\allowbreak Princ}, not a demand for one test
ring presenting all members at once.
\end{proof}
% END TDR13

\subsection{Exact interfaces and the remaining geometry}

The endpoint \texttt{Test\allowbreak Descent.\allowbreak exists\_\allowbreak fg\_\allowbreak descent} is supplied
by TDR11, for every commutative coefficient algebra, without
reducedness.  The endpoint
\texttt{Frames.\allowbreak all\allowbreak Test\_\allowbreak of\_\allowbreak sec\allowbreak Base\allowbreak Change} is supplied by TDR13,
for every reduced coefficient algebra.  The source's intermediate
open-chart, universal-frame, frame-classification, injective
base-change equality and principal-cover gluing results are
proved in TDR3--TDR9, rather than left as new supplier assumptions.
TDR7 records why the same concrete frame functor is used throughout.

The hypothesis \(\mathrm{BC}\) remains named.  This chapter does
not prove its original instance from the full package of abelian
scheme laws, and does not assert existence of the supplied
\texttt{Moduli} record.  No \texttt{Sec\allowbreak Twist\allowbreak St} premise was needed.
The remaining normalized analytic frame family still has to
satisfy the exact \texttt{Frame\allowbreak OK}, Borel-weight, arithmetic
invariance and holomorphic-reading conditions.  TDR13 supplies its
\texttt{From\allowbreak Test} field as soon as that same family has been
constructed.  The covering/fiber theorems, \texttt{Is\allowbreak Unf},
\texttt{exists\_\allowbreak unf\_\allowbreak has\allowbreak Princ}, spanning and common-family
obligations are not inferred from test-ring descent alone.

\paragraph{Pinned source correspondence.}
The primary sources are the pinned algebraic definitions of
\texttt{Fr\allowbreak Pt}, \texttt{Fr\allowbreak Map}, \texttt{TRing} and
\texttt{All\allowbreak Test}, the unitary-existence proof files, and the
original section-base-change predicate.  The accompanying
\texttt{source\_\allowbreak manifest.\allowbreak json} supplies complete immutable paths,
exact line ranges and file hashes; the claim provenance records
map every TDR claim to its source declarations.

The principal endpoint is in
\texttt{src/\allowbreak Development/\allowbreak Abelian\allowbreak Variety\allowbreak Theory/\allowbreak Proofs/\allowbreak Unitary\allowbreak Existence/\allowbreak Framed\allowbreak Point\allowbreak Reduced\allowbreak From\allowbreak Test\allowbreak Ring.\allowbreak lean},
at lines 227--269 for finite-subalgebra descent and 276--290
for the return to all test rings.  TDR1--TDR7 prove the geometric
and universal-frame suppliers; TDR8--TDR11 prove finite coefficient
capture and gluing; TDR12--TDR13 perform the radical-quotient
presentation and exact source return.  The optional representability
shortcut in TDR7 and TDR11 is a human simplification; its literal
identification with the source frame records is proved before use.


\section{Actual normalized theta points and the complete single-frame theorem}
\label{sec:normalized-frames-v050}

\status{NFR1--NFR13 construct the original \texttt{Theta\allowbreak Data},
\texttt{Is\allowbreak Norm\allowbreak Un} and every clause of \texttt{Frame\allowbreak OK}, on the actual
family and the actual universal line.  They retain the original supplied
\texttt{Moduli}; HUF21 and PUN11 supply the family and its models.
The Gaussian mean is exactly one.  The three further fields of
\texttt{Frames\allowbreak OK}---triangular weight covariance, arithmetic-group
invariance and holomorphic readings---are not asserted here.
No principal-uniformization or whole RH/GRH obligation is closed.}

Fix the original number field $E$, CM host $H$, datum $D_w$, level $l$,
moduli record $X_o$, coefficient field $c.F$, and embedding
$\sigma:c.F\hookrightarrow\mathbb C$ satisfying \texttt{Good\allowbreak Emb}.
Let
\[
 e:\{1,\ldots,g\}\simeq
       \coprod_{\rho:E\hookrightarrow\mathbb C}\{0,\ldots,D_w.\mathrm{mult}(\rho)-1\},
 \qquad \mathrm{ty}_j=\operatorname{relab}(c,\sigma)((ej)_1).
 \tag{NFR.1}
\]
Use the original $F:\texttt{Family\allowbreak Unif}$ and
$M:\texttt{Models\allowbreak Of}(F,r)$, with \texttt{Model\allowbreak Base} and \texttt{Neg\allowbreak Xi}.
These are exactly the conclusions of HUF21 and PUN11 in their stated
scope, not new arithmetic hypotheses.  Fix a component $k$ and write
\[
\begin{gathered}
D=F.D_i(k),\quad \Lambda=D.\mathrm{La},\quad
\phi=M.\phi_k:\Lambda\simeq_{\mathcal O_E}M.X_k.\Lambda,\\
G=M.X_k.G_m,\quad \xi_0=M.X_k.\xi_0,\quad q=D.q.
\end{gathered}
\]
The original trace identity is
\[
 E_D(a,b)=\operatorname{Tr}_{E/\mathbb Q}
                  \bigl(-\xi_0\,\operatorname{herm}_G(\phi a,\phi b)\bigr).
 \tag{NFR.2}
\]
Every Hermitian form in this chapter is linear in its first argument:
$H(z,w)=z^tH\overline w$.  Thus the first $q$ diagonal signs of
the source's \texttt{sgn} are $+1$, and the remaining $r-q$ signs
are $-1$.  This convention is essential for (NFR.2).

\subsection{The exact frame domain and a simultaneous theta datum}

% BEGIN NFR1
\paragraph{NFR1: the original embedding and column dictionary.}
\label{claim:NFR1}
Put $\Psi=M.X_k.\Psi$, and let $\mathcal G$ be its actual frame
domain.  There is a bijection
$j\leftrightarrow(\tau,a)\in\Psi\times\{0,\ldots,r-1\}$ such that
\[
\begin{array}{c|c|c}
 &\mathrm{ty}_j& a\\ \hline
 a<q&\overline\tau&q-1-(ej)_2\\
 a\ge q&\tau&q+(ej)_2
\end{array}                                                     \tag{NFR.3}
\]
For a frame $x=(X_\tau)_{\tau\in\Psi}$ and $v\in E^r$, its
literal \texttt{per\allowbreak Fr} coordinates are
\[
 p_x(v)_j=
 \begin{cases}
  (\operatorname{col}_aX_\tau)^tG_\tau\overline{\tau(v)},&a<q,\\
  (X_\tau^{-1}\tau(v))_a,&a\ge q.
 \end{cases}                                                    \tag{NFR.4}
\]
In particular $P_x(a)=p_x(\phi a)$ defines the unique additive
map with the source's \texttt{Is\allowbreak Per\allowbreak Fr} formula.

\begin{proof}
The map $\operatorname{relab}(c,\sigma)$ permutes the embeddings
of $E$: it comes from applying the coefficient-field embedding to
the chosen realization of every embedding in its splitting field.
Its inverse is the inverse field automorphism on that splitting field.
It commutes with conjugation, since conjugation on a CM field is
the unique involution over its maximal real subfield.  The
\texttt{Good\allowbreak Emb} equation says it preserves the test
$\rho(H.\mathrm{rt})=H.\mathrm{rt}_0$.
Consequently the type has multiplicity $D_w.p$ at each member of
$\Psi$ and $D_w.q$ at its conjugate.  Since
$r=D_w.p+D_w.q$ and $q=D_w.q$, send the former letters to
columns $q,q+1,\ldots,r-1$ and the latter, in reverse order,
to $q-1,q-2,\ldots,0$.  This is (NFR.3), including the empty
blocks when $p$ or $q$ is zero.

Substitution in the definition of \texttt{per\allowbreak Fr} gives (NFR.4).
The two cases are additive in $v$; multiplication of $v$ by
$a\in\mathcal O_E$ multiplies the first case by $\overline\tau(a)$
and the second by $\tau(a)$.  These are precisely $\mathrm{ty}_j(a)$.
The inclusion of $M.X_k.\Lambda$ in $E^r$ composed with $\phi$
is injective, has that exact range, and spans $E^r$.
The trace, conjugation and signature data are the corresponding
\texttt{Model\allowbreak Base}, \texttt{Data} and \texttt{Neg\allowbreak Xi} fields.
No different rational lattice or reordered type is substituted.
\end{proof}
% END NFR1

% BEGIN NFR2
\paragraph{NFR2: these are holomorphic periods on an open connected domain.}
\label{claim:NFR2}
For every $x\in\mathcal G$, the map $P_x$ of NFR1 satisfies the
entire original \texttt{Is\allowbreak Period} predicate.  Each function
$x\mapsto P_x(a)$ is holomorphic on $\mathcal G$, and
$\mathcal G$ is open and connected.

\begin{proof}
Write $A_\tau=X_\tau^tG_\tau\overline{X_\tau}$.
The defining inequalities of $\mathcal G$ say that its leading
minors are nonzero with successive signs
$s_a=1$ for $a<q$, $s_a=-1$ otherwise.
Hermitian elimination gives an upper unitriangular $U_\tau$ with
\[
 U_\tau^tA_\tau\overline{U_\tau}
       =\operatorname{diag}(d_{\tau,a}),\qquad s_a d_{\tau,a}>0.
 \tag{NFR.5}
\]
Here one successively replaces column $a$ by itself minus its
projections onto earlier orthogonal columns, and makes the
conjugate row operations.  The $a$th diagonal entry is the ratio
of consecutive leading determinants.  This proves the stated signs
and also proves that the algorithm never divides by zero.
Rescaling each column by $|d_{\tau,a}|^{-1/2}$ gives $X_{0,\tau}$
with Gram matrix $\operatorname{diag}(s_a)$ and
$X_\tau=X_{0,\tau}b_\tau$ for an invertible upper triangular $b_\tau$.

In an orthonormal frame, put $c=X_{0,\tau}^{-1}\tau(v)$.
Formula (NFR.4) is $\overline c_a$ for $a<q$ and $c_a$ otherwise.
It is therefore a real linear isomorphism.  For arbitrary $b_\tau$,
the period vector changes, in the unreversed column order, by
\[
 Q_\tau=
 \begin{pmatrix} (b_{\tau,++})^t&0\\0&(b_{\tau,--})^{-1}\end{pmatrix}.
 \tag{NFR.6}
\]
Indeed an early column involves only earlier columns; a late row
of $b_\tau^{-1}$ involves only later rows.  Thus (NFR.6) is
invertible and the whole period map remains a real linear isomorphism.

For $\xi=-\xi_0$, write $\tau(\xi)=i\alpha_\tau$,
$\alpha_\tau=|\tau(\xi)|>0$.
The trace identity, pairing conjugate embeddings, becomes
\[
 E_D(a,b)=-2\sum_{\tau\in\Psi}\alpha_\tau
          \operatorname{Im}\operatorname{herm}_{G_\tau}
                         (\tau\phi a,\tau\phi b).
 \tag{NFR.7}
\]
In the orthonormal coordinates this is the imaginary part of the
positive diagonal form with entries $2\alpha_\tau$ on the
coordinates (NFR.4).  Conjugating that form by $Q^{-1}$ supplies
the positive Hermitian form for arbitrary $x$.
The image of the actual $\mathcal O_E$-lattice is discrete and
spans the full real space: after choosing an $E$-basis, it lies
between two nonzero integral multiples of $\mathcal O_E^r$,
whose Minkowski embedding is a full real lattice.  The real
isomorphism just proved preserves these properties.
Together with NFR1 this checks every \texttt{Is\allowbreak Period} clause.

The negative-type coordinates in (NFR.4) are complex linear in
the entries of $X_\tau$; the other coordinates are rational with
denominator $\det X_\tau$.  They are holomorphic on the domain.
Its definition is a finite conjunction of open inequalities in
continuous minors and nonzero determinants, so it is open.

For connectedness, (NFR.5) expresses every frame as an
orthonormal frame times an upper triangular matrix with positive
real diagonal.  The latter set is connected.  The orthonormal
frames for a fixed $G_\tau$ form a translate of $U(q,r-q)$.
This group is connected: a positive $q$-plane for the standard
form is uniquely the graph of a matrix $Z$ with $Z^*Z<I$;
the matrix ball is convex, and the positive square roots of
$I-Z^*Z$ and $I-ZZ^*$ give a continuous orthonormal graph frame.
The remaining choices are $U(q)\times U(r-q)$, each connected
by diagonalizing a unitary matrix and varying its eigenvalue
arguments to zero.  This proves connectedness of each factor,
then of their finite product and of $\mathcal G$.
The argument also includes $q=0$, $q=r$, and $r=0$.
\end{proof}
% END NFR2

% BEGIN NFR3
\paragraph{NFR3: an integral symplectic basis with the exact modulus eight.}
\label{claim:NFR3}
There is a basis
$\mathrm{bas}:\mathbb Z^g\oplus\mathbb Z^g\simeq\Lambda$ such that
\[
 E_D(\mathrm{bas}(u,v),\mathrm{bas}(u',v'))
                      =-u\cdot v'+v\cdot u',
 \tag{NFR.8}
\]
and, for the actual host element $w_0$,
\[
 w_0^4\mathrm{bas}(e_i,0)\equiv\mathrm{bas}(e_i,0)\pmod{8\Lambda},
 \qquad w_0^4\mathrm{bas}(0,e_i)\equiv0\pmod{8\Lambda}.
 \tag{NFR.9}
\]

\begin{proof}
The finite projective $\mathcal O_E$-module $\Lambda$ is finite
and torsion free over $\mathbb Z$, hence free.  Its rank is
$[E:\mathbb Q]r=2g$ by NFR1.  The form $E_D$ is unimodular
and alternating.  An integral symplectic basis exists by repeatedly
choosing a primitive vector, using unimodularity to pair it to one,
and splitting off the resulting hyperbolic plane.

The host equations $2w_0=1+\mathrm{rt}$,
$\mathrm{rt}^2=-d$ and $d\equiv7\pmod {16}$ imply, with
$m=2\lfloor d/16\rfloor+1$,
\[
 w_0^2=w_0-2m,\quad \overline{w_0}=1-w_0,\quad
 w_0^4+(1-w_0)^4\equiv1,\quad
 w_0^4(1-w_0)^4\equiv0\pmod8.
 \tag{NFR.10}
\]
Thus multiplication by $w_0^4$ on $\Lambda/8\Lambda$ is an
idempotent $e_+$; its adjoint for $E_D$ is $1-e_+$.
Its image and kernel are complementary isotropic direct summands.
The perfect form pairs them perfectly, so both are free of rank
$g$ over the local ring $R=\mathbb Z/8\mathbb Z$.
Choose a basis of the image and its dual basis in the kernel,
with the sign in (NFR.8).  This gives a symplectic basis modulo eight.

We justify lifting this basis integrally.  The reduction map
$\operatorname{Sp}_{2g}(\mathbb Z)\to\operatorname{Sp}_{2g}(R)$
is surjective by the following elementary elimination.
The symmetric shears
$\left(\begin{smallmatrix}I&S\\0&I\end{smallmatrix}\right)$,
$\left(\begin{smallmatrix}I&0\\S&I\end{smallmatrix}\right)$,
the simultaneous elementary changes
$\operatorname{diag}(C,C^{-t})$, and signed interchanges of
hyperbolic pairs all lift when $C$ is elementary.
The first column of a symplectic matrix is unimodular; over the
local ring $R$ it has a unit entry.  Move that entry to its first
upper coordinate.  A lower symmetric shear kills all lower
coordinates: choose its off-diagonal entries in the first column
to kill coordinates $2,\ldots,g$, and then its $(1,1)$ entry
to kill the first.  Elementary changes kill the remaining upper
coordinates.  A diagonal scaling on the first hyperbolic pair
normalizes the surviving unit.  This scaling also lifts, because
for every unit $u$ the product of the three two-dimensional shears
$U(u)L(-u^{-1})U(u)$, followed by the signed interchange, is
$\operatorname{diag}(u,u^{-1})$; lift the entries of the shears.

The first column is now $e_1$.  Its symplectic partner has lower
first coordinate one.  An upper symmetric shear fixing $e_1$
kills its upper coordinates.  An elementary change fixing $e_1$
then kills its remaining lower coordinates.  The other columns,
orthogonal to this pair, have zero first-pair coordinates; induction
finishes the reduction.  Reversing these lifted operations lifts
the given symplectic matrix.  Apply such a lift to the initial
integral symplectic basis.  Its reduction is the chosen image/kernel
basis, so (NFR.9) holds in the literal lattice, proving the claim.
\end{proof}
% END NFR3

% BEGIN NFR4
\paragraph{NFR4: the whole original \texttt{Theta\allowbreak Data} exists.}
\label{claim:NFR4}
For every component $k$ there is an actual
$T:\texttt{Theta\allowbreak Data}(e,F,M,k)$, with its original total period
function, coordinate isomorphism, adapted basis, normalizing
matrices, fixed half-characteristic and universal-line frames.

\begin{proof}
Take $T.Pf(x)=P_x$ from NFR1--NFR2 and the basis of NFR3.
Choose a complex linear coordinate isomorphism
$\mathrm{co}:\mathbb C^m\simeq\prod_{\tau\in\Psi}M_r(\mathbb C)$;
the preimage $B=\mathrm{co}^{-1}(\mathcal G)$ is open and connected.
Let $A(b),B_1(b)$ be the column matrices of the first and second
halves of $P_{\mathrm{co}(b)}\mathrm{bas}$.
The first real half-lattice spans a real isotropic $g$-plane $W$.
Its complex span is all of $\mathbb C^g$: otherwise $W\cap iW$
would contain $iw\ne0$ with $w\in W$, and
$\operatorname{Im}H(iw,w)=H(w,w)>0$ would contradict isotropy.
Thus $A$ is invertible.  Set $h=A^{-1}$ and $\tau=A^{-1}B_1$.

For clarity, the Riemann relations follow directly here.
The form $H'=A^tH\overline A$ is real symmetric and positive,
since its imaginary part vanishes on the first half-basis.
The mixed pairing in (NFR.8) says
$H'\operatorname{Im}\tau=I$.  Hence $Y=\operatorname{Im}\tau$
is positive and $H'=Y^{-1}$.  The imaginary part on the second
half-basis then says $\operatorname{Re}\tau$ is symmetric.
This gives every normalization equation in $T.\mathrm{norm}$.
The matrices $h,\tau$ are holomorphic because the period entries
are holomorphic and $\det A\ne0$.  Define them arbitrarily outside
$B$; all their asserted properties have the original domain guard.

Apply the actual $F.\mathrm{theta}$ field, supplied in full by
HUF20--HUF21, on this open connected base, with this prescribed
basis, sign $-1$, and these matrices.  It gives two constant
half-characteristics $c_1,c_2\in\{0,1/2\}^g$ and the actual
functions $fr_j(b,z)$ in the charts of the universal rigidified
line.  Its five clauses give exactly $T.fr_{\rm hol}$,
$T.fr_{\rm ne}$, $T.fr_{\rm coc}$, $T.fr_{\rm rig}$ and
$T.fr_{\rm sec}$.  Together with NFR1--NFR3 these fill every
field of \texttt{Theta\allowbreak Data}, including its modulo-eight
\texttt{adapt} field.  This proves the full source
\texttt{nonempty\_\allowbreak theta\allowbreak Data} conclusion.
\end{proof}
% END NFR4

\subsection{The actual algebraic section and its normalization}

% BEGIN NFR5
\paragraph{NFR5: theta convergence, nonvanishing and the section line.}
\label{claim:NFR5}
Let $\tau=X+iY$ be symmetric with $Y>0$, and let $a,b$ be real
half-characteristics.  With the source convention,
\[
 \theta[a,b](\tau,w)=\sum_{n\in\mathbb Z^g}
   \exp\{\pi i(n+a)^t\tau(n+a)+2\pi i(n+a)^t(w+b)\}.
 \tag{NFR.11}
\]
It is entire in $w$, is not identically zero, and
\[
 \theta(w+m+\tau n)=J_{m,n}(w)\theta(w),\quad
 J_{m,n}(w)=e^{\pi i(2a\cdot m-2n\cdot b-n^t\tau n-2n\cdot w)}.
 \tag{NFR.12}
\]
For a fiber of the $T$ in NFR4, its actual algebraic section
$s_0$ from $T.fr_{\rm sec}$ is nonzero and spans the whole space
of algebraic sections of the actual principal line.

\begin{proof}
On a compact set of $w$, the norm of a summand is at most
$\exp(-c|n|^2+C|n|+C)$ for $c>0$ depending on the least
eigenvalue of $Y$.  The same bound with a polynomial factor holds
after any derivative.  Thus the series and its derivatives converge
normally.  Translation by $m$ gives its characteristic factor;
translation by $\tau n$ follows by shifting the summation index.
This proves (NFR.12) with the stated sign.
Multiplying by $e^{-2\pi ia\cdot w}$ and integrating against
$e^{-2\pi ik\cdot u}$ for $u\in[0,1]^g$ extracts the nonzero
coefficient
$e^{\pi i(k+a)^t\tau(k+a)+2\pi i(k+a)\cdot b}$.
Hence the function is not zero.

We spell out why this yields the actual one-dimensional section
space, without adding a section-rank assumption.
Fix the fiber and pull an arbitrary algebraic section back by its
exponential.  Division by the nowhere-zero $T.fr_j$ gives a
single entire function $f$ on $\mathbb C^g$: the chart transition
identities ensure agreement on overlaps.  The section $s_0$
has scalar $\theta(hz)$.
Since exponential points $z$ and $z+\lambda$ are equal, its
values and (NFR.12) imply
$fr_j(z)=J_{m,n}(hz)fr_j(z+\lambda)$ whenever
$h\lambda=m+\tau n$.
Initially cancel $\theta(hz)$ on its nonzero set.  This set is
dense: if an entire function vanished on a ball, its restriction
to every complex line through that ball would vanish identically.
Continuity extends the identity across all its zeros.
The same identity applied to $f$ shows that $F(w)=f(h^{-1}w)$
satisfies (NFR.12).

Every entire function with the $\mathbb Z^g$-translation law
has the normally convergent Fourier expansion
$F(w)=\sum_{k\in\mathbb Z^g}c_k e^{2\pi i(k+a)\cdot w}$.
One may obtain it as the Laurent expansion on $(\mathbb C^*)^g$
after multiplying by $e^{-2\pi ia\cdot w}$; equivalently its
coefficients are the preceding cube integrals, independent of
the imaginary slice by Cauchy's theorem.
Substitution of the $\tau n$-translation law yields
\[
 c_{k+n}=c_k\exp\{\pi i((k+n+a)^t\tau(k+n+a)
                  -(k+a)^t\tau(k+a))+2\pi in\cdot b\}.
 \tag{NFR.13}
\]
All coefficients are therefore one scalar multiple of the
coefficients in (NFR.11), and $F$ is that multiple of $\theta$.
It follows that every algebraic section is a scalar multiple of
$s_0$: equality on complex points is equality of regular functions
in each chart, because the fiber is smooth, reduced and of finite
type over $\mathbb C$.  Nonvanishing of $\theta$ and of a chart
frame proves $s_0\ne0$, and hence uniqueness of the scalar.

This proves exactly \texttt{Sec\allowbreak Gen} for every nonzero multiple
of $s_0$.  It agrees with TDR2's section-line description wherever
\texttt{Sec\allowbreak Base\allowbreak Change\allowbreak St} is supplied.  The argument here only uses
the already supplied actual $\texttt{Theta\allowbreak Data}$ and the complex
fiber, so it also proves the original conclusions stated with
\texttt{Sec\allowbreak Twist\allowbreak St} as a premise, without asserting that premise
for any other base or line bundle.
\end{proof}
% END NFR5

% BEGIN NFR6
\paragraph{NFR6: the whole actual normalized-point constructor.}
\label{claim:NFR6}
For every $b\in T.co^{-1}(\mathcal G)$ and every
$\kappa\in\mathbb C^\times$, there is an actual
$y:\texttt{Fr\allowbreak Pt}(X_o,c,CAt(\sigma))$ satisfying
$\texttt{Is\allowbreak Theta\allowbreak Pt}(T,b,\kappa,y)$.
Consequently there is a single function
$un:\mathcal G\to\texttt{Fr\allowbreak Pt}(X_o,c,CAt(\sigma))$
satisfying the full original \texttt{Is\allowbreak Norm\allowbreak Un} predicate:
for every $x\in\mathcal G$ some $\kappa_x$ satisfies
\[
 \kappa_x^2=C_0\det h(T.co^{-1}x),\qquad
 C_0=\sqrt{\frac{2^g}{P_0}},\qquad
 P_0=\prod_{j=1}^g2|\mathrm{ty}_j(\xi_0)|,
 \tag{NFR.14}
\]
and $\texttt{Is\allowbreak Theta\allowbreak Pt}(T,T.co^{-1}x,\kappa_x,un(x))$.

\begin{proof}
Put $y.x=F.u_k(P_{T.co(b)})$.  Its coefficient map is the one in
$F.\mathrm{over}_\sigma$.  Differentiate the actual exponential
in the coordinate directions $e_j$.  The two $F.\mathrm{ex}_{\rm diff}$
clauses make this differential a bijection from $\mathbb C^g$
onto the tangent space at the identity.  Additivity of the
exponential makes it complex linear; $F.\mathrm{ex}_{\rm lin}$
makes direction $j$ an eigenvector with eigenvalue
$\mathrm{ty}_j(a)$.  Under the original embedding (NFR.1) these
are precisely the ordered eigenbases in \texttt{Fr\allowbreak Pt}.
HUF1 identifies these analytic derivatives with the literal
dual-number predicate \texttt{Is\allowbreak Deriv\allowbreak At0}, and TDR1 identifies
linear combinations with its literal \texttt{Is\allowbreak Comb\allowbreak At}.

Transport the actual $s_0$ of NFR5 along the canonical isomorphism
between the two pullback models of the same fiber, and set
$y.s=\kappa s_0$.  Pullback does not change a chart value at a
point with the two specified projections; multiplication by
$\kappa$ multiplies that value by $\kappa$.
Thus its value is exactly
$\kappa\theta(T.ta(b),T.c_1,T.c_2,T.h(b)z)T.fr_j(b,z)$.
NFR5 proves the literal \texttt{Sec\allowbreak Gen} condition.  These data
fill every field of $y$ and every \texttt{Is\allowbreak Theta\allowbreak Pt} clause.

Each factor in $P_0$ is positive, since $\xi_0\ne0$ and field
embeddings are injective.  Thus $C_0>0$.  The determinant in
(NFR.14) is nonzero.  Choose a square root of its nonzero
product with $C_0$, and apply the preceding constructor.
Choice at every $x$ produces one function $un$.
No continuity of this square-root choice is required or claimed;
\texttt{Is\allowbreak Norm\allowbreak Un} has precisely this pointwise existential
quantifier.  This proves the entire original
\texttt{exists\_\allowbreak is\allowbreak Norm\allowbreak Un} conclusion, and not a prescribed-sign
or simply-connected-base variant.
\end{proof}
% END NFR6

\subsection{The two determinant identities and the Gaussian mean}

% BEGIN NFR7
\paragraph{NFR7: the normalized period matrix determines the Hermitian form.}
\label{claim:NFR7}
For any original $T$, put $P=T.Pf(T.co(b))$, $h=T.h(b)$,
$\tau=T.ta(b)$ and $Y=\operatorname{Im}\tau$.
If $H$ is a matrix with
$\operatorname{Im}H(Pa,Pb)=E_D(a,b)$, then
\[
 H=h^tY^{-1}\overline h,
 \qquad \det H\det Y=|\det h|^2.
 \tag{NFR.15}
\]

\begin{proof}
Let $A,B$ be the two column matrices of the prescribed integral
basis.  Then $hA=I$, $hB=\tau$, so $B=A\tau$.
Put $R=A^tH\overline A$.  The first-half pairing says
$\operatorname{Im}R=0$, while the mixed pairing says
$\operatorname{Im}(R\overline\tau)=-I$.
Thus $RY=I$ and $R=Y^{-1}$.  Conjugating back gives the first
identity.  Taking determinants and using
$\det\overline h=\overline{\det h}$ gives the second.
This argument does not silently replace $H$ by a different
polarization form; it determines the given form itself.
\end{proof}
% END NFR7

% BEGIN NFR8
\paragraph{NFR8: the exact determinant for unit-minor frames.}
\label{claim:NFR8}
Suppose $x\in\mathcal G$ has the source normalization
$|\operatorname{dmin}(G,x,\tau,i)|=1$ for every $\tau,i$.
Then every $H$ with the original imaginary period pairing satisfies
\[
                     \det H=P_0.                     \tag{NFR.16}
\]
This is the full original \texttt{Unit\allowbreak Det.\allowbreak det\_\allowbreak eq\_\allowbreak prod} conclusion.

\begin{proof}
Every leading Gram determinant is real, has its prescribed sign,
and has absolute value one.  Hence it equals the product of its
first $i$ signs.  In the unitriangular elimination (NFR.5), the
successive determinant ratios are therefore exactly $s_a$.
No diagonal rescaling is needed: $X=X_0b$ with each $b_\tau$
upper unitriangular.  Formula (NFR.6) shows that
$P_x=Q P_{x_0}$ with $\det Q=1$; in particular $|\det Q|=1$.
All identities hold after the exact permutation (NFR.3).

At $x_0$, (NFR.7) says the imaginary pairing is given by
\[
              H_0=\operatorname{diag}
                           (2|\mathrm{ty}_j(\xi_0)|)_{j=1}^g.
\]
The form $Q^tH\overline Q$ has that same imaginary pairing on
the period lattice of $x_0$.  Two complex sesquilinear forms
whose imaginary parts agree on a real spanning set agree
everywhere: extend by real bilinearity to all vectors, then
apply the equality also to $iz,w$ to recover their real parts.
Thus $Q^tH\overline Q=H_0$.
Taking determinants gives $|\det Q|^2\det H=\det H_0=P_0$,
as asserted.  The reverse order of the first $q$ columns and
the sign $-\xi_0$ in the trace were used in (NFR.7); dropping
either would invalidate the positive diagonal formula.
\end{proof}
% END NFR8

% BEGIN NFR9
\paragraph{NFR9: an actual fundamental domain and its exact theta mean.}
\label{claim:NFR9}
For $h\in\operatorname{GL}_g(\mathbb C)$, $\tau=X+iY$ with
$Y>0$ symmetric, and the exact lattice
$\mathcal L=h^{-1}(\mathbb Z^g+\tau\mathbb Z^g)$, define
\[
 F_d=\{h^{-1}(u+\tau v):u,v\in[0,1)^g\}.
\]
It is an additive fundamental domain for $\mathcal L$, has finite
positive volume, and for real $a,b$ (in particular the source's
half-characteristics),
\[
 \frac1{\operatorname{vol}F_d}\int_{F_d}
   |\theta[a,b](\tau,hz)|^2
    e^{-2\pi\operatorname{Im}(hz)^tY^{-1}\operatorname{Im}(hz)},dz
                       =\frac1{\sqrt{2^g\det Y}}.
 \tag{NFR.17}
\]

\begin{proof}
The real linear map $(u,v)\mapsto h^{-1}(u+\tau v)$ is invertible:
the imaginary part recovers $v$ using $Y^{-1}$, and then the
real part recovers $u$.  Its image of the half-open unit cube
tiles the whole space by the exact image of $\mathbb Z^{2g}$.
Its constant nonzero Jacobian is both $\operatorname{vol}F_d$
and the change-of-variables factor, so it cancels in the mean.

At $w=u+\tau v$, integrate first in $u$.
The normally convergent series from NFR5 permits termwise
integration of its square on the compact cube.
Fourier orthogonality eliminates every pair of different indices.
After multiplying by $e^{-2\pi v^tYv}$, the remaining summand is
\[
                       e^{-2\pi(n+a+v)^tY(n+a+v)}.
\]
All these terms are nonnegative; Tonelli's theorem permits their
integration and summation in either order.  The cubes
$n+a+[0,1)^g$ partition $\mathbb R^g$ up to their null boundaries.
Thus the double cube integral is
$\int_{\mathbb R^g}e^{-2\pi t^tYt}\,dt$.
Diagonalize the positive real matrix $Y$ orthogonally and scale
each coordinate; the product of the one-dimensional Gaussian
integrals is $2^{-g/2}(\det Y)^{-1/2}$.
For completeness, $\int_{\mathbb R}e^{-\pi t^2}\,dt=1$
follows by squaring the integral and using polar coordinates;
scaling gives the factors just used.
This proves (NFR.17), with the original Lebesgue measure and
normalization by the volume of this very domain.
\end{proof}
% END NFR9

% BEGIN NFR10
\paragraph{NFR10: the holomorphic gauge and its unit lattice multiplier.}
\label{claim:NFR10}
With $H=h^tY^{-1}\overline h$, set
\[
 q_0(z)=\frac\pi2(hz)^tY^{-1}(hz),\qquad
 \Theta(z)=e^{q_0(z)}\theta[a,b](\tau,hz).
\]
For $h\lambda=m+\tau n$ define
\[
 \chi(\lambda)=e^{\pi i(m\cdot n+2a\cdot m-2n\cdot b)}.
\]
Then $|\chi(\lambda)|=1$ and
\[
 \Theta(z+\lambda)=\chi(\lambda)
       e^{\pi H(z,\lambda)+\pi H(\lambda,\lambda)/2}\Theta(z),
 \tag{NFR.18}
\]
\[
 |e^{q_0(z)}|^2e^{-\pi H(z,z)}
             =e^{-2\pi\operatorname{Im}(hz)^tY^{-1}\operatorname{Im}(hz)}.
 \tag{NFR.19}
\]
Here $H(z,z)$ is real.

\begin{proof}
Write $w=hz$, $\mu=m+\tau n$ and $S=Y^{-1}$.
The relation $\overline\mu=\mu-2iYn$ gives
$w^tS\overline\mu=w^tS\mu-2iw\cdot n$.
Expand $q_0(z+\lambda)-q_0(z)$ and insert (NFR.12).
The $w$-dependent terms become $\pi w^tS\overline\mu$.
The constant difference from
$\pi\mu^tS\overline\mu/2$ is
$\pi i(m\cdot n+2a\cdot m-2n\cdot b)$.
This proves (NFR.18); its exponent is purely imaginary, proving
the unit norm.  For (NFR.19), write $w=u+iv$.
Then $\operatorname{Re}(w^tSw)=u^tSu-v^tSv$ and
$w^tS\overline w=u^tSu+v^tSv$.
Subtracting gives $-2v^tSv$, exactly the exponent in (NFR.19).
\end{proof}
% END NFR10

% BEGIN NFR11
\paragraph{NFR11: normalized square roots give mean exactly one.}
\label{claim:NFR11}
At a unit-minor frame $x$, let $\kappa$ be any root in (NFR.14).
Then
\[
 \frac{|\kappa|^2}{\sqrt{2^g\det Y}}=1,\qquad
 \frac1{\operatorname{vol}F_d}\int_{F_d}
  |\kappa\Theta(z)|^2e^{-\pi H(z,z)},dz=1.            \tag{NFR.20}
\]

\begin{proof}
NFR7--NFR8 give $P_0\det Y=|\det h|^2$.
Taking norms of $\kappa^2=C_0\det h$ gives
$|\kappa|^2=C_0|\det h|$.
Every factor is positive, so
\[
 C_0|\det h|=\sqrt{(2^g/P_0)|\det h|^2}
                        =\sqrt{2^g\det Y}.
\]
This proves the first equality, independently of the chosen sign
of $\kappa$.  Multiply (NFR.17) by $|\kappa|^2$ and use
(NFR.19) to obtain the second.
\end{proof}
% END NFR11

% BEGIN NFR12
\paragraph{NFR12: every literal clause of the original \texttt{Frame\allowbreak OK}.}
\label{claim:NFR12}
For every original $T:\texttt{Theta\allowbreak Data}(e,F,M,k)$ and every
function $un$ satisfying \texttt{Is\allowbreak Norm\allowbreak Un}, one has
\[
                  \forall x\in\mathcal G,\quad
                     \texttt{Frame\allowbreak OK}(e,F,M,k,x,un(x)).
 \tag{NFR.21}
\]
The statement retains exactly \texttt{Good\allowbreak Emb}, \texttt{Model\allowbreak Base}
and \texttt{Neg\allowbreak Xi}.  In particular it proves the original
\texttt{frame\allowbreak OK\_\allowbreak of\_\allowbreak is\allowbreak Norm\allowbreak Un} theorem with its printed
\texttt{Sec\allowbreak Twist\allowbreak St} antecedent, which the proof of this implication
does not need.

\begin{proof}
Fix $x$ and the $\kappa$ supplied by \texttt{Is\allowbreak Norm\allowbreak Un}; put
$b=T.co^{-1}x$ and $P=T.Pf(x)$.
The period formula, \texttt{Is\allowbreak Period}, equality of the moduli
point, and all prescribed derivatives are respectively $T.per$,
$T.isPer$ and the first two \texttt{Is\allowbreak Theta\allowbreak Pt} clauses.
This gives the unconditional part of \texttt{Frame\allowbreak OK}.

Now impose precisely its unit-minor antecedent.
Choose the positive Hermitian $H$ already supplied by
\texttt{Is\allowbreak Period}.  NFR7 identifies it, and NFR8 computes its
determinant.  The basis normalization in $T$ identifies its exact
period subgroup with $h^{-1}(\mathbb Z^g+\tau\mathbb Z^g)$:
write any lattice vector in the integral basis for one direction;
for the converse use that same basis and invertibility of $h$.
Thus the $F_d$ in NFR9 is a fundamental domain for $P.\mathrm{range}$,
as required by the source, not a finite-index sublattice.
Use the source chart frames and scalar
\[
 \widetilde{fr}_j(z)=T.fr_j(b,z)e^{-q_0(z)},\qquad
 \widetilde\theta(z)=\kappa e^{q_0(z)}\theta(T.ta(b),T.c_1,T.c_2,T.h(b)z).
 \tag{NFR.22}
\]
Define $\chi$ on the period lattice by NFR10 and arbitrarily
off it.  Uniqueness of $m,n$ follows from $Y>0$.

The matrix $H$ is Hermitian and positive and has exactly the
original imaginary period pairing.  Each $\widetilde{fr}_j$ is
holomorphic and nonzero on the exact inverse image of its
algebraic chart.  Multiplication by the common exponential
preserves every transition equation.  Since $q_0(0)=0$, its
value at zero is exactly the given rigidification value.
These verify the first seven displayed frame clauses.

For the translation clause, let $\lambda=Pa$.
The exponential satisfies $\mathrm{ex}(z+\lambda)=\mathrm{ex}(z)$,
using its actual kernel and group law.  Thus the scalar values
of the algebraic section $s_0$ at these two presentations of
the same fiber-product point agree.  Combining that equality
with (NFR.18) yields
\[
 \widetilde{fr}_j(z)=\chi(\lambda)
    e^{\pi H(z,\lambda)+\pi H(\lambda,\lambda)/2}
                           \widetilde{fr}_j(z+\lambda).
 \tag{NFR.23}
\]
Cancellation first proves this away from the zeros of $\Theta$.
Those zeros have empty interior by NFR5, so continuity extends
it on the whole open chart.  The translated chart is the same
open set because the exponential points are equal.  Thus this
extension is legitimate even where the theta section vanishes.
NFR10 also gives $|\chi(Pa)|=1$.

The last \texttt{Is\allowbreak Theta\allowbreak Pt} clause and
$e^{q_0(z)}e^{-q_0(z)}=1$ give, for every original chart and
every fiber-product point with the two prescribed projections,
the exact identity
$\operatorname{val}(un(x).s)=\widetilde\theta(z)
\widetilde{fr}_j(z)$.  Finally NFR9 gives the literal additive
fundamental-domain condition and NFR11 gives its mean equal to
one.  These are all remaining clauses of \texttt{Frame\allowbreak OK}.
Nothing in this argument changes the section of $un(x)$ or its
ordered tangent frame.
\end{proof}
% END NFR12

% BEGIN NFR13
\paragraph{NFR13: simultaneous actual choices on the finite family.}
\label{claim:NFR13}
Under the original input scope above, choose one $T_k$ for each
component and one function $un_k$ such that
\[
  \forall k,\quad\texttt{Is\allowbreak Norm\allowbreak Un}(T_k,un_k),\qquad
  \forall k\ \forall x\in M.X_k.\mathrm{dom},\quad
                       \texttt{Frame\allowbreak OK}(e,F,M,k,x,un_k(x)).
 \tag{NFR.24}
\]
With the existing \texttt{Sec\allowbreak Base\allowbreak Change\allowbreak St} premise of TDR13,
every one of these very framed points is \texttt{From\allowbreak Test}.

\begin{proof}
NFR4 supplies each $T_k$, then NFR6 supplies its $un_k$.
Choose them once for the finite component set, before the
universal frame variable.  NFR12 applies to every chosen
$un_k$; it is valid for any choices of the pointwise square roots.
TDR13 is universal over actual complex \texttt{Fr\allowbreak Pt} values,
so apply it directly to $un_k(x)$, preserving its moduli point,
its ordered eigenbasis and its actual generating section.
The finitely presented reduced test ring and its frame may depend
on $k,x$, exactly as the original \texttt{From\allowbreak Test} permits.
\end{proof}

The full original \texttt{Frames\allowbreak OK} is a conjunction of four
fields.  Equation (NFR.24) supplies its first field, on one
simultaneous family.  Its other three fields need the
transformation of theta squares under the exact integral
level group, the resulting triangular character identity,
and holomorphic evaluation of the original algebraic section
rules.  They are not consequences of choosing a square root
or of the mean-one calculation.  Likewise the independent
\texttt{Cover\allowbreak Tri}, \texttt{Fib\allowbreak Tri} and complete \texttt{Is\allowbreak Unf}
assembly remain open in this packet.
% END NFR13

\paragraph{Source map and inherited scope.}
The exact original contracts and source ranges are recorded in
this component's \texttt{contracts.\allowbreak json}, \texttt{provenance.\allowbreak json}
and \texttt{source\_\allowbreak manifest.\allowbreak json}.
The main source bodies are:\par
\noindent\texttt{Norm\allowbreak Const.\allowbreak lean};\par
\noindent\texttt{Nonempty\allowbreak Theta\allowbreak Data\allowbreak Flag\allowbreak Domain\allowbreak Models.\allowbreak lean};\par
\noindent\texttt{Exists\allowbreak Theta\allowbreak Pt.\allowbreak lean};\par
\noindent\texttt{Exists\allowbreak Normalised\allowbreak Framed\allowbreak Point\allowbreak On\allowbreak Each\allowbreak Component.\allowbreak lean};\par
\noindent\texttt{Herm\allowbreak Eq\allowbreak Of\allowbreak Normalised.\allowbreak lean};
\texttt{Eq\allowbreak Of\allowbreak Im\allowbreak Herm\allowbreak Form\allowbreak Eq.\allowbreak lean};\par
\noindent\texttt{Exists\allowbreak Fundamental\allowbreak Domain\allowbreak Mean\allowbreak Norm\allowbreak Sq\allowbreak Theta\allowbreak Char.\allowbreak lean};\par
\noindent\texttt{Exp\allowbreak Mul\allowbreak Theta\allowbreak Char\allowbreak Automorphy.\allowbreak lean};\par
\noindent\texttt{Normalised\allowbreak Points\allowbreak Satisfy\allowbreak Single\allowbreak Frame.\allowbreak lean}.\par
The proofs use HUF1, HUF20--HUF21, PUN11 and the exact TDR
frame dictionary as indicated; their source hashes are separately
bound.  No assertion about existence of the original representing
moduli scheme or the remaining global frame laws is added.


\section{Uniform character seeds and the common abelian CM field}
\label{sec:common-arithmetic-frame-v50}

This section constructs the character-and-field input to the ordinary
automorphy argument.  Its principal endpoint is the full original
\texttt{seed\_\allowbreak cl}: one modulus is chosen before the later abelian CM
field and auxiliary prime, and every compatible later choice receives
the entire original \texttt{Seed\allowbreak Clause}.  We also construct the actual
fixed field in \texttt{field\_\allowbreak ex}, the auxiliary prime in
\texttt{auxw\_\allowbreak ex}, and install the seed on an already supplied common
frame without changing its marked element.  The seed here is a triple
of characters.  Constructing the rank-two and rank-three automorphic
representations from that triple remains a separate task.

\subsection{Objects, hypotheses, and the order of choices}

Fix the supplied algebraic closure \(\Omega/\mathbb Q\), put
\(G=\operatorname{Gal}(\Omega/\mathbb Q)\), and let \(\ell\) be an
odd prime.  The coefficient ring \(O=c.O\subset\overline{\mathbb Q}_\ell\)
is the original finite \(\mathbb Z_\ell\)-algebra which is a DVR.
Write \(k=O/\mathfrak m_O\), \(q=|k|\), and
\[
 \epsilon:G\longrightarrow\mathbb Z_\ell^\times,\qquad
 \omega:G\longrightarrow k^\times,
 \qquad \alpha=\overline\Psi:G\longrightarrow k^\times .
 \tag{CAF.1}
\]
Here \(\omega\) is the reduction of the actual cyclotomic character,
and \(\Psi:G\to O^\times\) and
\(\eta:G\to\{\pm1\}\) are the specified characters.  A complex
conjugation \(\sigma\in G\) and a valuation \(B\) over \(\ell\)
are fixed.  For a positive integer \(n\), the source's
\(\operatorname{KillsMu}(n,\chi)\) means
\[
 s|_{\mu_n}=1\quad\Longrightarrow\quad \chi(s)=1.
 \tag{CAF.2}
\]
It is a factorization through the actual cyclotomic extension inside
\(\Omega\), rather than an assertion about an unrelated character of
the same order.  Decomposition and inertia subgroups at a valuation
\(A\) are denoted by \(D_A,I_A\).

The full original \(\operatorname{Gen}(\ell,\eta,\Psi,\sigma,B,M,H_b)\)
is an allowed starting input.  For the seed construction only the
following clauses of that record are needed:
\begin{enumerate}
\item \(k\) is finite, \(M>0\), \(\ell\mid M\), \(\ell\ne2\),
      \(\sigma\) is a complex conjugation, and \(B\) lies over \(\ell\).
\item There are \(D>0\) and \(\delta\in\Omega\) such that
      \(\delta^2=-D\) and \(s\delta=\eta(s)\delta\) for every \(s\in G\).
\item For one integer \(a\ge0\), both \(\eta\) and \(\alpha\)
      satisfy \(\operatorname{KillsMu}(M^a,-)\).
\item \(\eta|_{I_A}=1\) for every \(A\) over \(\ell\),
      \(\alpha(\sigma)=1\), and there is one \(j\ge0\) such that
      \(\alpha|_{I_A}=\omega^j|_{I_A}\) for every such \(A\).
\item \(\ker\omega\cap\ker\alpha\not\subseteq\ker\eta\).
\end{enumerate}
Call these explicit clauses \(\mathcal S\).  Every use of
\(\mathcal S\) below is therefore available from the original
\texttt{Gen} without an extra hypothesis.  Conversely, proving the
seed from these fewer clauses does not assert the other
\texttt{Gen} fields.  In particular, \(H_b\), the local degree bounds,
and the strict cohomology classes are not constructed by the seed
argument.  Claims CAF9 and CAF10 verify \(\mathcal S\) directly for
the original \texttt{Dat} and \texttt{Host\allowbreak Data} applications.

For an abelian CM field \(L\subset\Omega\), let
\(U_L=G_{L^+}\), the source's \texttt{upl}, and
\(G_L=\operatorname{Gal}(\Omega/L)\).  Put
\[
 H_\ell(w)=\{s\in G:\ s|_{\mu_w}\text{ is an }\ell
                  \text{-th power in }(\mathbb Z/w\mathbb Z)^\times\}.
 \tag{CAF.3}
\]
This equals the literal \texttt{Hth}: the source additionally writes
the same condition for \(s^{-1}\), which is automatic because the
\(\ell\)-th powers form a subgroup.

A \emph{witness prime} is exactly the original
\texttt{Frobenius\allowbreak Witness\allowbreak Prime}: a prime \(w\nmid M\) with
\(w\equiv1\pmod\ell\), such that for every \(A\) above \(w\),
\[
          D_A\subseteq\ker\eta\cap\ker\alpha\cap\ker\omega.
 \tag{CAF.4}
\]
All assertions below use the same \(\Omega,O,\eta,\Psi,\sigma\).
The embedding \(\iota:\overline{\mathbb Q}_\ell\simeq\mathbb C\)
needed for later automorphy is not used or changed here.

\subsection{Local algebra and the two cyclotomic constructions}

% BEGIN CAF1
\paragraph{CAF1. Teichmuller lifts with the original norm congruence.}
\label{claim:CAF1}
There is an injective homomorphism
\(T:k^\times\to\overline{\mathbb Q}_\ell^\times\), landing in
\(O^\times\), such that for every \(u\in O^\times\), every
\(j\ge0\), and every \(s\in G\),
\[
 \left|\epsilon(s)^j
       T\bigl(\overline u\,\omega(s)^{-j}\bigr)-u\right|_\ell<1.
 \tag{CAF.5}
\]
Moreover, \(|x|_\ell\le1\) for every \(x\in O\), and
\(x\in\mathfrak m_O\) if and only if \(|x|_\ell<1\).
These assertions need only that \(\ell\) is prime, including
\(\ell=2\); the parity observation at the end of the proof uses
oddness.

\emph{Proof.}
Integrality over \(\mathbb Z_\ell\) gives the first norm bound: in a
monic integral equation, a root of norm greater than one would make
the leading term strictly larger than every other term, contradicting
the ultrametric inequality.  A unit and its inverse both have norm at
most one, so both have norm one.  If \(\varpi\) is a uniformizer,
write \(\ell=u\varpi^r\) in the DVR.  Since
\(|\ell|_\ell=\ell^{-1}<1\) and \(|u|_\ell=1\),
\(|\varpi|_\ell<1\).  The factorization of any nonzero element as a
unit times a power of \(\varpi\) now proves the asserted equivalence,
including zero separately.

The ring \(O\) is complete: it is finite over the complete DVR
\(\mathbb Z_\ell\), and its maximal-ideal topology is the same as
its \(\ell\)-adic topology.  For \(b\in k^\times\), Hensel's lemma
gives a unique root \(T(b)\in O\) of \(X^{q-1}-1\) reducing to
\(b\), since the derivative is a unit modulo \(\mathfrak m_O\).
Uniqueness shows \(T(bb')=T(b)T(b')\), and reduction proves
injectivity.  The reduction of
\(\epsilon(s)^jT(\overline u\,\omega(s)^{-j})\) is precisely
\(\overline u\).  Its difference from \(u\) belongs to
\(\mathfrak m_O\), proving (CAF.5).  Since \(\ell\) is odd,
\(-1\ne1\) in \(k\), and uniqueness also gives \(T(-1)=-1\).
This proves the complete original \texttt{sc\_\allowbreak norm} and
\texttt{sc\_\allowbreak teich} conclusions.
% END CAF1

% BEGIN CAF2
\paragraph{CAF2. Removing the cyclotomic conductor at \(\ell\).}
\label{claim:CAF2}
Let \(C\) be any group and \(\chi:G\to C\) a homomorphism.
If \(N_0>0\), \(\ell\nmid N_0\),
\(\operatorname{KillsMu}(\ell^rN_0,\chi)\), and
\(\chi|_{I_B}=1\), then
\(\operatorname{KillsMu}(N_0,\chi)\).
If also \(w\equiv1\pmod {N_0}\), then the source's full
\(\operatorname{SplitAt}(N_0,w,\chi)\) holds.

\emph{Proof.}
Inertia at \(\ell\) acts trivially on roots of unity of order prime
to \(\ell\).  Indeed, those roots are the Hensel lifts of the roots
in finite residue-field extensions, and hence lie in the maximal
unramified extension.  In contrast,
\(\mathbb Q_\ell(\mu_{\ell^r})/\mathbb Q_\ell\) is totally
ramified: for \(r\ge1\),
\(\Phi_{\ell^r}(1+X)\) is Eisenstein at \(\ell\), of degree
\(\ell^{r-1}(\ell-1)\).  Inertia therefore surjects onto the full
group of its cyclotomic automorphisms.  The case \(r=0\) is immediate.

Suppose \(s\) fixes \(\mu_{N_0}\).  Choose \(t\in I_B\) with the
same action as \(s\) on \(\mu_{\ell^r}\).  It fixes
\(\mu_{N_0}\), so \(t^{-1}s\) fixes both sets of roots.  Since the
orders are coprime, these roots generate \(\mu_{\ell^rN_0}\).
Thus \(\chi(t^{-1}s)=1\), and \(\chi(t)=1\) gives \(\chi(s)=1\).
No commutativity of \(C\) was used.  Finally, an element acting as
\(\zeta\mapsto\zeta^w\) on \(\mu_{N_0}\) fixes it when
\(w\equiv1\pmod {N_0}\), so \(\chi\) is one on every such
element.  These are exactly the two \texttt{Split\allowbreak At} clauses.
The proof preserves the one given valuation \(B\) and supplies the
original \texttt{sc\_\allowbreak strip} for every target group.
% END CAF2

% BEGIN CAF3
\paragraph{CAF3. The actual imaginary quadratic field and its discriminant.}
\label{claim:CAF3}
The quadratic clause of \(\mathcal S\) alone makes the subfield
\(K_0=\mathbb Q(\delta)\) of \(\Omega\) imaginary quadratic with
\(G_{K_0}=\ker\eta\).  With the inertia clause as well,
\[
                 G_{K_0}=\ker\eta,\qquad
                 \ell\nmid\operatorname{disc}(K_0).
 \tag{CAF.6}
\]
More generally, the discriminant conclusion holds for any supplied
quadratic field \(K\) for which every inertia subgroup above
\(\ell\) lies in \(G_K\).  No oddness assumption is needed for
this assertion.

\emph{Proof.}
The polynomial \(X^2+D\) has no rational root because \(D>0\).
It is therefore the minimal polynomial of \(\delta\), and all its
complex roots are nonreal.  An element \(s\in G\) fixes \(K_0\)
exactly when it fixes \(\delta\).  As \(\delta\ne0\) and
\(s\delta=\eta(s)\delta\), this is equivalent to \(\eta(s)=1\).
Every inertia subgroup above \(\ell\) lies in that kernel by
\(\mathcal S\), so \(K_0\) is unramified above \(\ell\).

For completeness, the discriminant consequence uses the ordinary
integral trace pairing.  The completion of the ring of integers at
\(\ell\) is a product of the integer rings of the local extensions.
Trivial inertia makes each of these unramified; its residue extension
is a finite separable field extension.  Reduction of its trace pairing
is the nondegenerate field-trace pairing over the finite residue field.
Its determinant is consequently a unit.  Taking the product over the
places shows that the global discriminant is a unit at \(\ell\).
This is the source's \texttt{sc\_\allowbreak quad} and \texttt{sc\_\allowbreak disc}, on the
actual embedded quadratic field.
% END CAF3

% BEGIN CAF4
\paragraph{CAF4. The prescribed order-\(\ell\) character at \(w\).}
\label{claim:CAF4}
For primes \(\ell,w\) with \(w\equiv1\pmod\ell\), there is a
character \(\theta:G\to\overline{\mathbb Q}_\ell^\times\) such that
\[
 \ker\theta=H_\ell(w),\qquad \theta\ne1,
 \qquad\theta^\ell=1,
 \qquad\operatorname{KillsMu}(w,\theta).
 \tag{CAF.7}
\]
The group \(H_\ell(w)\) is open, contains \([G,G]\), has index
\(\ell\), contains all inertia away from \(w\), and does not
contain inertia at any place over \(w\).  These assertions hold also
for \(\ell=2\); only the further real-place observation below uses
oddness.

\emph{Proof.}
The cyclotomic map \(G\twoheadrightarrow
(\mathbb Z/w\mathbb Z)^\times\) is onto; the target is cyclic of
order \(w-1\).  The subgroup of \(\ell\)-th powers has index
\(\ell\).  Compose this map, its quotient by the power subgroup,
and an isomorphism of that cyclic quotient with
\(\mu_\ell(\overline{\mathbb Q}_\ell)\).  This gives \(\theta\)
and all four identities in (CAF.7).  It also proves openness and
the commutator assertion.  Roots of order \(w\) are unramified away
from \(w\), while total ramification of
\(\mathbb Q_w(\mu_w)\) makes inertia at \(w\) surject onto the
entire cyclotomic group.  Its image in the order-\(\ell\) quotient
is therefore nontrivial at every prolongation of the place.  For odd
\(\ell\), complex conjugation belongs to \(H_\ell(w)\), because
its cyclotomic image \(-1\) is \((-1)^\ell\).
This proves \texttt{hth} and \texttt{sc\_\allowbreak theta}, including the
literal kernel equality needed in the later field choice.
% END CAF4

% BEGIN CAF5
\paragraph{CAF5. A witness prime satisfying every prescribed congruence.}
\label{claim:CAF5}
Given the original \texttt{Gen} and any \(Q>0\), there is a witness
prime \(w\) satisfying \(w\equiv1\pmod Q\).  In fact, the
weaker clauses \(\mathcal S\) suffice.

\emph{Proof.}
Choose \(a\) in \(\mathcal S\) and put \(R=M^{a+1}Q\).  There
are arbitrarily large primes congruent to one modulo \(R\).
Only this residue class is needed, so the elementary cyclotomic proof
suffices: for \(R>1\), take a sufficiently large integer \(A\)
divisible by \(R\)
and by any proposed finite list of primes.  A prime divisor of
\(\Phi_R(A)>1\) divides neither \(A\) nor \(R\), since
\(\Phi_R(0)=1\).  The order of \(A\) modulo that prime is exactly
\(R\): the cyclotomic factors of \(X^R-1\) have disjoint roots
in characteristic prime to \(R\).  Hence the prime is one modulo
\(R\) and is outside the proposed list.  The case \(R=1\) is
Euclid's argument.  Select such a prime \(w>R\).

Then \(w\nmid M\), \(w\equiv1\pmod\ell\), and
\(w\equiv1\pmod Q\).  For \(A\) above \(w\), decomposition
acts on \(\mu_n\), \(n\mid R\), through powers of the Frobenius
\(\zeta\mapsto\zeta^w\).  These roots are unramified there, and
\(w\equiv1\pmod n\), so the entire decomposition subgroup fixes
them.  In particular it fixes \(\mu_{M^a}\) and \(\mu_\ell\).
The two \texttt{Kills\allowbreak Mu} conditions and the definition of \(\omega\)
give (CAF.4).  This proves \texttt{auxw\_\allowbreak ex} without a stronger
distribution theorem for primes or a field-dependent modulus.
% END CAF5

\subsection{One modulus before every later field and prime}

% BEGIN CAF6
\paragraph{CAF6. Fixed core characters and a prime-to-\(\ell\) modulus.}
\label{claim:CAF6}
Choose \(a,j,\delta\) from \(\mathcal S\), put
\[
 M^a\ell=\ell^rN_0,\quad \ell\nmid N_0,\quad N_0>0,
 \qquad \phi=\alpha\omega^{-j},\quad
 \chi'=T\phi,\quad \eta'=\eta\text{ viewed in }
                  \overline{\mathbb Q}_\ell^\times.
 \tag{CAF.8}
\]
All of \(K_0,N_0,j,\chi',\eta'\) are fixed before \(L,w\).
They satisfy
\[
 \begin{gathered}
 \operatorname{KillsMu}(N_0,\chi'),\qquad
 \operatorname{KillsMu}(N_0,\eta'),\qquad
 \ell\nmid\operatorname{ord}(\chi'),\qquad \chi'\ne\eta',\\
 \ker\eta'=G_{K_0},\qquad \chi'(\sigma)=(-1)^j,\qquad
 |(\epsilon^j\chi')(s)-\Psi(s)|_\ell<1\quad(s\in G).
 \end{gathered}
 \tag{CAF.9}
\]

\emph{Proof.}
The characters \(\eta,\alpha\) factor through \(\mu_{M^a}\),
and \(\omega\) factors through \(\mu_\ell\).  Thus \(\chi'\)
and \(\eta'\) satisfy \(\operatorname{KillsMu}(M^a\ell,-)\).
At every inertia group over \(\ell\),
\(\phi=\alpha\omega^{-j}=1\), and \(\eta'=1\).
CAF2 strips the factor \(\ell^r\) and proves the two factorizations
through \(\mu_{N_0}\).

The image of \(\chi'\) lies in the Teichmuller group of order
\(q-1\).  Its order divides \(q-1\), which is prime to \(\ell\).
The natural embedding of \(\{\pm1\}\) in the characteristic-zero
field is injective, so \(\ker\eta'=\ker\eta=G_{K_0}\).
At \(\sigma\), \(\alpha=1\) and \(\omega=-1\), giving
\(\chi'(\sigma)=T((-1)^{-j})=(-1)^j\).  Equation (CAF.5), with
\(u=\Psi(s)\), gives the last assertion of (CAF.9) for every
\(s\), with the original integral embedding of \(O\).

Finally, choose \(s\in\ker\omega\cap\ker\alpha\) outside
\(\ker\eta\), using \(\mathcal S\).  At this element
\(\chi'(s)=T(1)=1\) and \(\eta'(s)=-1\).  Hence
\(\chi'\ne\eta'\).  No later field or representation has been
chosen in this proof.
% END CAF6

% BEGIN CAF7
\paragraph{CAF7. The complete original uniform \texttt{Seed\allowbreak Clause}.}
\label{claim:CAF7}
Set \(Q=N_0\) from CAF6.  For every later abelian CM field
\(L\subset\Omega\) and every later witness prime \(w\), assume
\[
 G_L=U_L\cap\ker\eta,\qquad w\equiv1\pmod Q,
 \qquad U_L\subseteq H_\ell(w),\qquad
 \exists s_0\in G_L:\alpha(s_0)^2\ne1.
 \tag{CAF.10}
\]
Then the full original
\(\operatorname{SeedClause}(c,\eta,\Psi,\sigma,L)\) holds.
Consequently this proves \texttt{seed\_\allowbreak cl} for every original
\texttt{Gen}, with its entire \(\forall L\,\forall w\) scope.

\emph{Proof.}
Use \(K_0,N_0,j,\chi',\eta'\) already fixed in CAF6, and choose
\(\theta\) from CAF4 for this \(w\).  The numerical conditions in
\texttt{Slope\allowbreak Test\allowbreak Remaining\allowbreak Choices} are all satisfied: \(\ell\ne2\),
\(N_0>0\), \(\ell\nmid N_0\), \(w\) is prime, and
\(w\nmid N_0\).  For the last assertion, \(N_0\mid M^a\ell\),
whereas \(w\nmid M\) and \(w\ne\ell\); the latter follows from
\(w\equiv1\pmod\ell\).  The remaining three auxiliary conditions
\(\theta\ne1\), \(\theta^\ell=1\), and
\(\operatorname{KillsMu}(w,\theta)\) are exactly (CAF.7).

We check every field condition.  The given \(\sigma\) is a complex
conjugation.  CAF3 gives \([K_0:\mathbb Q]=2\) and imaginary
quadratic \(K_0\).  Abelianity of \(L\) is supplied in (CAF.10).
The equality there gives
\(G_L\subseteq\ker\eta=G_{K_0}\), which is the original
\texttt{KL} inclusion, and CAF6 gives
\(\ker\eta'=G_{K_0}\).  Since \(w\equiv1\pmod {N_0}\),
CAF2 and CAF6 give both
\(\operatorname{SplitAt}(N_0,w,\eta')\) and
\(\operatorname{SplitAt}(N_0,w,\chi')\).  Furthermore,
\[
        G_L\subseteq U_L\subseteq H_\ell(w)=\ker\theta,
\]
which is the original \texttt{theta\allowbreak L} condition, and the parity
condition is \(\chi'(\sigma)=(-1)^j\) from CAF6.

For the remaining square condition, let
\(z=(\epsilon^j\chi')(s_0)\) and \(u=\Psi(s_0)\).
Equation (CAF.9) gives \(|z-u|_\ell<1\).  Since \(u\) is a
unit, both \(|u|_\ell=1\) and \(|z|_\ell=1\).
If \(|z^2-1|_\ell<1\), then
\[
 u^2-1=(u-z)(u+z)+(z^2-1)
\]
has norm less than one.  CAF1 would give
\(\alpha(s_0)^2=1\), contrary to (CAF.10).  Therefore
\(\neg\operatorname{Cong}(z^2,1)\), precisely the source's
\texttt{sq} clause.  This proves every field of
\texttt{FieldPair\allowbreak{}Character\allowbreak{}TripleConditions}.

The four additional \texttt{Seed\allowbreak Clause} requirements are
\(\chi'\ne\eta'\), \(\ell\nmid\operatorname{ord}(\chi')\),
the literal equality of values of \(\eta'\) with the integer-valued
\(\eta\), and the all-\(s\) norm congruence.  They are (CAF.8)--(CAF.9).
The final discriminant requirement \(\ell\nmid\operatorname{disc}K_0\)
is CAF3.  Nothing in these choices depends on a later automorphic
representation, coefficient extension, or universal-ring point.
The proof of the square implication also proves the original
\texttt{sc\_\allowbreak cong} on its full stated inputs.
% END CAF7

\subsection{The actual CM field, and the two arithmetic applications}

% BEGIN CAF8
\paragraph{CAF8. An explicit common abelian CM field from the prescribed subgroup.}
\label{claim:CAF8}
Let \(U\subseteq G\) be open, contain \([G,G]\), and contain
every complex conjugation.  Under the quadratic clause of
\(\mathcal S\), define the actual subfield
\[
                   L=\Omega^{U\cap\ker\eta}.
 \tag{CAF.11}
\]
Then \(L/\mathbb Q\) is finite abelian and CM, and literally
\[
             G_L=U\cap\ker\eta,\qquad U_L=U.
 \tag{CAF.12}
\]
Thus the original \texttt{field\_\allowbreak ex} is proved for every original
\texttt{Gen}, every \(m\), and every supplied
\(\operatorname{UFacts}(\ell,\eta,\Psi,M,m,U)\).

\emph{Proof.}
A subgroup containing the commutator is normal, and its quotient is
abelian.  Both \(U\) and \(\ker\eta=G_{K_0}\) are open, so
infinite Galois correspondence identifies their fixed fields with
finite Galois extensions in the given \(\Omega\).  Put
\(F=\Omega^U\).  Every complex embedding of \(F\) extends to an
embedding of \(\Omega\) in \(\mathbb C\).  The corresponding
complex conjugation lies in \(U\), so fixes \(F\).  Hence every
such embedding has real image and \(F\) is totally real.

The fixed field of an intersection of these closed subgroups is the
compositum, so \(L=FK_0\) inside \(\Omega\).  Since the imaginary
quadratic field \(K_0\) has no nontrivial real subfield,
\(F\cap K_0=\mathbb Q\).  Thus \([L:F]=2\), and every complex
embedding of \(L\) contains a nonreal root of \(X^2+D\).
The extension \(L/F\) is totally imaginary quadratic, so it is CM
with maximal real subfield exactly \(F\).  Consequently
\(U_L=G_F=U\), and the first equality in (CAF.12) follows directly
from the definition (CAF.11).  Both groups and their embeddings are
the original subgroups of \(G\); there is no replacement by an
abstract isomorphic CM field.  In \texttt{UFacts}, precisely
\texttt{is\allowbreak Open}, \texttt{ab}, and \texttt{real} provide the three
subgroup hypotheses used here.  Its remaining local clauses are
retained for subsequent frame construction.
% END CAF8

% BEGIN CAF9
\paragraph{CAF9. The seed hypotheses for the original \texttt{Dat}.}
\label{claim:CAF9}
For the original \(D:\operatorname{Dat}(\Omega,K,c,e,\psi)\),
use its characters \(\eta=D.\eta\), \(\Psi=\epsilon^e\psi\),
its conjugation \(\sigma=D.\mathrm{cc}\), and any supplied \(B\)
over \(\ell\).  Set
\[
 M=\ell N|\operatorname{disc}K|,
 \qquad Q=|\operatorname{disc}K|\,N_{\ell'},
 \tag{CAF.13}
\]
where \(N_{\ell'}\) is the prime-to-\(\ell\) part of \(N\).
Then \(\mathcal S\) holds, and CAF7 supplies the original
\texttt{Seed\allowbreak Clause} for every later \(L,w\) satisfying (CAF.10)
with this \(Q\).  The quadratic field in that clause may be the
original embedded \(K\).

\emph{Proof.}
The original \(GD\) conditions give imaginary quadratic \(K\),
\(\ell\ne2\), even positive \(N\),
\(\ell\nmid\operatorname{disc}K\), prime-to-\(\ell\) order of
\(\psi\), \((\operatorname{comp}_\ell\psi)^2=1\), and
\(\psi(-1)=(-1)^e\).  The ring \(O\) is finite over
\(\mathbb Z_\ell\), so its residue field is finite.
Write \(\Delta=\operatorname{disc}K<0\).  The discriminant
determinant of an integral basis is an element \(\delta\in K\)
with \(\delta^2=\Delta=-|\Delta|\).  It generates the quadratic
field, and its two conjugates are \(\delta,-\delta\).  Since
\(\ker\eta=G_K\), \(s\delta=\eta(s)\delta\).

The elementary quadratic conductor theorem places
\(K\) in \(\mathbb Q(\mu_{|\Delta|})\), so
\(\operatorname{KillsMu}(|\Delta|,\eta)\).  One usual explicit
proof uses the primitive quadratic Kronecker character \(\chi_\Delta\):
its Gauss sum \(g=\sum_{a\bmod |\Delta|}\chi_\Delta(a)\zeta^a\)
satisfies \(s_b(g)=\chi_\Delta(b)g\) and \(g^2=\Delta\).
The latter identity follows from finite additive-character
orthogonality, giving \(|g|^2=|\Delta|\), and
\(\overline g=\chi_\Delta(-1)g=-g\).
Thus the same quadratic field generated by \(\delta\) is contained
in that cyclotomic field.  This is the usual elementary quadratic
conductor input; no general potential modularity is used.
Also \(\psi\) factors through \(\mu_N\), and \(\omega\) through
\(\mu_\ell\).  Hence \(\eta\) and
\(\alpha=\omega^e\overline\psi\) both satisfy
\(\operatorname{KillsMu}(M,-)\), so \(a=1\) is allowed.

The character \(\eta\) is unramified at \(\ell\).
At \(\ell\), the residual character of the finite local component
of \(\psi\) factors through \(\mathbb F_\ell^\times\): the
\(\ell\)-Sylow subgroup maps trivially into the prime-to-\(\ell\)
group \(k^\times\).  Therefore it is \(\omega^b\) for one
\(0\le b<\ell-1\).  The same exponent works at every prolongation
of the place, by conjugacy and commutativity of the target.  Taking
\(j=e+b\) proves the inertia clause, and the stated parity gives
\(\alpha(\sigma)=1\).

It remains to verify the kernel noncontainment.  The original
\(GD.\eta\) gives an \(s\) with \(\omega(s)=1\),
\(\chi_P(s)=1\), and \(\eta(s)=-1\), where, since \(\psi\)
has order prime to \(\ell\),
\(\chi_P=\psi(\operatorname{comp}_\ell\psi)^{-1}\).
The local component has order prime to \(\ell\), and therefore
factors through the mod-\(\ell\) cyclotomic character rather than
its pro-\(\ell\) kernel.  It is one at this \(s\).
Consequently \(\psi(s)=1\), and \(\alpha(s)=1\), as required.
CAF6 with \(a=1\) yields
\(Q=(M\ell)_{\ell'}=|\Delta|N_{\ell'}\), exactly (CAF.13).
All constructions precede \(L,w\).  Taking the original \(K\)
in CAF3 instead of the canonically isomorphic \(\mathbb Q(\delta)\)
preserves its original embedding and all the field conditions.
This claim verifies the seed-relevant part of \texttt{gen\_\allowbreak gd}; it
does not claim the unused full \texttt{Gen} fields as new results.
% END CAF9

% BEGIN CAF10
\paragraph{CAF10. The separate CountV host has its own seed modulus.}
\label{claim:CAF10}
For the original \(\operatorname{HostData}(\ell,e,\psi,o_1)\),
take
\[
 \eta=\eta_4,\qquad \Psi=\epsilon^e\psi,
 \qquad M=\ell N,
 \qquad Q=N_{\ell'}^{\,2}.
 \tag{CAF.14}
\]
For every complex conjugation \(\sigma\) and every \(B\) over
\(\ell\), these data satisfy \(\mathcal S\).  Thus CAF7 supplies
\texttt{Seed\allowbreak Clause} for every compatible later \(L,w\), with its
original host characters.  This assertion includes \(\ell=3\).

\emph{Proof.}
The source \texttt{Host\allowbreak Data} has \(\ell\ne2\), \(e,o_1\ge1\),
even \(N>0\), the parity \(\psi(-1)=(-1)^e\), and dyadic
conductor exponent \(c_2(\psi)\ge4\).  Its additional order and
size inequalities are retained but are unnecessary for this seed
construction.  The character \(\eta_4\) is the character of the
actual quadratic field \(\mathbb Q(i)\subset\Omega\); choose
\(\delta=i\), so \(D=1\).  It is unramified at odd \(\ell\)
and factors through \(\mu_4\).
Both \(\eta_4\) and \(\alpha\) factor through
\(\mu_{M^2}\), because \(4\mid M^2\), \(N\mid M^2\), and
\(\ell\mid M^2\).  Take \(a=2\).
The residue field is finite because \(O\) is finite over
\(\mathbb Z_\ell\), and the stated parity gives
\(\alpha(\sigma)=1\).  On inertia at \(\ell\), the residual
character of \(\psi\) kills the \(\ell\)-Sylow subgroup of
\((\mathbb Z/\ell^r\mathbb Z)^\times\), because \(k^\times\)
has order prime to \(\ell\).  It is therefore a power
\(\omega^b\), with one \(0\le b<\ell-1\), and the exponent
\(j=e+b\) works at every prolongation by conjugacy.  Thus no
prime-to-\(\ell\) order assumption on \(\psi\) is required.

Here is the needed noncontainment without using the
\(GD.\mathrm{gen}\) guard.  Write \(\psi_2\) for the dyadic
CRT component.  Its restriction to the subgroup of units congruent
to one modulo four is nontrivial; otherwise its conductor would
divide four, contradicting \(c_2\ge4\).  The image of that subgroup
has nontrivial power-of-two order.  Since the coefficient field has
characteristic zero, its unique element of order two is \(-1\).
Hence there is a dyadic unit \(v\equiv1\pmod4\) with
\(\psi_2(v)=-1\).  If \(\psi_2(-1)=1\), take \(u=-1\);
otherwise take \(u=-v\).  In both cases
\[
                  u\equiv-1\pmod4,\qquad \psi_2(u)=1.
 \tag{CAF.15}
\]
Total ramification of the dyadic cyclotomic extension supplies an
element \(s\) of inertia at two with this dyadic cyclotomic action.
All prime-to-two roots of unity are unramified there.  Thus the other
CRT components of \(\psi\) are one at \(s\), and the
\(\ell\)-adic cyclotomic character is one because \(\ell\) is odd.
Equation (CAF.15) gives
\(\omega(s)=\alpha(s)=1\) and \(\eta_4(s)=-1\), proving
\(\mathcal S\).  CAF6 gives
\(Q=(M^2\ell)_{\ell'}=N_{\ell'}^2\).
This verifies the seed-relevant part of \texttt{gen\_\allowbreak host} and
preserves the complete host small-prime scope.  It does not replace
the host by the \texttt{Dat} field or assume
\(\ell-1\nmid2e\).
% END CAF10

% BEGIN CAF11
\paragraph{CAF11. Installing the seed on the same marked frame.}
\label{claim:CAF11}
Choose \(Q\) by CAF6, CAF9, or CAF10 before selecting the frame.
Suppose the original \texttt{Has\allowbreak Frame\allowbreak W} at that \(Q\) has been
supplied.  Its very same
\(L,g_1,\Phi,C_0,w\), with the same \texttt{Gal\allowbreak Link}, then
satisfies the full original \texttt{Has\allowbreak Frame}, including its
\texttt{Seed\allowbreak Clause}.  In particular, an eventual proof of
\texttt{frame\_\allowbreak ex} can be consumed in both the \texttt{Gd\allowbreak Frame\allowbreak Ex}
and \texttt{Host\allowbreak Frame\allowbreak Ex} applications without a new seed choice or
a new marked element.

\emph{Proof.}
Unpack the supplied \texttt{Has\allowbreak Frame\allowbreak W}.  It contains a
\texttt{Linked} witness, hence a literal
\(k:\operatorname{GalLink}(\eta,\Psi,\sigma,g_1,B,M,L,\Phi)\),
the existing \(\Phi.\mathrm{Ok}(C_0)\) and
\(\Phi.\mathrm{Big}(C_0)\), a witness prime \(w\),
\(w\equiv1\pmod Q\), and \(U_L\subseteq H_\ell(w)\).
The \texttt{ab} and \texttt{gal} fields of \(k\) give abelian
\(L\) and \(G_L=U_L\cap\ker\eta\).

The \texttt{g0} field gives \(g_1\in U_L\) with
\(p(g_1)=\Phi.g_0\).  The \texttt{del} field and the fact
\(\Phi.g_0\in\ker\Phi.\mathrm{del}\) imply
\(\eta(g_1)=1\).  Hence \(g_1\in G_L\).
The \texttt{psi} field identifies
\(\alpha(g_1)=\Phi.\Psi(\Phi.g_0)\), whose square is not one
by \(\Phi.\mathrm{Ok}.\mathrm{baseSubgroup\_\allowbreak spec}\).
Thus the final witness in (CAF.10) is the already selected \(g_1\).
CAF7 provides the full \texttt{Seed\allowbreak Clause}.  Keep the same
\(L,g_1,\Phi,C_0,k\), its \texttt{Ok} and \texttt{Big} proofs,
and add that seed.  These are exactly all fields of
\texttt{Has\allowbreak Frame}; none has been reselected.

The source's \texttt{gd\_\allowbreak ex} and \texttt{host\_\allowbreak ex} use precisely
this installation after \texttt{frame\_\allowbreak ex}.  CAF9 and CAF10 give
their seed moduli directly, so their \texttt{Gen} constructions are
not additional seed hypotheses.  The existence of the large
\texttt{Has\allowbreak Frame\allowbreak W} is still required for either full frame-existence
conclusion.
% END CAF11

\subsection{The exact remaining interface}

The completed internal producers are \texttt{sc\_\allowbreak norm},
\texttt{sc\_\allowbreak teich}, \texttt{sc\_\allowbreak strip}, \texttt{sc\_\allowbreak quad},
\texttt{sc\_\allowbreak disc}, \texttt{hth}, \texttt{sc\_\allowbreak theta},
\texttt{auxw\_\allowbreak ex}, \texttt{sc\_\allowbreak cong}, \texttt{seed\_\allowbreak cl}, and
\texttt{field\_\allowbreak ex}, together with the two actual arithmetic
specializations of the seed hypotheses and the same-frame
installation.  The original \texttt{Gen} and \texttt{UFacts}
antecedents of the first two major source endpoints have been
preserved or weakened by proved implications; they have not been
silently assumed to exist.

For the remaining large frame, the source next chooses a real
quadratic subgroup \(H_0\), sets
\[
 U_m=H_0\cap H_\ell(w)\cap H_b\cap
          \ker\bigl(\epsilon\bmod\ell^{m+1}\bigr)^{\ell-1},
 \tag{CAF.16}
\]
proves its full \texttt{UFacts}, obtains one uniform class-rank
bound \(C_0\), and takes \(m=3C_0+7\).  Then
\(\ell^m>3C_0+7\).  CAF8 identifies the actual field belonging
to that very \(U_m\).  The still required pieces are the strict
cocycle constructions, the actual quotient/local linking frame,
and the cohomological \texttt{Ok} fields, with their exact
\texttt{GC.\allowbreak Hypotheses\_\allowbreak 2} input.  The elementary inequality alone
does not prove those fields.

The CountV lane must keep \(M=\ell N\), \(\eta=\eta_4\), and
\(H_b=\ker(\kappa_B)\), where \(\kappa_B\) is the away-from-\(\ell\)
part of the \(\ell\)-power-order character.  The \texttt{Dat}
lane uses \(M=\ell N|\operatorname{disc}K|\) and \(H_b=G\).
Moreover, its auxiliary \(\Phi\) is the cohomological frame over
\(G_{L^+}\).  It is not the original CountV host frame over
\(G_{\mathbb Q,S}\).  Transport between them must retain their
actual projection maps, coefficient modules, local conditions, and
normalized Selmer lengths.  No \texttt{HOST\_\allowbreak FRAME\_\allowbreak LENGTH},
\texttt{HOST\_\allowbreak VAN\_\allowbreak E}, or \texttt{HOST\_\allowbreak XPP} closure is claimed here.

Finally, \texttt{Linked\allowbreak End\allowbreak EAlt} consumes both rank-two and rank-three
automorphic seeds, the general \texttt{Mod\allowbreak UAlt} supplier, the same
frame and its common coefficient extension.  CAF7 supplies the
character triple those seeds require.  It does not assert that the
automorphic representations have been constructed.  Returning from
\texttt{Gd\allowbreak End\allowbreak Alt} to the full original \texttt{Gd\allowbreak End} likewise
retains its local and descent matching.  Thus the selected whole
obligations \texttt{AUT\_\allowbreak UNRAM\_\allowbreak COEF}, \texttt{THM\_\allowbreak S6\_\allowbreak S}, and
the CountV joins remain open, while their common character-seed
producer is now completely written.


\section{Constructing one marked arithmetic frame and its two strict classes}
\label{sec:marked-frame-construction}

\paragraph{The result and its precise boundary.}
This chapter constructs the actual auxiliary field, marked element,
profinite quotient, local groups and two nonzero strict cohomology
classes used by \texttt{frame\_\allowbreak ex}.  It proves the complete source
suppliers \texttt{exists\_\allowbreak twist\allowbreak Character\_\allowbreak continuous\allowbreak Cohomology},
\texttt{hlay}, \texttt{l15}, \texttt{ufacts}, \texttt{b0},
\texttt{link\_\allowbreak ex}, \texttt{ok\_\allowbreak basic}, \texttt{strict} and
\texttt{frobenius\_\allowbreak galois\allowbreak Field\_\allowbreak four\_\allowbreak compatibility}.
The construction uses one field and one marked element throughout.
The real quadratic auxiliary field admits a useful elementary
simplification: a Chinese-remainder construction replaces the
source's prime in a prescribed progression.

The entry premise is the \emph{full} original \texttt{Gen} record,
not just its seed-relevant subset from CAF.  The strict-class step
retains exactly the original \(P:\texttt{GC.\allowbreak Hypotheses}_{2}\).
Its finite local and global cohomology laws and Greenberg--Wiles
identity are the laws used below.  The uniform class-rank step
uses the already reconstructed TW1 and its named classical inputs.
No automorphic rank-two or rank-three seed is assumed or produced
by this construction.

At the end, the resulting frame has \texttt{Big}, both strict
classes, and twenty-one of the twenty-eight literal \texttt{Ok}
fields.  The seven remaining fields are named explicitly at the
end of the chapter.  Thus the chapter produces several complete
upstream theorems and an actual shared frame, but does not yet
claim the complete \texttt{Has\allowbreak Frame\allowbreak W}, \texttt{Gd\allowbreak Frame\allowbreak Ex},
\texttt{Host\allowbreak Frame\allowbreak Ex} or \texttt{Gd\allowbreak End} endpoints.

\subsection{Notation and the exact arithmetic input}

Let \(\Omega\) be the specified algebraic closure of \(\mathbb Q\)
and \(G=\operatorname{Gal}(\Omega/\mathbb Q)\).  Fix an odd prime
\(\ell\), a commutative local ring \(O\) with finite residue field
\(k\) of characteristic \(\ell\), and characters
\[
 \eta:G\longrightarrow\{\pm1\},\qquad
 \Psi:G\longrightarrow O^\times,\qquad
 \alpha=\overline\Psi:G\longrightarrow k^\times.
\]
Write \(\epsilon:G\to\mathbb Z_\ell^\times\) for the actual
cyclotomic character, \(\epsilon_r\) for its reduction modulo
\(\ell^r\), and \(\omega\) for its reduction in \(k^\times\).
The prime-field map is injective, so
\(\ker\omega=\ker\epsilon_1\).
For a valuation subring \(A\subset\Omega\), write \(D_A,I_A\)
for its decomposition and inertia groups, and write \(A\mid n\)
for \(A(n)<1\), with the source's multiplicative valuation.

We fix \(\sigma,B,M,H_b\) satisfying the original \texttt{Gen}:
\(M>0\), \(2\ell\mid M\), \(\sigma\) is a complex conjugation,
\(B\mid\ell\), and there are \(D>0,\delta\in\Omega\) with
\[
 \delta^2=-D,\qquad s\delta=\eta(s)\delta,
 \qquad p\mid D\Longrightarrow p\mid M
 \quad\text{for every rational prime }p.                  \tag{MFC.1}
\]
The characters \(\eta,\alpha\) factor through the action on
\(\mu_{M^a}\) for some \(a\); \(\eta\) is unramified at
\(\ell\); \(\alpha(\sigma)=1\); and one integer \(j\ge0\)
satisfies \(\alpha|_{I_A}=\omega^j|_{I_A}\) for every
\(A\mid\ell\).  The two separation hypotheses are
\[
 \ker\omega\cap\ker\alpha\not\subseteq\ker\eta,
 \qquad
 \exists t\in D_B\cap\ker\eta:\ \alpha(t)^2\ne1.          \tag{MFC.2}
\]
Finally \(H_b\) is open, contains \([G,G]\), has index
\(\ell^b\) for some \(b\ge0\), contains every \(I_A\) with
\(A\mid\ell\), and satisfies the exact lifting condition
\[
 A\nmid\ell,\quad s\in I_A\cap H_b,\quad\alpha(s)=1
                  \quad\Longrightarrow\quad\Psi(s)=1.      \tag{MFC.3}
\]
These are all nineteen source fields, with the finite residue-field
and odd-prime conditions included.  In particular (MFC.2)'s second
condition and (MFC.3) are not inferred from the weaker CAF entry
conditions.

The subgroup \(H_\theta=H_\theta(w)\) is the inverse image of
the \(\ell\)-th powers in \((\mathbb Z/w\mathbb Z)^\times\).
CAF4 proved the literal source \texttt{Hth} identification and its
open, abelian, index-\(\ell\), unramified and inertia properties.
A witness prime \(w\) means exactly
\[
 w\text{ prime},\quad w\nmid M,\quad w\equiv1\pmod\ell,
 \quad D_A\subseteq\ker\eta\cap\ker\alpha\cap\ker\omega
                     \quad(A\mid w).                       \tag{MFC.4}
\]
CAF5 supplies such a prime with \(w\equiv1\pmod Q\), for every
already prescribed positive \(Q\).

\subsection{The real quadratic subgroup and the entire tower}

% BEGIN MFC1
\paragraph{MFC1. An elementary real quadratic auxiliary subgroup.}
\label{claim:MFC1}
There is an open subgroup \(H_0\subset G\) of index two such that
every complex conjugation belongs to \(H_0\),
\[
 D_A\cap H_0=D_A\cap\ker\eta\quad(A\mid M),                \tag{MFC.5}
\]
and some valuation \(A\nmid M\) has \(I_A\not\subseteq H_0\).
These are all fields of \texttt{Frobenius\allowbreak Split\allowbreak Prime\allowbreak Datum}.

\begin{proof}
Choose an odd prime \(p\nmid2DM\).  Such a prime exists by
Euclid's argument.  The Chinese remainder theorem gives a positive
integer \(n\) with
\[
                  n\equiv-1\pmod{4M},\qquad
                  n\equiv p\pmod{p^2}.                    \tag{MFC.6}
\]
Choose \(x\in\Omega\) with \(x^2=-n\), put \(z=\delta x\),
and let \(H_0\) be the stabilizer of \(z\).  Since
\(z^2=Dn>0\) and \(v_p(Dn)=1\), \(Dn\) is not a rational
square.  Hence \(\mathbb Q(z)/\mathbb Q\) is quadratic and
\(H_0\) is open of index two.  Every complex embedding sends
\(z\) to one of the real roots of \(X^2-Dn\), so every complex
conjugation fixes \(z\).

Write \(-n=1+4Mt\) with \(t\in\mathbb Z\).  We claim that
every \(s\in D_A\), with \(A\mid M\), fixes \(x\).
It sends \(x\) to \(x\) or \(-x\).  In the latter case put
\(y=(1+x)/2\).  Then \(sy=1-y\) and
\(y(1-y)=-Mt\).  Decomposition-group invariance makes the
valuations of \(y\) and \(1-y\) equal, while their product
has valuation less than one because integers have valuation at
most one.  Both valuations are therefore less than one.
Their sum is one, contradicting the nonarchimedean inequality.
Thus \(sx=x\).  Since \(z\ne0\), (MFC.1) now gives
\(sz=z\) if and only if \(\eta(s)=1\), proving (MFC.5).

At the prime \(p\), the polynomial \(X^2-Dn\) is Eisenstein.
The quadratic extension is totally ramified there, so its inertia
group has a nonidentity image and sends \(z\) to \(-z\).
For any corresponding extension \(A\) of the \(p\)-adic
valuation, \(I_A\not\subseteq H_0\) and \(A\nmid M\).
This proves the final field.  No requirement that the auxiliary
integer \(n\) itself be prime was needed.
\end{proof}
% END MFC1

% BEGIN MFC2
\paragraph{MFC2. Exact cyclotomic layers and their local indices.}
\label{claim:MFC2}
For \(m\ge0\), define the literal source subgroup
\[
 H_m=\ker(\epsilon_{m+1}^{\ell-1}).                         \tag{MFC.7}
\]
It is open, contains \([G,G]\), has index \(\ell^m\), and
\[
 H_m\cap\ker\epsilon_1=\ker\epsilon_{m+1}.                \tag{MFC.8}
\]
It contains inertia at every valuation not above \(\ell\).
For every \(A\mid\ell\),
\[
 [I_A:I_A\cap H_m]=\ell^m,
 \qquad\epsilon_r(I_A)=(\mathbb Z/\ell^r\mathbb Z)^\times
                         \quad(r\ge0).                    \tag{MFC.9}
\]
These are the full conclusions of \texttt{hlay}.

\begin{proof}
For \(r\ge1\), \(\Phi_{\ell^r}(1+T)\) is Eisenstein at
\(\ell\): its reduction is a power of \(T\) and its constant
term is \(\ell\).  Thus \(\mathbb Q_\ell(\mu_{\ell^r})\)
is totally ramified of its full degree.  The inertia action on
these roots is consequently the whole unit group, at every
prolongation of the valuation.  The assertion at \(r=0\) is
trivial.  This also proves global surjectivity of \(\epsilon_r\).

For odd \(\ell\), the units modulo \(\ell^{m+1}\) have a
subgroup of order \(\ell-1\) consisting of the unique lifts
of the nonzero residues satisfying \(u^{\ell-1}=1\).
Existence and uniqueness follow from Hensel's lemma, since the
derivative \((\ell-1)u^{\ell-2}\) is a unit.  Its complement
is the subgroup \(1+\ell\mathbb Z/\ell^{m+1}\mathbb Z\),
of order \(\ell^m\).  Raising to \(\ell-1\) kills the first
factor and is an automorphism of the second.  Hence its image
has order \(\ell^m\), proving the index assertion and (MFC.9).
If a unit is also one modulo \(\ell\), its order is an
\(\ell\)-power.  If its \((\ell-1)\)-st power is one, its
order is one.  This proves (MFC.8).

The kernel of the finite cyclotomic action is open, so \(H_m\)
is open.  Its quotient is abelian.  Away from \(\ell\), all
\(\ell\)-power roots are unramified: their defining polynomials
have invertible derivative over the residue field, and Hensel
lifts their roots from an unramified extension.  Hence such
inertia is contained even in \(\ker\epsilon\), as asserted.
\end{proof}
% END MFC2

% BEGIN MFC3
\paragraph{MFC3. Residual-character identities cannot appear only in the tower.}
\label{claim:MFC3}
Fix \(H_0\) as in MFC1 and a witness prime \(w\).  Put
\[
 U_m=H_0\cap H_\theta(w)\cap H_b\cap H_m.                 \tag{MFC.10}
\]
For every \(i,j\in\mathbb Z\), the character
\(\chi=\omega^i\alpha^j\) has the following properties:
an identity \(\chi=1\) on \(U_m\) holds on \(G\);
an identity \(\chi=\eta_k\) on \(U_m\) holds on \(G\);
and an identity \(\chi=1\) on \(U_m\cap\ker\eta\) holds
on \(\ker\eta\).  Here \(\eta_k\) is the same quadratic
character viewed in \(k^\times\).  This is the full
\texttt{l15} conclusion.

\begin{proof}
Let \([G:H_b]=\ell^b\).  Each of \(H_\theta,H_b,H_m\)
is normal, with abelian quotient of order respectively
\(\ell,\ell^b,\ell^m\).  Lagrange's theorem gives
\[
 s^{\ell^{1+b+m}}\in U_m\qquad(s\in H_0).                 \tag{MFC.11}
\]
An element of \(k^\times\) whose \(\ell^N\)-th power is one
is one: \((u-1)^{\ell^N}=u^{\ell^N}-1=0\) in a field.
Therefore a character trivial on \(U_m\) is trivial on
\(H_0\), and the same argument applies on intersections with
any subgroup.

Choose \(A\nmid M\) and \(t\in I_A\setminus H_0\) from
MFC1.  The cyclotomic conductor hypotheses imply
\(\eta(t)=\alpha(t)=1\): inertia away from \(M\) fixes
\(\mu_{M^a}\).  Since \(\ell\mid M\), also
\(\omega(t)=1\).  Thus \(\chi(t)=\eta_k(t)=1\).
Every element outside the index-two subgroup \(H_0\) is the
product of an element in \(H_0\) with \(t\), or its inverse.
This extends the first identity to all of \(G\).  Apply the
same argument to \(\chi\eta_k^{-1}\) for the second identity.
For the third, note that \(t\in\ker\eta\) and apply the
identical coset argument inside \(\ker\eta\).
\end{proof}
% END MFC3

% BEGIN MFC4
\paragraph{MFC4. All twelve fields of the original \texttt{UFacts}.}
\label{claim:MFC4}
For every \(m\ge0\), \(U=U_m\) satisfies the following
complete source record:
\begin{enumerate}
\item \(U\) is open and contains \([G,G]\).
\item Every complex conjugation lies in \(U\).
\item \(I_A\cap U\subseteq\ker\eta\) for every valuation \(A\).
\item \(D_A\cap U\subseteq\ker\eta\) whenever \(A\mid M\).
\item If \(A\nmid\ell\), \(s\in I_A\cap U\), and
      \(\alpha(s)=1\), then \(\Psi(s)=1\).
\item If \(A\mid\ell\), some \(s\in I_A\cap U\) has
      \(\epsilon_1(s)\ne1\).
\item If \(A\mid\ell\), \([D_A:D_A\cap U]\ge\ell^m\).
\item For \(j=1,-1\), some \(s\in U\cap\ker\eta\) has
      \(\omega(s)\alpha(s)^j\ne1\).
\item The same assertion holds for \(j=2,-2\).
\item At every \(A\mid\ell\), some \(s\in D_A\cap U\)
      has \(\alpha(s)^2\ne1\).
\item There is no \(j\in\mathbb Z\) with
      \(\omega(s)=\alpha(s)^j\) for every
      \(s\in U\cap\ker\eta\).
\end{enumerate}
The first item contains the two separate fields
\texttt{is\allowbreak Open} and \texttt{ab}; the remaining items are
\texttt{real}, \texttt{unr}, \texttt{split}, \texttt{teich},
\texttt{mu}, \texttt{deg}, \texttt{f2b}, \texttt{f2c},
\texttt{f3} and \texttt{f5}.

\begin{proof}
Intersections of finitely many open subgroups are open.
Every index-two subgroup contains the commutator, as do the other
three factors in (MFC.10).  A complex conjugation has order two.
Its image in a group of odd order is trivial.  Hence it lies in
the three \(\ell\)-power-index factors and, by MFC1, in \(H_0\).

At a valuation above \(M\), (MFC.5) gives both the inertia and
decomposition assertions.  At every other valuation, \(\eta\)
is already trivial on inertia by its cyclotomic conductor.
The lifting assertion follows from \(U\subseteq H_b\) and
(MFC.3).

For \(A\mid\ell\), choose \(s\in I_A\) with
\(\epsilon_{m+1}(s)=-1\), using MFC2.  It belongs to \(H_m\)
because \(\ell-1\) is even.  Since \(\eta\) is unramified
at \(\ell\), (MFC.5) puts it in \(H_0\); here \(A\mid M\)
because \(\ell\mid M\).  It lies in \(H_b\) by hypothesis
and in \(H_\theta\) because \(w\ne\ell\).  Thus \(s\in U\)
and \(\epsilon_1(s)=-1\ne1\).  The index estimate follows
from \(I_A\subseteq D_A\), \(U\subseteq H_m\), and
\([I_A:I_A\cap H_m]=\ell^m\).

By (MFC.1), \(\eta(\sigma)=-1\), since a complex conjugation
negates a nonzero square root of a negative rational number.
Choose \(t\in\ker\omega\cap\ker\alpha\) with
\(\eta(t)=-1\) from (MFC.2).  Then
\[
 s_*=t\sigma\in\ker\eta,\qquad
                  \omega(s_*)=-1,\quad\alpha(s_*)=1.       \tag{MFC.12}
\]
If \(\omega\alpha^j=1\) throughout \(U\cap\ker\eta\),
MFC3 extends this identity to \(\ker\eta\), contradicting
(MFC.12).  In fact this proves the stated \texttt{f2b} and
\texttt{f2c} assertions for every integer \(j\).  Applying it
to \(\omega\alpha^{-j}\) also proves \texttt{f5}.

Finally take the element in the second assertion of (MFC.2).
Conjugate it so that its decomposition group is \(D_A\) for
the requested \(A\mid\ell\).  Characters into abelian groups
are invariant under conjugation, so it still has \(\eta=1\)
and \(\alpha^2\ne1\).  Equation (MFC.5) puts it in \(H_0\).
Raise it to \(\ell^{1+b+m}\).  Equation (MFC.11) puts the
result in \(U\), without leaving \(D_A\).  Its residual
square is still nonidentity by the characteristic-\(\ell\)
argument following (MFC.11).  This proves \texttt{f3} at every
prolongation, completing the record.
\end{proof}
% END MFC4

% BEGIN MFC5
\paragraph{MFC5. One class-rank bound before the layer is chosen.}
\label{claim:MFC5}
For the fixed \(H_0,w,H_b\), there is \(C_0\in\mathbb N\)
such that, for every \(m\ge0\),
\[
 \operatorname{clRk}_\ell
  (U_m\cap\ker\eta\cap\ker\alpha\cap\ker\epsilon_1)
                         \le C_0.                         \tag{MFC.13}
\]
Here \(\operatorname{clRk}\) is exactly \texttt{GC.\allowbreak cl\allowbreak Rk}:
the \(\mathbb F_\ell\)-rank of the locally constant additive
characters of the indicated subgroup which vanish on every
intersection with inertia.  This is the full \texttt{b0}
supplier, with \(C_0\) chosen before \(m\).

\begin{proof}
Set
\[
 V=H_0\cap H_\theta\cap H_b\cap\ker\eta\cap\ker\alpha
                                      \cap\ker\epsilon_1.
\]
This subgroup is open: the two character kernels contain the
kernel of a finite cyclotomic action.  It contains \([G,G]\).
Thus \(E=\Omega^V\) is a finite abelian number field, and it
contains \(\mu_\ell\).  By (MFC.8),
\[
 U_m\cap\ker\eta\cap\ker\alpha\cap\ker\epsilon_1
                 =V\cap\ker\epsilon_{m+1}
                 =G_{E(\mu_{\ell^{m+1}})}.                \tag{MFC.14}
\]
These are equal subgroups of the original \(G\), not merely
abstractly isomorphic Galois groups.

The final decomposition condition in the source's
\(\texttt{Kum.\allowbreak XDual}\ 1\) is vacuous because no valuation
has \(A(1)<1\).  Its inertia intersections are the inertia
groups for the actual fixed field in (MFC.14).  Finite class-field
reciprocity therefore identifies its additive character space with
\[
 \operatorname{Hom}(\operatorname{Cl}(E(\mu_{\ell^{m+1}})),
                    \mathbb F_\ell).
\]
There is no real-place qualification: \(E\) contains
\(\mu_\ell\) for odd \(\ell\), hence is totally complex.
This Hom space has dimension
\(\log_\ell|\operatorname{Cl}(E(\mu_{\ell^{m+1}}))/\ell|\),
by elementary finite abelian group theory.  TW1 gives one bound
\(\ell^{C_0}\) for these class-group quotients at all levels,
including repetitions at the beginning of the tower.  This proves
(MFC.13) at its exact fixed-field and coefficient scope.
\end{proof}
% END MFC5

% BEGIN MFC6
\paragraph{MFC6. The field and the marked local element are chosen together.}
\label{claim:MFC6}
After MFC5 choose \(m=3C_0+7\), put \(U=U_m\), and take
the actual field
\[
                  L=\Omega^{U\cap\ker\eta},\qquad F=L^+.
                                                               \tag{MFC.15}
\]
Then \(L/\mathbb Q\) is finite abelian CM,
\(G_F=U\), \(G_L=U\cap\ker\eta\), and there is
\[
              g_1\in D_B\cap G_L,\qquad\alpha(g_1)^2\ne1.
                                                               \tag{MFC.16}
\]
Moreover \(\ell^m>3C_0+7\).

\begin{proof}
MFC4 supplies precisely the \texttt{UFacts} input to CAF8.
That theorem gives (MFC.15), including its literal Galois-group
equalities and maximal real subfield.  Use MFC4's \texttt{f3}
at the original \(B\) to choose \(g_1\in D_B\cap U\) with
the required residual square.  Since \(B\mid M\), the
\texttt{split} field gives \(\eta(g_1)=1\), proving (MFC.16).
Finally \(m\ge7\), and \(\ell^m>m\) follows by induction
from \(\ell\ge3\).  No constant has been chosen after its
own layer: the order is \(H_0,w,C_0,m,L,g_1\).
\end{proof}
% END MFC6

\subsection{The literal global quotient and local frame}

% BEGIN MFC7
\paragraph{MFC7. All local places, embeddings and degrees on the same field.}
\label{claim:MFC7}
Let \(L,F,U,g_1\) satisfy MFC6; the same assertions below hold
for any original \(L\) and \(g_1\) satisfying the hypotheses
of \texttt{link\_\allowbreak ex}.  Let \(\mathcal V\) be the set of primes
of \(F\) above \(\ell\).  One may choose prolongations
\(A_v\) and elements \(\gamma_v\in\ker\eta\) with
\[
 A_v=\gamma_v B,\qquad
 G_v=D_{A_v}\cap U,\qquad I_v=I_{A_v}\cap U.              \tag{MFC.17}
\]
The set \(\mathcal V\) is finite and nonempty.  Every valuation
above \(\ell\) is \(U\)-conjugate to exactly one \(A_v\),
and, with the natural embeddings \(j_v:G_v\hookrightarrow U\),
\[
 d_v=[F_v:\mathbb Q_\ell]=[D_{A_v}:D_{A_v}\cap U],
 \qquad \sum_vd_v=[F:\mathbb Q].                           \tag{MFC.18}
\]
There are also complex conjugations \(c_\nu\in G_F\) at
every real place \(\nu\) of \(F\).

\begin{proof}
Primes above a fixed nonzero rational prime are finite and nonempty
in a number field.  Choose an extension \(A\) of each.  All
prolongations of the rational \(\ell\)-adic valuation to the
normal algebraic closure are conjugate, so \(A=\gamma B\)
for some \(\gamma\in G\).  If \(\eta(\gamma)=-1\), replace
\(\gamma\) by \(\sigma\gamma\) and \(A\) by \(\sigma A\).
Here \(\sigma\in U=G_F\), so the prime of \(F\) is unchanged,
and \(\eta(\sigma\gamma)=1\).  This gives the required choices.

Prolongations of one valuation of \(F\) are conjugate by \(G_F\).
Conversely a member of \(G_F\) cannot change the restricted
valuation on \(F\).  This proves both existence and uniqueness
of the index \(v\) in the asserted orbit description.  Restriction
identifies the decomposition and inertia groups over \(F\) with
the intersections in (MFC.17).  Completion of the finite extension
\(F/\mathbb Q\) identifies its local Galois degree with the
decomposition quotient in (MFC.18).  Summing the local degrees
gives the global degree, for example by the decomposition of
\(F\otimes_\mathbb Q\mathbb Q_\ell\) into its finite local
factors.

For a real embedding \(F\hookrightarrow\mathbb R\), extend
it to \(\Omega\hookrightarrow\mathbb C\).  Conjugation on
the image preserves the algebraic closure and fixes \(F\),
so it pulls back to the required \(c_\nu\in G_F\).  This
also records the exact real place at which it is a conjugation.
\end{proof}
% END MFC7

% BEGIN MFC8
\paragraph{MFC8. The complete original \texttt{Linked} frame.}
\label{claim:MFC8}
Under the full hypotheses of \texttt{link\_\allowbreak ex}, there is an
actual \(\Phi:\texttt{Frame}\ \ell\ k\) satisfying
\(\texttt{Linked}\ \ell\ \eta\ \Psi\ \sigma\ g_1\ B\ M\ L\ \Phi\).
In particular it contains every \texttt{Gal\allowbreak Link} and
\texttt{GC.\allowbreak FLink} field, with the same marked \(g_1\), field,
characters and local maps.

\begin{proof}
Identify \(G_F\) with \(U=\operatorname{upl}(\Omega,L)\)
by restriction to the same \(\Omega\); this is a topological
group isomorphism.  Its inverse merely regards an automorphism
fixing the embedded \(F\) as an \(F\)-automorphism.  Define
\[
 K_0=U\cap\ker\eta\cap\ker\alpha,
 \qquad
 J=\overline{\langle\!\langle
       I_A\cap K_0:\ A\nmid\ell
                   \rangle\!\rangle}_{,U},\qquad
                         \Gamma=U/J.                      \tag{MFC.19}
\]
The normal closure and closure are taken in this particular \(U\).
In the \(G_F\) notation, this is exactly
\(\texttt{tame\allowbreak N}(S,K_0)\), with \(S\) the full finite set
of primes above \(\ell\).  Every nontrivial valuation restricts
to a finite prime of \(F\); the trivial valuation has trivial
inertia and adds only the identity to the generating set.
Thus the two descriptions have equal normal closures, including
the valuation quantifier in \texttt{Gal\allowbreak Link.\allowbreak p\_\allowbreak ker}.

The subgroup \(K_0\) is open and normal in \(U\).  Its index there divides
\(2|k^\times|\), which is prime to \(\ell\), since its
quotient is the image of \((\eta,\alpha)\).  Because \(K_0\)
is closed and normal, \(J\subseteq K_0\).  Also
\(J\subseteq\ker\epsilon\), since \(\epsilon\) is
unramified away from \(\ell\).  Consequently \(\eta,\alpha\)
and \(\epsilon\) descend to characters
\[
 \delta_\Phi:\Gamma\to\{\pm1\},\quad
 \psi_\Phi:\Gamma\to k^\times,\quad
 \epsilon_\Phi:\Gamma\to\mathbb Z_\ell^\times.            \tag{MFC.20}
\]
Let \(p:U\twoheadrightarrow\Gamma\) be the canonical quotient.
It is continuous and onto.  The quotient is profinite: a closed
normal subgroup of a profinite group is the intersection of its
open normal overgroups, and the quotient is their inverse limit
of finite quotients.  This supplies its compactness, Hausdorff
property and total disconnectedness.

Use MFC7's local groups, with \(\operatorname{loc}_v=pj_v\),
the indicated \(I_v\), degrees \(d_v\), real-place index set,
and \(c_w=p(c_\nu)\).  Since \(\ell\mid M\), \texttt{UFacts.\allowbreak split}
says \(\eta=1\) on each \(G_v\).  Hence the source's
\(\operatorname{loc0}_v:G_v\to\ker\delta_\Phi\) is well
defined.  Set
\[
                    c=p(\sigma),\qquad g_0=p(g_1).
                                                               \tag{MFC.21}
\]
The latter is regarded as an element of \(\ker\delta_\Phi\),
since \(g_1\in G_L\).  Together (MFC.17)--(MFC.21) specify
every field of the original \texttt{Frame} record.

We verify the full linking records, to retain the actual maps.
The field is abelian and \(G_L=U\cap\ker\eta\) by hypothesis.
The \texttt{unr}, \texttt{split} and \texttt{teich} fields are
the corresponding \texttt{UFacts} fields.  The map \(p\), its
continuity, surjectivity and exact kernel are (MFC.19).  Its three
character identities are precisely
\[
 \delta_\Phi(pu)=\eta(u),\quad
 \epsilon_\Phi(pu)=\epsilon(u),\quad
 \psi_\Phi(pu)=\alpha(u).                                 \tag{MFC.22}
\]
The \texttt{c} and \texttt{g0} equations are (MFC.21).
The chosen \(A_v,\gamma_v\) give \texttt{pl}, \texttt{pl\_\allowbreak conj}
and \texttt{pl\_\allowbreak all}.  Inclusion gives continuous injective
\(j_v\), its exact decomposition-group range, the equality
\(\operatorname{loc}_v=pj_v\), and the exact inertia preimage.
Equation (MFC.18) gives \texttt{dv}.  Finally
\texttt{UFacts.\allowbreak mu} produces an element in \(I_{A_v}\cap U\)
whose residual cyclotomic character is nontrivial.  Viewed in
\(G_v\), it gives \texttt{Gal\allowbreak Link.\allowbreak mu}.

For \texttt{FLink}, take \(S=\mathcal V\), the original
\(G_F\)-version of \(K_0\), the canonical quotient map
\(\pi:G_F\to\Gamma\), and the decomposition inclusions
\(\iota_v:G_v\hookrightarrow G_F\).  Take the prime map
\(v\mapsto v\), the identity real-place equivalence, and the
actual \(c_\nu\) from MFC7.  The \texttt{GLink} fields are:
odd \(\ell\), all primes above \(\ell\), normal open
\(K_0\) of prime-to-\(\ell\) index, continuous surjective
\(\pi\), and kernel exactly \texttt{tame\allowbreak N}.  All have just
been verified.  The \texttt{LLink} fields are MFC7's place,
range and inclusion statements.  Its remaining equations are
the definitions of \(I_v,d_v,\epsilon_\Phi,c_w\).
The quadratic kernel is open because its inverse image under
the quotient map is open; \(\delta_\Phi(c)=-1\) follows
from \(\eta(\sigma)=-1\).  Finally both \(\delta_\Phi\)
and \(\psi_\Phi\) kill \(\pi(K_0)\).
The two final comparison equations in \texttt{Linked} say
that \(\pi\) and \(p\), and \(\iota_v\) and \(j_v\),
are these same maps under the fixed topological identification
\(G_F=U\).  They hold by construction.  This proves the full
source endpoint without an additional linking assumption.
\end{proof}
% END MFC8

% BEGIN MFC9
\paragraph{MFC9. All original basic frame conclusions and the size bound.}
\label{claim:MFC9}
For any linked frame in MFC8 and any marked
\(g_1\) with \(\alpha(g_1)^2\ne1\), all conclusions of
\texttt{ok\_\allowbreak basic} hold.  They include the nineteen
\texttt{Ok} fields other than \texttt{f4}, \texttt{f4'}
and the seven generic cohomological fields listed below.
They also give \(d_v\ge\ell^m\) for every \(v\) and
\(\mathcal V\ne\varnothing\).  With MFC6's \(m\),
\[
       3C_0+7<d=\sum_vd_v,\qquad5<d_v\quad\text{for all }v,
                                                               \tag{MFC.23}
\]
which is exactly \texttt{Big}\(\ C_0\).

\begin{proof}
For an arbitrary linked frame, each equality used here is one
of its \texttt{Gal\allowbreak Link} or \texttt{FLink} fields; the argument
does not require that its underlying type be the particular
quotient model constructed in MFC8.
The characters \(\delta_\Phi,\psi_\Phi\) have open kernel
because their pullbacks in (MFC.22) have open kernel and \(p\)
is a quotient map.  The same quotient property gives continuity
of \(\epsilon_\Phi\); the local maps are continuous by their
definitions.  Complex conjugation has square one, giving
\(c^2=1\), \(c_w^2=1\).  Its action on \(\delta\) gives
\(\delta_\Phi(c)=\delta_\Phi(c_w)=-1\), and the original
evenness gives \(\psi_\Phi(c)=1\).  The marked equation is
\(\psi_\Phi(g_0)^2=\alpha(g_1)^2\ne1\).
Oddness of \(\ell\) and finiteness of \(k\) are the input
\texttt{Gen} fields.  Since \(F\) is totally real, MFC7 gives
\(\#\operatorname{Inf}=[F:\mathbb Q]=\sum_vd_v\).

Reduction of (MFC.22)'s cyclotomic identity gives
\(\operatorname{om}_\Phi(pu)=\omega(u)\).  This is the
source's exact residual character: the two maps
\(\mathbb Z_\ell\to k\) reduce through the same prime field.
An element of \(I_B\cap U\) supplied by \texttt{UFacts.\allowbreak mu}
also lies in \(\ker\eta\), so its image in
\(G_0=\ker\delta_\Phi\) proves \texttt{f2a}.
Push the \texttt{UFacts.\allowbreak f2b} witnesses to \(G_0\) to prove
\texttt{f2b}.  Push the \texttt{f2c} witnesses to \(\Gamma\);
their \(\delta_\Phi\)-value is one, so the additional
\(\delta_{\Phi,k}\) factor in the source is also one.
Every local witness in \texttt{UFacts.\allowbreak f3} lifts uniquely into
the indicated local decomposition group, proving \texttt{f3}.
An identity \(\operatorname{om}_\Phi=\psi_\Phi^j\) on
\(G_0\) would pull back to the identity forbidden by
\texttt{UFacts.\allowbreak f5}; this proves \texttt{f5}.

The source's
\(N_v=I_v\cap\ker(\epsilon_\Phi\circ\operatorname{loc}_v)\)
has \(\omega=1\).  On inertia above \(\ell\), the original
relation \(\alpha=\omega^j\) then makes \(\alpha=1\).
Equation (MFC.22) proves exactly \texttt{Psi\_\allowbreak N}.
These are all the asserted basic fields, including the actual
marked element and local residual character.

By MFC7 and \texttt{UFacts.\allowbreak deg}, \(d_v\ge\ell^m\).
There is at least one prime above \(\ell\), hence
\(d\ge\ell^m\).  MFC6's strict inequality now gives both
inequalities (MFC.23), since \(3C_0+7\ge7>5\).
\end{proof}
% END MFC9

\subsection{Two strict classes on this same quotient}

% BEGIN MFC10
\paragraph{MFC10. A nonzero residual Selmer class with one allowed prime.}
\label{claim:MFC10}
Assume the finite local and global laws and the Greenberg--Wiles
identity contained in \(P:\texttt{GC.\allowbreak Hypotheses}_{2}\).
Let \(\kappa:G\to k^\times\) have open kernel,
\(\kappa(\sigma)\ne1\), and \(\kappa\ne\omega\).
Let \(S\) be a finite set of rational primes containing
\(\ell\) and every ramified prime of \(\kappa\).
Suppose \(w\in S\), \(w\ne\ell\),
\(\kappa|_{D_A}=\omega|_{D_A}=1\) for \(A\mid w\),
and \(\kappa|_{D_A}\ne1\) at every other place in \(S\).
Then there is a locally constant cocycle
\[
 y(st)=y(s)+\kappa(s)y(t)                                  \tag{MFC.24}
\]
which is not a global coboundary, vanishes on inertia outside
\(S\), and is a coboundary on every decomposition group above
\(S\setminus\{w\}\).  This gives the full
\texttt{st\_\allowbreak sel} conclusion, including every prolongation.

\begin{proof}
Descend \(\kappa\) to \(G_{\mathbb Q,S}\): its kernel
contains each excluded inertia group and hence their closed
normal closure.  Regard \(k(\kappa)\) as a finite
\(\mathbb F_\ell\)-module.  Its global invariants vanish
because \(\kappa(\sigma)\ne1\).  Its cyclotomic twisted
dual has no global invariants either.  Indeed an invariant
\(\mathbb F_\ell\)-linear functional \(f:k\to\mathbb F_\ell\)
would satisfy \(f(\kappa(s)x)=\omega(s)f(x)\).
Choose \(s\) with \(\kappa(s)\ne\omega(s)\).  Multiplication
by their nonzero difference is bijective on \(k\), forcing
\(f=0\).  The real-place invariant group also vanishes.
At every \(v\in S\setminus\{w\}\), the assumed local
nontriviality gives zero local invariants by the same scalar
argument.

Impose local condition zero at those places and the full
\(H^1\) at \(w\).  Let \(\mathcal S\) and
\(\mathcal S^\perp\) be the resulting finite Selmer groups
and their dual conditions.  At \(w\), both the representation
and the residual cyclotomic character are trivial.  Local Tate
duality gives \(|H^2(\mathbb Q_w,k)|=|k|\), while
\(|H^0(\mathbb Q_w,k)|=|k|\).  The local Euler characteristic
away from \(\ell\) gives
\(|H^1(\mathbb Q_w,k)|=|k|^2\).
In the Greenberg--Wiles identity every global and real invariant
factor is one, and every local factor other than \(w\) is
one.  Its remaining equality is therefore
\[
                      |\mathcal S|=|k|\,|\mathcal S^\perp|.
                                                               \tag{MFC.25}
\]
The dual group contains zero, so its cardinality is at least
one, and \(|k|>1\).  Thus \(\mathcal S\) has a nonzero
class.  Choose a continuous representative and inflate it to
\(G\), obtaining \(y\).  Surjectivity of
\(G\to G_{\mathbb Q,S}\) ensures that this inflated cocycle
cannot be a global coboundary.  It is zero on excluded inertia,
and membership in the zero local conditions supplies the required
local primitives.

These primitives hold for every prolongation.  If
\(A'=\gamma A\) and \(y(s)=(\kappa(s)-1)b\) on \(D_A\),
the cocycle identity gives
\[
 y(\gamma s\gamma^{-1})
       =(\kappa(s)-1)(\kappa(\gamma)b-y(\gamma)).
\]
Thus the primitive at \(A'\) is
\(\kappa(\gamma)b-y(\gamma)\).  This verifies the exact
valuation scope, without a new independent local choice.
\end{proof}
% END MFC10

% BEGIN MFC11
\paragraph{MFC11. Restriction preserves nontriviality, and the auxiliary inertia is killed.}
\label{claim:MFC11}
Two elementary facts will be used.  First, if a subgroup
\(U\subseteq G\) contains \([G,G]\) and
\(\kappa|_U\ne1\), a \(\kappa\)-cocycle which is a
coboundary on \(U\) is already a coboundary on \(G\).
Second, at a rational prime \(p\ne\ell\), if an open
normal subgroup \(H\) has index \(\ell\) and
\(I_A\not\subseteq H\), then every locally constant additive
homomorphism \(y:I_A\to k\) vanishes on \(I_A\cap H\).

\begin{proof}
For the first assertion, subtract the supplied coboundary and
assume \(y|_U=0\).  The subgroup \(U\) is normal because it
contains the commutator.  Choose \(u\in U\) with
\(\kappa(u)\ne1\).  Comparing the two expressions
\(su=(sus^{-1})s\) in the cocycle identity gives
\(y(s)=\kappa(u)y(s)\).  The nonzero scalar
\(\kappa(u)-1\) is invertible, so \(y(s)=0\) for every
\(s\), as required.

For the second assertion, the quotient map
\(G\to G/H\simeq\mathbb F_\ell\) restricts to a nonzero
additive character \(\varphi\) of \(I_A\), with kernel
\(I_A\cap H\).  The maximal pro-\(\ell\) abelian quotient
of inertia at \(p\ne\ell\) is \(\mathbb Z_\ell\): the
wild inertia is pro-\(p\), and the tame inertia quotient has
one \(\mathbb Z_\ell\) factor.  Hence its continuous
\(\mathbb F_\ell\)-valued character space has dimension one.
Every coordinate of the additive map \(y\), in an
\(\mathbb F_\ell\)-basis of \(k\), is a multiple of
\(\varphi\).  It therefore vanishes on the same kernel.
All coordinates vanish there, giving the assertion for \(k\).
This is the source's \texttt{st\_\allowbreak tame} mechanism.  A trivial
valuation has trivial inertia and causes no exception.
\end{proof}
% END MFC11

% BEGIN MFC12
\paragraph{MFC12. Both complete \texttt{Strict\allowbreak Cl} suppliers on one \(U_m\).}
\label{claim:MFC12}
For each \(a\in\{1,2\}\), put
\(\kappa_a=\alpha^{-a}\eta_k\).  Under \(P\), there is
a locally constant \(\kappa_a\)-cocycle \(y_a:G\to k\)
such that
\begin{align}
 &y_a|_{U_m}\text{ is not a coboundary};\notag\\
 &\forall A\mid\ell\ \exists b_A\quad
       y_a(s)=(\kappa_a(s)-1)b_A\quad(s\in D_A);\notag\\
 &A\nmid\ell,\quad s\in I_A\cap U_m,\quad
       \eta(s)=\alpha(s)=1\quad\Longrightarrow\quad y_a(s)=0.
                                                               \tag{MFC.26}
\end{align}
This is the full original \texttt{strict} theorem for each
\(m\) and each of its two allowed characters.  The two choices
are made on the same \(H_0,w,U_m\).

\begin{proof}
The character \(\kappa_a\) factors through \(\mu_{M^b}\)
for some \(b\), so its kernel is open and it is unramified
away from \(M\).  It is odd:
\(\kappa_a(\sigma)=-1\), since \(\alpha(\sigma)=1\).
It differs from \(\omega\): otherwise, on
\(U_m\cap\ker\eta\), one would have
\(\omega\alpha^a=1\), contradicting MFC4's appropriate
\texttt{f2b} or \texttt{f2c} witness.

At every \(A\mid\ell\), MFC4 gives \(s\in D_A\cap U_m\)
with \(\alpha(s)^2\ne1\).  Its \(\eta\)-value is one,
so \(\kappa_a(s)=\alpha(s)^{-a}\ne1\) for both \(a=1,2\).
At \(w\), the entire decomposition group kills
\(\eta,\alpha,\omega\) by (MFC.4).  Apply MFC10 with
\(S=\{\ell,w\}\) together with the ramified primes of
\(\kappa_a\).  At a ramified prime, the character is
nontrivial on inertia, and therefore on decomposition.  All
the hypotheses of MFC10 hold.  It yields a global noncoboundary
\(y_a\) with the indicated local primitives above \(\ell\).

At any prime other than \(w\), if
\(s\in I_A\) and \(\kappa_a(s)=1\), the cocycle is zero:
outside \(S\) it vanishes on inertia, while in
\(S\setminus\{w\}\) it is a local coboundary.
At \(w\), \(\kappa_a\) is trivial on decomposition, so
\(y_a|_{I_A}\) is additive.  The subgroup \(H_\theta(w)\)
has index \(\ell\), and inertia at \(w\) maps onto its
quotient, by CAF4.  MFC11 makes \(y_a\) vanish on
\(I_A\cap H_\theta(w)\), hence on \(I_A\cap U_m\).
This proves the last clause of (MFC.26) also at \(w\),
because \(\eta=\alpha=1\) implies \(\kappa_a=1\).

Finally \(\kappa_a\) is nontrivial on \(U_m\), witnessed
by the same local element used above.  Since \(U_m\) contains
the commutator, MFC11 says that a coboundary restriction to
\(U_m\) would make \(y_a\) a global coboundary.  This is
impossible.  All assertions of \texttt{Strict\allowbreak Cl} are now verified.
Nothing in choosing \(y_1,y_2\) changes the already chosen
subgroup, field, local groups or marked element.
\end{proof}
% END MFC12

% BEGIN MFC13
\paragraph{MFC13. Exact descent of each strict class to the frame cone.}
\label{claim:MFC13}
Let \(\Phi\) be any actual linked frame as in MFC8, and let
\(y_a\) satisfy \texttt{Strict\allowbreak Cl} for \(U=\operatorname{upl}(\Omega,L)\).
Then
\[
 H^1_C\!\left(\Phi.\operatorname{loc},
    k(\psi_\Phi^{-a}\delta_{\Phi,k}),\,b_v=0\right)
                         \ne0\qquad(a=1,2).               \tag{MFC.27}
\]
These are exactly the two source \texttt{Ok} fields
\texttt{f4} and \texttt{f4'}; the quotient and all local
maps are the same ones used in MFC8.

\begin{proof}
Write \(\chi_a=\psi_\Phi^{-a}\delta_{\Phi,k}\).
Equation (MFC.22) gives \(\chi_a\circ p=\kappa_a|_U\).
The strict inertia clause makes \(y_a\) zero on every
generator of the exact kernel \(J\) in (MFC.19).
To justify passage to its closed normal closure, consider
\[
                    H=\{u\in U:p(u)=1,\ y_a(u)=0\}.
\]
The cocycle identity makes \(H\) a subgroup.  If \(h\in H\),
then \(\kappa_a(h)=1\), and the same identity gives
\(y_a(uhu^{-1})=\kappa_a(u)y_a(h)=0\); hence it is normal.
It is closed because \(p\) is continuous and \(y_a\) is
locally constant.  It contains the kernel generators, so it
contains \(J\).  Thus \(y_a\) is constant on fibers of
\(p\), and defines \(\varphi_a:\Gamma\to k\).
The quotient-map property of \(p\) shows that
\(\varphi_a\) is locally constant.  Its cocycle identity
follows by lifting two elements of \(\Gamma\) to \(U\).

For each \(v\), the chosen \(A_v\) lies above \(\ell\).
Choose its local primitive \(b_v\) from (MFC.26).  Under
the exact inclusion \(j_v\), this gives
\[
 \varphi_a(\operatorname{loc}_v(s))
      =(\chi_a(\operatorname{loc}_v(s))-1)b_v
                                  \quad(s\in G_v).        \tag{MFC.28}
\]
By the literal source definition, the pair
\((\varphi_a,(b_v)_v)\) is therefore a member of
\(Z^1_C\) with all \(b_v\)-subspaces equal to zero.
If its class in \(H^1_C\) were zero, the source's
\texttt{continuous\allowbreak Coboundaries1} condition would supply one
\(b\in k\) with
\(\varphi_a(g)=(\chi_a(g)-1)b\) for every \(g\).
Pulling back by \(p\) would make \(y_a|_U\) a coboundary,
contrary to MFC12.  Thus the class is nonzero, proving (MFC.27).
The argument applies to both characters without changing any
frame data.
\end{proof}
% END MFC13

% BEGIN MFC14
\paragraph{MFC14. The actual shared marked construction, with an exact remaining boundary.}
\label{claim:MFC14}
Fix the full original \texttt{Gen}, the original \(P\), and
any prescribed \(Q>0\).  There exist one
\[
                  L,\quad g_1,\quad\Phi,\quad C_0,\quad w
\]
with the complete original \texttt{Linked} relation,
\texttt{Big}\(\ C_0\), the witness-prime conditions
(MFC.4), \(w\equiv1\pmod Q\),
\(\operatorname{upl}(\Omega,L)\subseteq H_\theta(w)\),
and all the twenty-one \texttt{Ok} fields proved in MFC9
and MFC13.  The bound (MFC.13) holds at the selected layer.
When \(O=c.O\) is the original coefficient DVR, choose
\(Q\) first by CAF7.  This same constructed \(L\) then
has the full original \texttt{Seed\allowbreak Clause}, with the same
marked square witness \(g_1\).

\begin{proof}
Choose \(H_0\) by MFC1 and choose \(w\) by CAF5 for the
already fixed \(Q\).  Choose \(C_0\) by MFC5, before
selecting \(m=3C_0+7\).  MFC4 and MFC6 give the actual
embedded CM field and marked \(g_1\).  MFC8 constructs
\(\Phi\) with every linking field.  MFC9 proves the basic
fields and \texttt{Big}; MFC12 and MFC13 prove the two
remaining strict-class fields on this very quotient.
The subgroup inclusion follows from the definition of \(U_m\),
and the class-rank bound is the selected instance of MFC5.

For the coefficient-DVR assertion, the full \texttt{Gen}
implies the seed input used by CAF7.  Make its choice of
\(Q\) before the preceding construction.  The resulting
\(L,w\) satisfy all its field, splitting and congruence
hypotheses.  Equation (MFC.16) supplies
\(g_1\in G_L\) with \(\alpha(g_1)^2\ne1\), so CAF7
returns the original \texttt{Seed\allowbreak Clause}.  Its character
triple is installed on this field; it does not replace the
field, marked element, quotient, local maps or either strict class.
This completes the claimed construction at its stated boundary.
\end{proof}
% END MFC14

\subsection{The seven remaining \texttt{Ok} fields, and the original exits}

The seven fields still needed on the \emph{same} \(\Phi,C_0\)
are \texttt{phi}, \texttt{gen}, \texttt{cl}, \texttt{ec},
\texttt{b1p}, \texttt{b1} and \texttt{h1}.
Their source scopes remain unchanged: \texttt{phi} is the
global \texttt{Phi\allowbreak Cond}; \texttt{gen} chooses \(d_v\) local
inertia generators with the prescribed exceptional index and
character restrictions, uniformly for every finite abelian
\(\ell\)-group character; \texttt{cl} gives one exponent
annihilating every everywhere-unramified finite abelian
\(\ell\)-group character of \(\Phi.G0\); and \texttt{ec}
is the full cone Euler inequality for every finite coefficient
field, finite representation and stable collection of local
subspaces satisfying its stated dual-invariant condition.

The last three retain their exact coefficient modules and local
conditions:
\begin{align}
 \operatorname{rank}_k H^1_C(\operatorname{loc},
       k(\psi_\Phi^{-2}\delta_{\Phi,k}),0)&\le C_0,
          &&\texttt{b1p},\notag\\
 \operatorname{rank}_k H^1_C(\operatorname{loc2},
       k(\psi_\Phi^{-1}|_{\Phi.G0}),0)&\le C_0,
          &&\texttt{b1},\notag\\
 \operatorname{rank}_k H^1_C(\operatorname{loc},k,\top)
                                                &\le1+C_0,
          &&\texttt{h1}.                                  \tag{MFC.29}
\end{align}
The bare uniform unramified character bound in MFC5 is not
silently identified with these twisted cone bounds.  Their
restriction, coefficient and local-condition comparisons still
have to be supplied.  Filling these seven fields in the same
constructed record gives the literal \texttt{Has\allowbreak Frame\allowbreak W};
MFC14's seed installation then gives the literal
\texttt{Has\allowbreak Frame}.

To obtain the original \texttt{Gd\allowbreak Frame\allowbreak Ex} and
\texttt{Host\allowbreak Frame\allowbreak Ex}, their full entry suppliers
\texttt{gen\_\allowbreak gd} and \texttt{gen\_\allowbreak host} must also be connected
at their original scopes.  CAF9--CAF10 supply their seed subset,
not every \texttt{Gen} field.  In particular the host's
\(H_b=\ker\kappa_B\), its index, local lifting property
and nonquadratic decomposition witness must be retained; a
\(\top\) subgroup or the Dat numerical restriction cannot
be substituted.  The rank-two/rank-three automorphic seeds,
common coefficient extension and original/alternative endpoint
matching remain further work beyond these arithmetic frames.

\paragraph{Pinned correspondence.}
The immutable primary sources are the \texttt{Frobenius\allowbreak Smooth\allowbreak Vectors}
files named in this chapter, the literal frame and linking definitions
in \texttt{Continuous\allowbreak Coboundaries2}, \texttt{Host\allowbreak Frame\allowbreak Ex} and
\texttt{GLink\allowbreak Quot}, and the exact class-character definition
\texttt{Kum.\allowbreak XDual}.  Complete file snapshots, statement ranges,
source hashes and claim-block hashes accompany this chapter.
The elementary CRT choice in MFC1 is an intentional human
simplification with the same complete output record.  No changed
field model, local identification or order of choices is hidden
in that simplification.


\section{The Bessel nonvanishing base, its parameter walk, and its literal return}
\label{sec:bessel-nonvanishing-v50}

This section isolates the original nonvanishing problem, proves its analytic
base case and its elementary parameter walk, and returns the result to the
endomorphism-valued integral actually used by the local-to-global argument.
The remaining input is identified exactly: the original two-sided
\texttt{Step\allowbreak Statement}, which propagates nonvanishing by one unit in a
Harish--Chandra parameter.  We do not replace that input by an abstract
branching assertion, and do not claim to have proved it here.  In particular,
the final implication to \texttt{Bessel\allowbreak NVGap} and the supplied-form version of
\texttt{Concrete\allowbreak Setting\allowbreak Exists} is conditional on that one producer.

The proofs use DC1--DC13 for the actual large-gap discrete series and its
coefficient kernel, BWP1--BWP3 for the actual unipotent group and compact
stabilizer, and AMD8 for passage from the source's weak \texttt{Is\allowbreak Haar\allowbreak On}
predicate to ordinary Haar measure.  They do not use a multiplicity-one
theorem, a formal-degree formula, or the desired Bessel nonvanishing itself.
Source locations and their exact pinned bytes are recorded in the accompanying
manifest.  References to a source declaration specify the target of a human
proof; the declaration is not being used in place of that proof.

\subsection{The original integral and the analytic notation}

Put $n=2r+k+1$ and use the ordered Witt decomposition
$E\oplus V\oplus F$ of BWP1, with $\dim E=\dim F=r$ and $V$ of signature
$(k,1)$.  The dual $F$ coordinates are reversed when necessary, exactly as in
the source matrix $H=H_{pk}$.  Write
\[
 G=U(H),\qquad N=N_0=\texttt{Nset}(k,r),\qquad
 H_z=\texttt{UWset}(k,r,z)\simeq U(k).
\]
Here $z\in V$ satisfies $T_z=-z^*Hz>0$.  The group $H_z$ fixes $E,F,z$
pointwise; $N$ acts trivially on the successive pieces of the prescribed
isotropic flag.  If $\ell_a$ are the original simple-root coordinates, their
last member is the $E_{r-1}$ coordinate of $nz$ and the others are the
$(a,a+1)$ entries.  For $c_a\ne0$ the original character is
\[
 \psi_{z,c}(n)=
 \exp\left(i\operatorname{Re}\sum_{a=0}^{r-1}c_a\ell_a(n)\right).
 \tag{BNV.1}
\]
For $r=0$ the sum is empty, $N=\{1\}$, and this character is $1$.

The signature $\epsilon=\epsilon^{48}$ has $k$ initial positive coordinates,
then alternating negative and positive coordinates.  Thus
$p=k+r$, $q=r+1$, and, for $k\ge1$, $q_0=\min(p,q)=r+1$.
The target parameter is
\[
 \lambda=(v_0,\ldots,v_{k-2},-b_0,-b_1,\ldots,-b_{2r+1}).
 \tag{BNV.2}
\]
A Harish--Chandra parameter of length $m$ is a strictly decreasing vector
whose coordinates belong to $(m-1)/2+\mathbb Z$.  The interlacing and
half-integrality assumptions are exactly
\[
 u_i>v_i>u_{i+1}\quad(0\le i<k-1),\qquad u_{k-1}>-b_0,
 \qquad u_i-\lambda_s\in\mathbb Z+\tfrac12.
 \tag{BNV.3}
\]

The original \texttt{Bessel\allowbreak NVGap} has the following scope.  For every even
$k\ge2$, every $r\ge0$, every $z,c,u,v,b$ as above with $u,\lambda$
Harish--Chandra parameters and with the original \texttt{Gap\allowbreak Big}, take
\emph{any} complete complex Hilbert space $\mathcal H$ and representation
$\sigma$ satisfying the literal predicate
$\texttt{Is\allowbreak DSHpk}(k,r,\epsilon,\lambda,\sigma)$.  Take any finite-dimensional
complex Hilbert space $V_\pi$ and representation $\pi$ satisfying the literal
$\texttt{Is\allowbreak HCType\allowbreak W}(k,r,z,u,\pi)$; take any lowest-type unit vector $w$ of
$\sigma$; and take any original \texttt{Is\allowbreak Haar\allowbreak On} measures $\mu_W,\mu_N$ on
$H_z,N$.  The conclusion is the existence of $x,x_0\in G$ such that
\[
 \mathcal B_{\sigma,w}(x_0,x):=
 \int_{H_z\times N}
 \langle\sigma(x_0)w,\sigma(hn)\sigma(x)w\rangle
 \overline{\psi_{z,c}(n)}\,\pi(h)^{-1}\,d\mu_W(h)d\mu_N(n)
 \ne0
 \quad\text{in }\operatorname{End}_{\mathbb C}(V_\pi).
 \tag{BNV.4}
\]
Inner products are conjugate linear in the first variable.  The measures in
the formal statement live on ambient matrices and are supported on the
displayed groups; writing them on the groups makes no change.  The inverse
in (BNV.4) is the operator inverse, equal on $H_z$ to $\pi(h^{-1})$.
No Haar normalization, chosen model of $\sigma$, chosen lowest vector, or
chosen nondegenerate character may be silently substituted for these
quantifiers.

For a compact lowest type $T$ of weight
$\Lambda=\lambda+\rho_n-\rho_c$, let $\Psi_{T,\mu}$ be the normalized
coefficient kernel of DC2--DC3, with $\mu=\lambda+\rho$.  A Cayley matrix
$P$ means $P^*J_\epsilon P=H$; Euclidean unitarity of $P$ is not required.
Define the exact vector-valued source integral
\[
 A(v_1,v_2,\xi;x_0,x)=
 \int_{H_z\times N}
 \langle v_2,\Psi_{T,\mu}(Px_0^{-1}hnxP^{-1})v_1\rangle
 \overline{\psi_{z,c}(n)}\,\pi(h)^{-1}\xi\,d\mu_Wd\mu_N.
 \tag{BNV.5}
\]
This is \texttt{aa\allowbreak Fun}, not an unspecified equivariant functional.
$\texttt{Bessel\allowbreak NZ}(\epsilon,\lambda,u)$ means that there exist all the data
in (BNV.5), including $T,P,z,c,\pi$, the two Haar measures, and
$v_1,v_2,\xi,x_0,x$, for which (BNV.5) is nonzero.
$\texttt{Bessel\allowbreak All}(\epsilon,\lambda,u)$ means that for \emph{every} such
choice of data and \emph{every} $0\ne v\in T$, some $\xi,x_0,x$ make
$A(v,v,\xi;x_0,x)\ne0$.  These are the original predicates.

\subsection{A direct proof of the original subgroup-integrability input}

% BEGIN BNV1
\paragraph{BNV1. The whole original \texttt{Int\allowbreak SStatement}.}
\label{claim:BNV1}
For every $k,r\ge0$, admissible $z$, original Haar pair, $x_0,x\in G$, and
real $b> b_S(k,r):=2rn^2$, the function
\[
 (h,n')\longmapsto \|x_0^{-1}hn'x\|_{\rm F}^{-b}
 \tag{BNV.6}
\]
is integrable on $H_z\times N$.  Here $\|\cdot\|_{\rm F}$ is Frobenius
norm, the square root of the source's \texttt{hs}.

\begin{proof}
First suppose $r\ge1$.  In the real polynomial chart of BWP1 the dimension
is $d=2r^2+2rk+r$.  Its measure is a positive multiple of Lebesgue measure,
and the inverse-coordinate estimate of BWP2 is
\[
 1+|y(n')|\le C\|n'\|_{\rm op}^2.
\]
It follows that $\|n'\|_{\rm op}^{-b}$ is integrable whenever $b>2d$.
Indeed it is bounded by a constant times $(1+|y|)^{-b/2}$, whose radial
integral converges for $b/2>d$.  Since
\[
 2d=2r(2r+2k+1)\le2r(2r+k+1)^2=2rn^2,
\]
the asserted exponent suffices.  The elementary inequality used here holds
also when $k=0$, because $r\ge1$ and expansion gives nonnegative terms.

For $g=x_0^{-1}hn'x$,
\[
 \|n'\|_{\rm op}
 =\|h^{-1}x_0gx^{-1}\|_{\rm op}
 \le C_{x_0,x,z}\|g\|_{\rm F},
\]
uniformly for $h\in H_z$, since $H_z$ is compact.  Its Haar mass is finite.
Comparison and Tonelli give (BNV.6).  If $r=0$, $N$ is a point and the
continuous, strictly positive norm in (BNV.6) has a positive minimum on the
compact $H_z$, so the conclusion holds for every $b>0=b_S$.
AMD8 and Haar uniqueness allow every original weak Haar pair, with positive
scalars absorbed in the estimates.  Thus this is the original universal
\texttt{Int\allowbreak SStatement}, including its exponent and degenerate rank case.
\end{proof}
% END BNV1

\subsection{The finite parameter walk}

For this subsection let $k\ge1$, let (BNV.2) be Harish--Chandra, and assume
(BNV.3).  Write $\mathrm{Gap}_h(\eta)$ for
$|\eta_a-\eta_b|\ge h$ whenever $\epsilon_a\ne\epsilon_b$.
Suppose $b_{j+1}-b_j\ge h$.  Define
\[
 \lambda^{(t)}=
 (u_0-\tfrac12,\ldots,u_{k-1}-\tfrac12,
       -b_1-t,-b_2-2t,\ldots,-b_{2r+1}-(2r+1)t),\quad t\in\mathbb N.
 \tag{BNV.7}
\]
A predicate $\mathcal P(\eta)$ is \emph{step closed} at $(\epsilon,u,h)$ if
it is preserved by each move $\eta\mapsto\eta\pm e_s$ for which both
endpoints are Harish--Chandra and satisfy $\mathrm{Gap}_h$, and
\[
 u_j\ne\eta_s\pm\tfrac12\quad\text{for every }j.
 \tag{BNV.8}
\]
This is exactly \texttt{Step\allowbreak Closed}; there is no assertion here that a
particular representation-theoretic predicate has this property.

% BEGIN BNV2
\paragraph{BNV2. The sign, lattice, and interlacing arithmetic.}
\label{claim:BNV2}
The signature has $(p,q)=(k+r,r+1)$.  Every coordinate of $u$ belongs to
$(k-1)/2+\mathbb Z$, and $u$ is strictly decreasing.  Each displacement
\[
 d_i=u_i-\tfrac12-\lambda_i\quad(0\le i<k)
 \tag{BNV.9}
\]
is a nonnegative integer, where $\lambda_{k-1}=-b_0$.
Every adjacent $b$ difference is a positive integer.

\begin{proof}
Counting the alternating tail after the initial $k$ positives gives the
signature.  Interlacing gives strict decrease of $u$, and its difference
from any $\lambda_s$ belongs to $\mathbb Z+1/2$.  Since
$\lambda_s\in(n-1)/2+\mathbb Z$ and $n=k+2r+1$, this puts
$u_i$ in $n/2+\mathbb Z$, the same coset as $(k-1)/2+\mathbb Z$:
their difference is $r+1$.  Positivity and half-integrality of
$u_i-\lambda_i$ give (BNV.9).  The $-b_j$ are strictly decreasing
coordinates of the same integral coset, proving the last assertion.
\end{proof}
% END BNV2

% BEGIN BNV3
\paragraph{BNV3. Adjacent sign changes suffice for all noncompact gaps.}
\label{claim:BNV3}
If $\eta$ is strictly decreasing and
$\eta_a-\eta_{a+1}\ge h$ at every adjacent sign change, then
$\mathrm{Gap}_h(\eta)$ holds.

\begin{proof}
For $a<b$ of opposite signs, some $j$ with $a\le j<b$ changes sign.
The telescoping sum
$\eta_a-\eta_b=\sum_{i=a}^{b-1}(\eta_i-\eta_{i+1})$
is at least its $j$th positive summand, hence at least $h$.  Reversing
$a,b$ gives the absolute-value formulation.  This proves the source's full
\texttt{Adj\allowbreak Gap\allowbreak Statement}, for every signature and every real $h$.
\end{proof}
% END BNV3

% BEGIN BNV4
\paragraph{BNV4. Raising the deformed tail without a forbidden crossing.}
\label{claim:BNV4}
Starting at $\lambda^{(t)}$, for $j=1,\ldots,2r+1$ in that order raise
the coordinate $-b_j-jt$ to $-b_j$ in $jt$ unit steps.  Every endpoint is
Harish--Chandra, satisfies $\mathrm{Gap}_h$, and satisfies the upward
exclusion (BNV.8).

\begin{proof}
All moves are integral, so the lattice condition is preserved.  The
preceding tail coordinate has already reached $-b_{j-1}$ when $j>1$;
the current coordinate is at most $-b_j$.  Their difference is at least
$b_j-b_{j-1}$.  The next tail coordinate is still at
$-b_{j+1}-(j+1)t$, so its difference below the current coordinate is at
least $b_{j+1}-b_j+t$.  For $j=1$ the preceding coordinate is
$u_{k-1}-1/2\ge-b_0$, giving at least $b_1-b_0$.  Thus the ordering
is strict and every adjacent sign-changing gap is at least its original
$b$ difference.  BNV3 applies.

Immediately before an upward step the moving coordinate $a$ is at most
$-b_j-1$.  In particular $a+1/2<-b_0<u_{k-1}\le u_i$, so no forbidden
equality occurs.  This proves every segment of the original
\texttt{Raise\allowbreak Step\allowbreak Statement}, including a segment of zero length.
\end{proof}
% END BNV4

% BEGIN BNV5
\paragraph{BNV5. Lowering the initial block without a forbidden crossing.}
\label{claim:BNV5}
After BNV4, lower the initial coordinates in order
$i=k-1,k-2,\ldots,0$, taking $u_i-1/2$ to $\lambda_i$ in $d_i$ unit
steps.  All endpoints satisfy the same Harish--Chandra and gap conditions
and the downward exclusion (BNV.8).

\begin{proof}
The coordinate on the left, if any, is still $u_{i-1}-1/2$, strictly
above the current value.  Coordinates on the right have reached their
strictly smaller target values.  The lattice is unchanged.  The only
adjacent sign change involving the initial block is between positions
$k-1$ and $k$; its smallest possible difference is
$-b_0-(-b_1)=b_1-b_0$.  All other sign-changing gaps are the fixed tail
differences.  BNV3 again applies.

Before a downward step the current value $a$ obeys
$\lambda_i+1\le a\le u_i-1/2$.  If $j\le i$, then
$a-1/2\le u_i-1<u_j$.  If $j>i$, interlacing gives
$\lambda_i>u_{i+1}\ge u_j$, so
$a-1/2\ge\lambda_i+1/2>u_j$.  Thus equality is impossible in both
cases.  There are no $j>i$ when $i=k-1$.  This proves the full original
\texttt{Lower\allowbreak Step\allowbreak Statement}.
\end{proof}
% END BNV5

% BEGIN BNV6
\paragraph{BNV6. The whole original \texttt{Walk\allowbreak Statement}.}
\label{claim:BNV6}
For every predicate $\mathcal P$ step closed at $(\epsilon,u,h)$ and
every $t\in\mathbb N$,
\[
 \mathcal P(\lambda^{(t)})\quad\Longrightarrow\quad\mathcal P(\lambda).
 \tag{BNV.10}
\]

\begin{proof}
Apply step closure successively along the $\sum_{j=1}^{2r+1}jt$ moves of
BNV4 and then the $\sum_{i=0}^{k-1}d_i$ moves of BNV5.  These are finite
nonnegative integers.  The initial and terminal parameters are precisely
(BNV.7) and (BNV.2).  Zero-length segments need no invocation of step
closure, so the proof includes $t=0$, $r=0$, and $k=1$.
\end{proof}
% END BNV6

\subsection{The exact large-gap and large-parameter bridges}

For any signature put
\[
 \begin{split}
 h_{\rm need}(\epsilon)=\max\{&2+\max(3n/2,1+npq),\\
 &16nq_0+8\sqrt{2q_0}+2,\\
 &8\sqrt{2q_0}\,[4\max(2nq_0,2(q-1)n^2)+1]\}.
 \end{split}
 \tag{BNV.11}
\]
$\texttt{Gap\allowbreak Lap}(\epsilon,\eta,h)$ means that $\eta$ is strictly
decreasing, $\mathrm{Gap}_h(\eta)$ holds, and $h>h_{\rm need}(\epsilon)$.
The original $\texttt{Step\allowbreak Gap}(\epsilon,r,h)$ consists of
\[
 \max(3n/2,1+npq)<h-2,\qquad
 4\max(2nq_0,2rn^2)<\frac{h}{8\sqrt{2q_0}}-1.
 \tag{BNV.12}
\]
In applications here $q_0=r+1>0$.

% BEGIN BNV7
\paragraph{BNV7. The original cubic gap implies both analytic thresholds.}
\label{claim:BNV7}
Under the original \texttt{Gap\allowbreak Big}, there exists $h$ for which all
$b_{j+1}-b_j\ge h$, $\texttt{Gap\allowbreak Lap}(\epsilon,\lambda,h)$ holds, and
(BNV.12) holds.

\begin{proof}
The original gap supplies an odd natural $T$ such that $T\ge n+5$,
$n+T\ge1801$, and
\[
 b_{j+1}-b_j>2\bigl(n+(2r+2)T\bigr)^3+T+3.
 \tag{BNV.13}
\]
Choose $h$ to be the minimum of these finitely many differences.  Put
$Q=r+1=q=q_0$ and $B=2(n+2QT)^3$.  Since $n\ge2$ and $Q\ge1$,
$(n+5)^3\ge17n^2$; hence
\[
 B\ge16Q^3T^3\ge272Q^3n^2>160Q^2n^2.
\]
Also $\sqrt{2Q}\le2Q$ and
$\max(2nQ,2(Q-1)n^2)\le2Qn^2$.  The third member of (BNV.11) is
therefore at most
$128Q^2n^2+16Q\le132Q^2n^2$.
The second is at most $16nQ+16Q+2\le25Q^2n^2$ for $n\ge2$;
the first is at most $3+n^2Q+3n/2\le3Q^2n^2$ for $n\ge2$.
The loose bound $160Q^2n^2$ thus exceeds all three.  Equation (BNV.13)
gives $h>h_{\rm need}$, and BNV3 gives the noncompact gaps of $\lambda$.
Finally the first and third members of (BNV.11), with $q-1=r$, give
(BNV.12) directly.  The parity and $1801$ hypotheses are retained as
original input; this elementary numerical subproof needs only $T\ge n+5$.
\end{proof}
% END BNV7

% BEGIN BNV8
\paragraph{BNV8. The full large-parameter family and its lowest weights.}
\label{claim:BNV8}
Let $d_a=0$ for $a<k$ and $d_{k-1+j}=-j$ for $1\le j\le2r+1$.
For $\Lambda_t=\lambda^{(t)}+\rho_n-\rho_c$ one has
\[
 \Lambda_t=t d+\Lambda_0,\qquad
 (\Lambda_t)_a=u_a-(k-1)/2+a\quad(0\le a<k).
 \tag{BNV.14}
\]
For all sufficiently large integral $t$, $\lambda^{(t)}$ is
Harish--Chandra and
$\texttt{Gap\allowbreak Lap}(\epsilon,\lambda^{(t)},t/2)$ holds.  This is the original
\texttt{Lam\allowbreak TStatement}, including the compact highest-weight identification.

\begin{proof}
The affine formula is immediate from (BNV.7).  At an initial positive
coordinate $a<k$, there are $r+1$ noncompact positive roots leaving $a$,
and the compact half-sum is $(r+k-1-2a)/2$.  Thus
$(\rho_n-\rho_c)_a=(2-k+2a)/2$ and (BNV.14) follows.
This is the compact $U(k)$ highest weight $u-\rho_k$.  BNV2 proves that
$u$ is itself a Harish--Chandra parameter.

The lattice condition for (BNV.7) follows from half-integrality and the
integral displacements $jt$.  The initial block is strictly decreasing.
Every remaining adjacent difference is either a fixed real number plus
$t$, or (at its boundary with the initial block) a fixed real number plus
$t$.  Consequently all such differences are positive and at least $t/2$
for all sufficiently large $t$.  These are all adjacent sign changes;
BNV3 applies.  Increasing the finite threshold so that
$t/2>h_{\rm need}(\epsilon)$ proves the assertion.  This also gives
$q_0=r+1\ge1$, as required by the original statement.
\end{proof}
% END BNV8

\subsection{An explicit Cayley matrix matched to the character}

% BEGIN BNV9
\paragraph{BNV9. The cyclic isotropic frame by Vandermonde interpolation.}
\label{claim:BNV9}
Let $x_0>\cdots>x_{2r}$ be real, put
$D_x=\operatorname{diag}(x_s)$, and equip $\mathbb R^{2r+1}$ with the form
$[a,b]=\sum_s(-1)^{s+1}a_sb_s$.  There are vectors
$e_0,\ldots,e_{r-1},z,f_0,\ldots,f_{r-1}$ with
\[
 [e_a,e_b]=[f_a,f_b]=[e_a,z]=[f_a,z]=0,\quad
 [f_a,e_b]=\delta_{ab},\quad [z,z]=-1,
\]
such that $D_xe_a=e_{a+1}$ for $a<r-1$ and $D_xe_{r-1}=z$.

\begin{proof}
Let $P_s=\prod_{j\ne s}(x_s-x_j)$, which has sign $(-1)^s$, and set
$c_s=\sqrt{(-1)^s/P_s}>0$ and $g_i=(x_s^ic_s)_s$.
Comparison of leading coefficients in the Lagrange interpolation formula
for $X^m$, $0\le m\le2r$, gives
\[
 \sum_{s=0}^{2r}\frac{x_s^m}{P_s}
 =\begin{cases}0&m<2r,\\1&m=2r.\end{cases}
\]
Hence $[g_i,g_j]=0$ if $i+j<2r$ and $[g_i,g_j]=-1$ if $i+j=2r$.
Take $e_a=g_a$ and $z=g_r$.  The Vandermonde determinant and $c_s\ne0$
show that $g_0,\ldots,g_{2r}$ form a basis.  Nondegeneracy provides
$f'_a$ with $[f'_a,g_i]=\delta_{ai}$ for $0\le i\le r$; extend these
conditions arbitrarily to the basis.  If $Q_{ab}=[f'_a,f'_b]$, replace
\[
 f_a=f'_a-\tfrac12\sum_b Q_{ab}e_b.
\]
The symmetric matrix $Q$ then gives $[f_a,f_b]=Q_{ab}-Q_{ab}/2-Q_{ba}/2=0$;
all other required pairings are unchanged.  Finally $D_xg_i=g_{i+1}$.
For $r=0$ this is simply a negative unit vector $z$, so the same
construction includes that case.
\end{proof}
% END BNV9

% BEGIN BNV10
\paragraph{BNV10. The exact trace-character identity.}
\label{claim:BNV10}
There is a Cayley matrix $P$ sending the $k$ positive middle standard
vectors to the first $k$ positive standard vectors, such that, for every
$Y\in\mathfrak n=\operatorname{Lie}(N)$,
\[
 \operatorname{tr}(D PYP^{-1})
 =\sum_{a=0}^{r-1}\bigl(\ell_a(Y)-\overline{\ell_a(Y)}\bigr),
 \quad D=\operatorname{diag}(d).
 \tag{BNV.15}
\]
With $z_0$ the standard negative middle vector and $c_a(t)=-2it$, one has
\[
 \psi_{z_0,c(t)}(\exp Y)
   =\exp\bigl(t\operatorname{tr}(D PYP^{-1})\bigr).
 \tag{BNV.16}
\]
These are the full original \texttt{Pairing\allowbreak Properties\allowbreak Statement}, including
the source condition \texttt{WCols}.

\begin{proof}
Apply BNV9 to $x_s=-(s+1)$ in the last $2r+1$ coordinates of
$\mathbb C^n$, whose signs are $-,+,-,\ldots,-$.  Extend the frame by the
first $k$ positive standard vectors.  Use its $e_a$, these positive
vectors, $z$, and the appropriately reversed $f_a$ as columns of $P$.
Their Gram matrix is exactly $H$, so $P^*J_\epsilon P=H$, and the
positive columns have the required prescribed values.

Here is also a block verification of the trace, to fix its sign and
normalization.  Put $\Xi=P^{-1}DP$, and let $A$ send $e_a$ to
$e_{a+1}$ (or to $z$ for $a=r-1$), and vanish on the other columns.
Both $\Xi$ and $A+HA^*H$ are $H$-self-adjoint and have the same columns
on $E$ and on the positive middle space.  Their difference $R$ kills
these columns.  Self-adjointness consequently kills its rows on the
dual $F$ space and on the positive middle space.  In the block upper
nilpotent Lie algebra $\mathfrak n$, all entries that could contribute
to $\operatorname{tr}(RY)$ are therefore zero (the remaining $z,z$
entry also meets a zero diagonal entry of $Y$).

Now $\operatorname{tr}(AY)=\sum_a\ell_a(Y)$: these are just the entries
opposite the $e_a\mapsto e_{a+1},z$ arrows.  Since $HY^*H=-Y$,
\[
 \operatorname{tr}(HA^*HY)
 =-\overline{\operatorname{tr}(AY)}.
\]
This proves (BNV.15).  Each $\ell_a$ is a coordinate in the first
nilpotent filtration degree, so $(\exp Y-I)$ and $Y$ have the same
$\ell_a$ coordinates; terms $Y^j$, $j\ge2$, cannot contribute.
Finally, for $w\in\mathbb C$, $i\operatorname{Re}(-2itw)=t(w-\bar w)$.
Substituting into the original character (BNV.1) proves (BNV.16).
No unspecified scalar multiple of the character remains.
\end{proof}
% END BNV10

% BEGIN BNV11
\paragraph{BNV11. The exact exponential Haar measure on the actual $N$.}
\label{claim:BNV11}
The matrix exponential is a polynomial diffeomorphism
$\mathfrak n\longrightarrow N$.  Its inverse is the finite matrix logarithm.
The pushforward of Euclidean Lebesgue measure on $\mathfrak n$, for the
Frobenius inner product, is Haar measure on $N$, precisely the source's
$\texttt{haar\allowbreak Lie}(N,1)$.  Moreover
\[
 \|Y\|_{\rm F}\le C\|\exp Y-I\|_{\rm F}
             (1+\|\exp Y-I\|_{\rm F})^{n-2}
 \tag{BNV.17}
\]
when $n\ge2$; if the algebra is zero the assertion is immediate.

\begin{proof}
Choose the basis order refining the isotropic flag and its reversed dual.
The endomorphisms in $\mathfrak n$ are strictly upper triangular in the
corresponding block filtration, and $N-I$ has the same property.  Thus
$Y^n=0$ and $(g-I)^n=0$.  The finite exponential and logarithm polynomials
are inverse by their formal power-series identities, valid after truncation.
The logarithm of a group element satisfies $H(\log g)^*H=-\log g$:
apply the polynomial logarithm to $Hg^*H=g^{-1}$ and use
$\log(g^{-1})=-\log g$.  The flag conditions are preserved by these
polynomials.  This proves both inclusions and the inverse assertion.
The logarithm formula
$Y=\sum_{j=1}^{n-1}(-1)^{j+1}(g-I)^j/j$
and submultiplicativity of Frobenius norm give (BNV.17), with an enlarged
constant when necessary.

In exponential coordinates multiplication is the finite
Baker--Campbell--Hausdorff polynomial.  Order coordinates by a basis
adapted to the lower central filtration
$\mathfrak n\supset[\mathfrak n,\mathfrak n]\supset\cdots$.
For fixed $Y$, the coordinate of $\log(\exp Y\exp Z)$ in filtration
degree $j$ is $Z_j+Y_j$ plus a polynomial in lower-degree coordinates
of $Z$.  Its derivative in $Z$ is block triangular with identity diagonal,
so its real Jacobian is $1$.  Change of variables proves invariance of
the pushed-forward Lebesgue measure.  It is nonzero and locally finite,
as exp is a diffeomorphism, hence Haar.  The source realification uses
exactly the real and imaginary matrix entries, whose Euclidean norm is
Frobenius norm; its prescribed pushforward with scalar $1$ is therefore
this same measure.  This proves the whole original
\texttt{Nset\allowbreak Exp\allowbreak Injective\allowbreak Surjective} package, including its norm estimate.
\end{proof}
% END BNV11

% BEGIN BNV12
\paragraph{BNV12. A global Cayley majorant from unipotence and polar decomposition.}
\label{claim:BNV12}
For every Cayley matrix $P$ there is $c>0$ such that, whenever
$Y\in\mathfrak n$ and $P\exp(Y)P^{-1}=\kappa\exp X$ is its polar
decomposition in $G_\epsilon$, with $X^*=X$ and $J_\epsilon XJ_\epsilon=-X$,
\[
 c\log(1+\|Y\|_{\rm F}^2)
 \le\min\left(\|X\|_{\rm F}^2,
                  \frac{\|X\|_{\rm F}}{\sqrt{2q_0}}\right).
 \tag{BNV.18}
\]
This is the full original \texttt{Cayley\allowbreak Majorant\allowbreak Bound}.

\begin{proof}
Suppose first that $q_0>0$, and write $s=\|X\|_{\rm F}$ and
$g=P\exp(Y)P^{-1}$.  Unipotence gives $\operatorname{tr}g=n$.
The nonzero eigenvalues of $X$ occur as $s_i,-s_i$, so
\[
 \begin{split}
 \|g-I\|_{\rm F}^2
 &=\operatorname{tr}(g^*g)-n
   =\operatorname{tr}(e^{2X})-n
   =4\sum_i\sinh^2s_i
 \le2s^2e^{2s}.
 \end{split}
\]
The finite logarithm of $g$ is $PYP^{-1}$.  Its polynomial formula gives
$\|Y\|_{\rm F}\le C s e^{Cs}$ for a constant depending on $P,n$.
For $0\le s\le1$, this implies
$\log(1+\|Y\|_{\rm F}^2)\le C_1s^2$; for $s\ge1$ it implies
$\log(1+\|Y\|_{\rm F}^2)\le C_2s$.
The quantity $\min(s^2,s)$ is at most
$\max(1,\sqrt{2q_0})\min(s^2,s/\sqrt{2q_0})$.
Taking a sufficiently small positive $c$ proves (BNV.18).

If $q_0=0$, $G_\epsilon$ is compact.  Its only unipotent element is $I$,
since a unitary matrix is diagonalizable.  Thus $Y=0$ and $X=0$.
The asserted inequality is $0\le0$ with the source's zero-division
convention; equivalently it can be stated without division in this
degenerate case.  No noncompact estimate is needed there.
\end{proof}
% END BNV12

\subsection{The Gaussian limit of the actual highest-weight coefficient}

For an arbitrary signature $\epsilon$ and real vectors $d,\Lambda_0$, put
$D=\operatorname{diag}(d)$ and $X^\dagger=J_\epsilon X^*J_\epsilon$.
Let $\mathcal N_+$ consist of the matrix positions $(a,b)$ such that either
$a<b$, $\epsilon_a=\epsilon_b$, $d_a\ne d_b$, or
$b<a$, $\epsilon_a\ne\epsilon_b$.
Let $\mathcal N_-$ be its transpose, and let $\mathcal L$ consist of the
positions of equal sign and equal $d$ value.  These three sets partition
all positions.  Write $Y=Y_L+Y_++Y_-$ for the corresponding projections,
and put
\[
 \tau(Y)=\operatorname{tr}(DY),\qquad
 Q(Y)=-\operatorname{tr}(D[Y,Y_-]).
 \tag{BNV.19}
\]
The subscript $-$ here refers to this explicitly specified projection;
it is not a choice of a new positive-root system.

% BEGIN BNV13
\paragraph{BNV13. Uniform highest-weight estimates from the Casimir.}
\label{claim:BNV13}
Suppose $E_{ab}$ act on an inner product space, obey the
$\mathfrak{gl}_n$ commutator relations and
$E_{ab}^*=\epsilon_a\epsilon_b E_{ba}$ on the common invariant domain.
Let $w$ be unit, compact highest of weight
$\Lambda=\lambda+\rho_n-\rho_c$, and have Casimir scalar
$\|\lambda\|^2-\|\rho\|^2$.  If $t\ge1$ and
$\Lambda=td+\Lambda_0$, then a constant $C$ depending only on
$\epsilon,d,\Lambda_0$ satisfies
\[
 \begin{aligned}
 &\|E_{ab}w\|\le C &&((a,b)\in\mathcal N_+),\\
 &\|E_{ab}w-\delta_{ab}td_aw\|\le C
                  &&((a,b)\in\mathcal L),\\
 &\left\|E(X)w-t\operatorname{tr}(DX)w\right\|
       \le C\sqrt t\sum_{a,b}|X_{ab}| &&(X\in M_n(\mathbb C)).
 \end{aligned}
 \tag{BNV.20}
\]
These are the full original \texttt{Highest\allowbreak Weight\allowbreak Vector\allowbreak Uniform\allowbreak Bounds}.

\begin{proof}
For $a<b$ of the same sign, compact highest weight and the commutator
give $\|E_{ba}w\|^2=\Lambda_a-\Lambda_b$.
For opposite signs the adjoint relation and the same commutator give
\[
 \|E_{ab}w\|^2=\|E_{ba}w\|^2+\Lambda_a-\Lambda_b.
\]
The Casimir identity in DC8, specialized to this highest weight, is
\[
 2\sum_{a<b,\ \epsilon_a\ne\epsilon_b}\|E_{ba}w\|^2
       =\|\rho\|^2-\|\rho_n-\rho_c\|^2=:K_0.
 \tag{BNV.21}
\]
For completeness its cancellation uses
$\Lambda+\rho_c-\rho_n=\lambda$: substitution in DC8 cancels
$\|\lambda\|^2$ from both sides and leaves exactly (BNV.21).
Thus every noncompact lower operator has a bound independent of $t$.
The compact upper operators vanish on $w$.  These give the first line
of (BNV.20).  In $\mathcal L$, compact lower operators have squared norm
$\Lambda_{0,a}-\Lambda_{0,b}$ because $d_a=d_b$; the upper ones vanish,
and diagonal operators leave residual $\Lambda_{0,a}w$.  This gives the
second line.

For every remaining off-diagonal position, the displayed norm identities
give a bound of the form
$\|E_{ab}w\|^2\le |K_0|+t|d_a-d_b|+|\Lambda_{0,a}-\Lambda_{0,b}|$.
The diagonal residuals are constant.  Since there are finitely many
positions, one constant $C$ bounds every centered operator by $C\sqrt t$.
Expand $E(X)=\sum X_{ab}E_{ab}$ and use the triangle inequality; this
proves the last line.  The argument did not assume $d$ is decreasing,
and hence has the original universal $d,\Lambda_0$ scope.
\end{proof}
% END BNV13

% BEGIN BNV14
\paragraph{BNV14. The whole original tangent-pairing estimate.}
\label{claim:BNV14}
Let $\sigma$ be any unitary representation of $G_\epsilon$ and $w$ any
smooth unit vector satisfying the highest-weight and Casimir hypotheses
of BNV13, with $\Lambda=td+\Lambda_0$ and $t\ge1$.
For every $R$ there is a constant $C_R$, depending only on
$\epsilon,d,\Lambda_0,R$, such that for every
$Y\in\mathfrak g_\epsilon$ with $\|Y\|_{\rm F}\le |R|$,
\[
 \left|
 \langle w,\sigma(\exp(Y/\sqrt t))w\rangle
       e^{-\sqrt t\tau(Y)}-e^{-Q(Y)/2}
 \right|\le\frac{C_R}{\sqrt t}.
 \tag{BNV.22}
\]
In particular the constant is uniform in the representation and its
highest vector.  This is the full original \texttt{Tangent\allowbreak Pairing\allowbreak Statement}.

\begin{proof}
Differentiate unitarity along real one-parameter subgroups and complexify.
The derived operators on smooth vectors then obey exactly the adjoint
relation in BNV13.  The chain rule gives, for
$U(s)=\langle w,\sigma(\exp(sY/\sqrt t))w\rangle$,
\[
 U'(s)=t^{-1/2}\langle w,\sigma(g_s)E(Y)w\rangle,
 \quad g_s=\exp(sY/\sqrt t).
\]
The conjugation rule for the derived representation and the adjoint
identity give
\[
 \langle w,\sigma(g_s)E(Y_-)w\rangle
 =\langle E((g_sY_-g_s^{-1})^\dagger)w,\sigma(g_s)w\rangle.
 \tag{BNV.23}
\]
All vectors in these identities are smooth.  Finite-dimensional matrix
Taylor expansion, with its second derivative bounded on the fixed ball,
gives uniformly for $0\le s\le1$
\[
 g_sY_-g_s^{-1}
   =Y_-+\frac{s}{\sqrt t}[Y,Y_-]+R_s,
 \qquad \|R_s\|\le C_R/t.
 \tag{BNV.24}
\]
Any fixed matrix norm can be used here; equivalence of the finite
dimensional norms only changes $C_R$.

We spell out the error estimate, since it is what permits the large-$t$
limit for a varying family of representations.  The $Y_L$ contribution
to $U'$ is $\sqrt t\tau(Y)U$ with an $O_R(t^{-1/2})$ error, by the
second line of (BNV.20); $\tau(Y_L)=\tau(Y)$.
The $Y_+$ contribution is $O_R(t^{-1/2})$ by its first line.  In
(BNV.23), $Y_-^\dagger$ belongs to $\mathcal N_+$ and gives another
$O_R(t^{-1/2})$ error.  Put $Z_s=(s/\sqrt t)[Y,Y_-]+R_s$.
Its centered action has norm at most
$C\sqrt t\|Z_s^\dagger\|=O_R(1)$, and hence contributes
$O_R(t^{-1/2})$ to $U'$.  Finally its scalar action contributes
\[
 \sqrt t\,\tau(Z_s)U(s)
   =s\operatorname{tr}(D[Y,Y_-])U(s)+O_R(t^{-1/2}).
\]
Here conjugate linearity in the first inner-product variable cancels the
conjugation in $\tau(Z_s^\dagger)=\overline{\tau(Z_s)}$.
The norm of $\sigma(g_s)w$ is $1$, so these estimates require no further
operator bounds.  Thus
\[
 U'(s)=(\sqrt t\tau(Y)-sQ(Y))U(s)+\delta_t(s),
 \qquad |\delta_t(s)|\le C_R/\sqrt t.
 \tag{BNV.25}
\]
Since $Y^\dagger=-Y$ and $D$ is real diagonal, $\tau(Y)$ is purely
imaginary.  Multiply (BNV.25) by
$\exp(-s\sqrt t\tau(Y)+s^2Q(Y)/2)$ and integrate from $0$ to $1$.
The factor involving $\tau$ has absolute value $1$, $U(0)=1$, and
$|Q(Y)|$ is bounded on the fixed ball.  This gives (BNV.22), after
enlarging $C_R$ by $\exp(\sup|Q|)$.  This derivation also proves all
differentiation, conjugation, and second-order estimates used here,
rather than treating a Gaussian approximation as an independent premise.
\end{proof}
% END BNV14

% BEGIN BNV15
\paragraph{BNV15. A nondegenerate complex Gaussian has nonzero integral.}
\label{claim:BNV15}
Let $B:E\times E\to\mathbb C$ be real bilinear on a finite-dimensional
real space and suppose $\operatorname{Re}B(y,y)>0$ for $y\ne0$.
For every additive Haar measure $dy$,
\[
 \int_E e^{-B(y,y)}\,dy\ne0.
 \tag{BNV.26}
\]

\begin{proof}
Only the symmetric part matters.  In linear coordinates write it
$A+iC$ with $A,C$ real symmetric and $A$ positive definite.
Changing variables by $A^{-1/2}$ and then an orthogonal diagonalization
of $A^{-1/2}CA^{-1/2}$ gives eigenvalues $\alpha_j\in\mathbb R$ and
\[
 \int_E e^{-B(y,y)}dy
   =c\,\pi^{d/2}(\det A)^{-1/2}
           \prod_{j=1}^{d}(1+i\alpha_j)^{-1/2},\quad c>0.
 \tag{BNV.27}
\]
The square roots are chosen continuously on the right half-plane with
positive real part.  One elementary justification of the one-dimensional
formula is as follows.  For $\operatorname{Re}z>0$, dominated
differentiation of $F(z)=\int_{\mathbb R}e^{-zx^2}dx$ and integration
by parts give $F'(z)=-F(z)/(2z)$.  Hence $F(z)\sqrt z$ is constant
on that half-plane; at $z=1$ its value is $\sqrt\pi$, by the real
Gaussian integral.  Fubini applies because the real part of the quadratic
form is positive.  Every factor in (BNV.27) is nonzero.  In dimension
zero the integral is the positive mass of the one-point Haar space.
\end{proof}
% END BNV15

% BEGIN BNV16
\paragraph{BNV16. The whole original \texttt{Gaussian\allowbreak Limit\allowbreak Haar}.}
\label{claim:BNV16}
Let $E$ be a finite-dimensional real normed space with additive Haar
measure, and $f_t:E\to\mathbb C$, $t\in\mathbb N$, eventually
almost-everywhere strongly measurable.
Suppose the first assertion below is pointwise in every $y$, and the
second has one uniform threshold $t_1$: it holds for every $t\ge t_1$
and every $y\in E$:
\[
 f_t(y/\sqrt t)\longrightarrow e^{-B(y,y)},
 \qquad |f_t(y)|\le(1+c\|y\|^2)^{-\gamma t}
 \tag{BNV.28}
\]
Here $B$ is complex-valued real bilinear and $c,\gamma>0$ are fixed.
Then for all sufficiently large $t$, $f_t$ is integrable and
$\int_E f_t\ne0$.

\begin{proof}
Choose an integer $N\ge1$ beyond the eventual bounds and so large that
$2\gamma N>\dim E$.  The function $s\mapsto(1+a/s)^s$ is increasing
for $s>0$, $a\ge0$, since its logarithmic derivative is
$\log(1+a/s)-(a/s)/(1+a/s)\ge0$.
Consequently for $t\ge N$
\[
 |f_t(y/\sqrt t)|\le
       (1+(c/N)\|y\|^2)^{-\gamma N}.
 \tag{BNV.29}
\]
This is integrable by radial comparison and equivalence of norms.
Each positive scalar dilation is nonsingular for additive Haar measure,
so the scaled functions retain almost-everywhere strong measurability;
pointwise Borel measurability is not required.
Taking limits in (BNV.29) shows
\[
 e^{-\operatorname{Re}B(y,y)}
 \le(1+(c/N)\|y\|^2)^{-\gamma N}<1\quad(y\ne0).
\]
Thus the real part of the quadratic form is positive definite; no
positivity of $B$ was silently added to the original hypotheses.
Dominated convergence and BNV15 now give
\[
 \int_E f_t(y/\sqrt t)dy
   \longrightarrow \int_E e^{-B(y,y)}dy\ne0.
\]
These integrals are nonzero for all sufficiently large $t$.  Scalar
dilation of additive Haar measure gives their exact relation to the
unscaled integrals,
$\int f_t(y/\sqrt t)dy=t^{\dim E/2}\int f_t(y)dy$.
Both integrability and nonvanishing follow.  The same argument, with
dilation factor $1$, covers dimension zero.
\end{proof}
% END BNV16

\subsection{Compact orthogonality and the complete Laplace base}

% BEGIN BNV17
\paragraph{BNV17. The exact compact Schur factor in the Bessel integral.}
\label{claim:BNV17}
Let $P$ have the prescribed positive middle columns of BNV10.
Let $T$ be any compact highest-weight type whose initial $k$ weights
agree with the compact $U(k)$ highest weight of $\pi$, and let $w\in T$,
$\xi\in V_\pi$ be unit highest vectors of the corresponding weights.
Take probability Haar measure on $H_{z_0}$ and any $\sigma$-finite
measure $\mu_N$ on the ambient matrix space $M_n(\mathbb C)$, with
no support assumption.  Use the original kernel $\Psi$ on $G_\epsilon$,
extended by zero outside that closed group.  Suppose it is continuous
on $G_\epsilon$, compact bi-equivariant there, and
$m\mapsto\|\Psi(PmP^{-1})\|$ is $\mu_N$-integrable.  In this claim
$A$ means the formula (BNV.5) with its second integral over the ambient
matrix space; it agrees with the preceding notation when $\mu_N$ is
supported on $N$.  Then
\[
 \langle\xi,A(w,w,\xi;1,1)\rangle
 =\frac1{\dim V_\pi}
       \int_{M_n(\mathbb C)}\langle w,\Psi(PmP^{-1})w\rangle
                     \overline{\psi_{z_0,c}(m)}d\mu_N(m).
 \tag{BNV.30}
\]
This is the original \texttt{Whittaker\allowbreak Schur\allowbreak Orthogonality}, with its exact
normalization and its actual vector integral.

\begin{proof}
The subgroup $PH_{z_0}P^{-1}$ acts only on the first $k$ positive
coordinates.  The cyclic compact subrepresentation it generates from
$w$ is the irreducible highest-weight $U(k)$ type of $\pi$.
One can see the cyclic assertion directly from complete reducibility:
in each irreducible summand a nonzero component of $w$ is killed by the
positive-root operators and has the specified weight, so that summand
has that highest weight.  The cyclic span of this one vector is a single
diagonal copy of the type.  Finite compact highest-weight uniqueness
therefore provides an equivariant isometry $J:V_\pi\to T$ with $J\xi=w$.

For any $Y\in T$, average the rank-one operator $\xi\otimes\xi^*$
under conjugation by $\pi$.  The average commutes with the irreducible
$\pi$, hence is scalar; its trace is $1$, so the scalar is
$1/\dim V_\pi$.  Taking the appropriate matrix coefficient gives
\[
 \int_{H_{z_0}}
 \langle\xi,\pi(h)^{-1}\xi\rangle
 \langle w,T(PhP^{-1})Y\rangle,dh
      =\frac{\langle w,Y\rangle}{\dim V_\pi}.
\]
The projection of $Y$ onto $J(V_\pi)$ may be inserted without changing
either side.  Compact equivariance of the zero extension holds for
every ambient matrix: $\kappa g\in G_\epsilon$ if and only if
$g\in G_\epsilon$, so outside the group both sides are zero.
The extension is Borel measurable because the group is closed and
the kernel is continuous on it.  Moreover (BNV.1), as an exponential
of $i$ times a real number, has modulus one on every ambient matrix,
not just on $N$.  Apply the averaged identity with
$Y=\Psi(PmP^{-1})w$ for every $m$.  Unitarity bounds the norm of the
double integrand by $\|\Psi(PmP^{-1})\|$, so the hypotheses give
absolute integrability for the stated arbitrary ambient measure.
Fubini and continuity of the scalar pairing yield
(BNV.30).  Compact Haar probability exists by normalizing any Haar
measure, whose total mass is finite and positive.
\end{proof}
% END BNV17

% BEGIN BNV18
\paragraph{BNV18. The whole original \texttt{Step\allowbreak IBase\allowbreak Statement}.}
\label{claim:BNV18}
For every $k\ge1$, $r\ge0$, and $u,v,b$ satisfying the Harish--Chandra,
interlacing, and half-integrality assumptions (BNV.2)--(BNV.3), there is
$t_0\in\mathbb N$ such that
\[
 \texttt{Bessel\allowbreak NZ}(\epsilon,\lambda^{(t)},u)
          \quad\text{for every integer }t\ge t_0.
 \tag{BNV.31}
\]
No lower-gap assumption on the initial $b$ is added.

\begin{proof}
By BNV8, all sufficiently large $t$ have analytic gap $t/2$.
Construct the compact type $T_t$, normalized kernel $\Psi_t$, and
smooth unitary model of DC1--DC11 for this parameter.  After dividing
the embedding by the square root of the positive scalar in DC11,
the unit highest vector $w_t\in T_t$ gives the exact coefficient
\[
 \langle w_t,\Psi_t(g)w_t\rangle
    =\langle\widetilde w_t,\sigma_t(g)\widetilde w_t\rangle,
\]
where $\widetilde w_t$ is smooth, unit, of lowest compact highest weight
$\Lambda_t$, and has the required Casimir.  The gap hypothesis here
is stronger than DC's construction threshold by (BNV.11).

Fix the one matrix $P$ of BNV10, independently of $t$.
On the real vector space $\mathfrak n$ define
\[
 f_t(Y)=\langle w_t,\Psi_t(P\exp YP^{-1})w_t\rangle
       \exp\bigl(-t\operatorname{tr}(DPYP^{-1})\bigr).
 \tag{BNV.32}
\]
The phase has modulus $1$.  By BNV10 it is precisely the conjugate
of the character $\psi_{z_0,c(t)}(\exp Y)$ with every
$c_a(t)=-2it\ne0$ for $t\ge1$.  Applying BNV14 to $PYP^{-1}$ gives
\[
 f_t(Y/\sqrt t)\longrightarrow
     \exp\bigl(-Q(PYP^{-1})/2\bigr).
 \tag{BNV.33}
\]
The exponent is the diagonal of a real bilinear form on $\mathfrak n$:
one may take
$B(Y,Z)=-\operatorname{tr}(D[PYP^{-1},(PZP^{-1})_-])/2$.
All its data are fixed as $t$ varies.

For $P\exp YP^{-1}=\kappa e^X$, DC6 and BNV12 give the global estimate
\[
 \begin{split}
 \|\Psi_t(P\exp YP^{-1})\|
 &\le\exp\left[-\frac{t}{16}
       \min(\|X\|_{\rm F}^2,\|X\|_{\rm F}/\sqrt{2q_0})\right]\\
 &\le(1+\|Y\|_{\rm F}^2)^{-ct/16}.
 \end{split}
 \tag{BNV.34}
\]
This bound is uniform in $t$; in particular no dimension of $T_t$
occurs in it.  BNV16 applied to (BNV.32)--(BNV.34) proves that the
scalar integral of $f_t$ is nonzero for every sufficiently large $t$.
The same bound makes the full operator norm in (BNV.34) integrable
once $ct/8>\dim\mathfrak n$.

Choose $\mu_N$ to be the exact exponential Haar measure of BNV11,
$\mu_W$ the probability Haar on $H_{z_0}$, and $\pi$ the compact
$U(k)$ type with Harish--Chandra parameter $u$, realized on $H_{z_0}$.
BNV2 gives its integral highest weight, and BNV8 identifies that
weight with the initial $k$ entries of $\Lambda_t$.
Choose a unit highest vector $\xi$.  BNV17 gives
\[
 \langle\xi,A(w_t,w_t,\xi;1,1)\rangle
       =\frac1{\dim V_\pi}\int_{\mathfrak n}f_t(Y)dY\ne0.
\]
All these are legitimate witnesses for the original \texttt{Bessel\allowbreak NZ}:
the character is nondegenerate, the Cayley matrix has the prescribed
form, the measures have the original Haar predicates, and both group
points are $1$.  This proves (BNV.31) with every quantifier of the
source base statement.  If $r=0$, $\mathfrak n$ is zero-dimensional;
the same argument reduces to the nonzero compact Schur factor, with
no character coordinates to check.
\end{proof}
% END BNV18

\subsection{From one nonzero model period to every original datum}

The next four claims are proved for a Harish--Chandra parameter $\eta$,
a signature with $(p,q)=(k+r,r+1)$, $k\ge1$, and some
$\texttt{Gap\allowbreak Lap}(\epsilon,\eta,h)$.  They do not require that $\eta$
have the special form (BNV.2).  In this scope $q_0=r+1$.

% BEGIN BNV19
\paragraph{BNV19. Absolute convergence for every original coefficient period.}
\label{claim:BNV19}
For all data permitted in (BNV.5), every $v_1,v_2,\xi$ and every
$x_0,x\in G$, its integrand is Bochner integrable.  The analogous
endomorphism integrand is also integrable.  This holds for every original
Haar pair and every Cayley matrix $P$.

\begin{proof}
Let $c_0=\sqrt{2q_0}$ and $a=h/(8c_0)$.  In polar coordinates
$g=\kappa e^X$, $s=\|X\|_{\rm F}$, DC6 gives
\[
 \|\Psi(g)\|\le e^{-(h/8)\min(s^2,s/c_0)}
                 \le C e^{-as}.
\]
Indeed $s^2\ge s/c_0-1/(4c_0^2)$, so one may take
$C=e^{h/(32c_0^2)}$.  Also $\|g\|_{\rm F}\le\sqrt n e^s$.
Consequently
\[
 \|\Psi(PgP^{-1})\|\le C_P\|g\|_{\rm F}^{-a}.
 \tag{BNV.35}
\]
Conjugation by a fixed invertible $P$ changes Frobenius norm only by
fixed positive comparison constants.  The third member of (BNV.11)
implies
\[
 a>4\max(2nq_0,2rn^2)+1>2rn^2=b_S.
\]
Apply BNV1 to $g=x_0^{-1}hn'x$.  On the measure's support,
$|\psi(n')|=1$ and $\|\pi(h)^{-1}\xi\|=\|\xi\|$, while
$\|\pi(h)^{-1}\|=1$.  Thus each asserted norm is bounded by a fixed
constant times the integrable function in BNV1.  Continuity on the
groups gives strong measurability; finite-dimensional ranges of the
kernel and of $\pi$ suffice for the operator-valued assertion.
This proves actual absolute convergence, including every translation
and every subsequent finite linear combination used below.
\end{proof}
% END BNV19

% BEGIN BNV20
\paragraph{BNV20. Explicit conjugacy of all allowed Bessel characters and compact types.}
\label{claim:BNV20}
For any two allowed pairs $(z_1,c_1)$ and $(z_2,c_2)$ there is
$a\in G$ normalizing $N$ such that
\[
 aH_{z_1}a^{-1}=H_{z_2},\qquad
 \psi_{z_2,c_2}(an'a^{-1})=\psi_{z_1,c_1}(n')\quad(n'\in N).
 \tag{BNV.36}
\]
The transport preserves the compact Harish--Chandra label $u$.
Thus if $\pi_1,\pi_2$ have that label, there is an isometry
$J:V_{\pi_1}\to V_{\pi_2}$ intertwining
$\pi_1(h)$ with $\pi_2(aha^{-1})$.

\begin{proof}
Write $T_i=-z_i^*Hz_i$ and $s=\sqrt{T_1/T_2}>0$.
The negative unit vectors $z_i/\sqrt{T_i}$ have positive-definite
orthogonal complements in $V$.  Choosing orthonormal bases of those
complements constructs $M\in U(k,1)$ with $Mz_1=sz_2$.
Let $D_E=\operatorname{diag}(d_0,\ldots,d_{r-1})$, and use the Levi
element
\[
 a=\operatorname{diag}(D_E,M,D_E^{-*}),
\]
with the dual ordering adjusted as for $H$.  It normalizes the isotropic
flag and $N$.  For $r>0$ choose
\[
 d_{r-1}=s\,c_{1,r-1}/c_{2,r-1},\qquad
 d_j=(c_{1,j}/c_{2,j})d_{j+1}\quad(j=r-2,\ldots,0).
\]
All are nonzero.  Direct entry multiplication gives
\[
 \ell_j(an'a^{-1})=(d_j/d_{j+1})\ell_j(n')\quad(j<r-1),
 \qquad
 \ell_{r-1,z_2}(an'a^{-1})=(d_{r-1}/s)\ell_{r-1,z_1}(n').
\]
Thus $c_{2,j}\ell_{j,z_2}(an'a^{-1})=c_{1,j}\ell_{j,z_1}(n')$
for every $j$, proving the character identity exactly, before taking
real parts or exponentials.  The fixed spaces of $H_{z_i}$ also
transform as asserted.  For $r=0$ only $M$ is needed.

For the label assertion choose a positive orthonormal frame in
$z_1^\perp\cap V$; its image under $M$ is such a frame for $z_2$.
In these frames conjugation by $a$ is the identity identification of
the compact $U(k)$ factors, because $D_E$ acts on the fixed outer
coordinates.  Another allowed frame differs by an element of $U(k)$,
which induces an inner automorphism and does not change the compact
type.  Finite compact highest-weight uniqueness supplies the unitary
intertwiner $J$.  This checks the frame clause of the original
$\texttt{Is\allowbreak HCType\allowbreak W}$, rather than presuming that an arbitrary group
isomorphism preserves its label.
\end{proof}
% END BNV20

% BEGIN BNV21
\paragraph{BNV21. The full original transport $\texttt{Bessel\allowbreak NZ}\Rightarrow\texttt{Bessel\allowbreak All}$.}
\label{claim:BNV21}
Under the analytic gap specified above,
\[
 \texttt{Bessel\allowbreak NZ}(\epsilon,\eta,u)
       \Longrightarrow\texttt{Bessel\allowbreak All}(\epsilon,\eta,u).
 \tag{BNV.37}
\]
All choices in the latter predicate, including every nonzero vector of
the lowest compact type, are retained.

\begin{proof}
Start with any nonzero witness (BNV.5).  We make each change of data
explicit.  All integrals in the argument converge by BNV19.

First, compact highest-weight uniqueness identifies any two choices
of $T$ unitarily, carrying their highest lines to each other.  The
normalized compact average defining $\Psi$ in DC2 is then conjugated
by that isometry: its rank-one highest-line projector and its
representation are conjugated, and its dimension and Haar normalization
are unchanged.  Also $\mu$ is forced to be $\eta+\rho$.  Thus neither
choice changes the existence of a nonzero period.

Second, for two Cayley matrices $P,P'$ write $P'=Pq$; the identity
$P^*J_\epsilon P=P'^*J_\epsilon P'=H$ gives $q\in G$.
Direct substitution in (BNV.5) gives
\[
 A_{P'}(v_1,v_2,\xi;x_0q,xq)
           =A_P(v_1,v_2,\xi;x_0,x).
 \tag{BNV.38}
\]
No invariance of the subgroup under $q$ is being claimed or needed;
only the two free group points have changed.

Third, use $a$ from BNV20 to change $z,c$, and its isometry $J$ to
change the compact representation.  By AMD8 and uniqueness of Haar
measure, the target Haar measures are positive multiples
$\alpha,\beta$ of the pushforwards of the old measures under
$h\mapsto aha^{-1}$ and $n'\mapsto an'a^{-1}$.
Changing variables gives the exact vector identity
\[
 A_2(v_1,v_2,J\xi;ax_0,ax)
       =\alpha\beta\,J A_1(v_1,v_2,\xi;x_0,x).
 \tag{BNV.39}
\]
The matrix argument reduces exactly:
\[
 (ax_0)^{-1}(aha^{-1})(an'a^{-1})(ax)=x_0^{-1}hn'x.
\]
The character is (BNV.36), and the representation factor satisfies
$\pi_2(aha^{-1})^{-1}J=J\pi_1(h)^{-1}$.
This also allows arbitrary Haar normalizations.  We have now obtained
a nonzero period for the entire target datum, still with possibly
different vectors $v_1,v_2$.

Finally fix any nonzero $v\in T$.  The finite-dimensional span of its
$K_\epsilon$ orbit is a nonzero invariant subspace, hence all of the
irreducible compact type $T$.  Express $v_1$ and $v_2$ as finite linear
combinations of $T(k_1)v$ and $T(k_2)v$.  Compact bi-equivariance gives
\[
 \begin{split}
 &\langle T(k_2)v,\Psi(Px_0^{-1}hn'xP^{-1})T(k_1)v\rangle\\
 &\qquad=\langle v,
 \Psi\bigl(P(x_0P^{-1}k_2P)^{-1}hn'(xP^{-1}k_1P)P^{-1}\bigr)v\rangle.
 \end{split}
\]
Consequently the known nonzero vector period is a finite linear
combination of periods with the same vector $v$ in both slots and
with new free group points.  At least one term is nonzero, proving
(BNV.37).  This proves the original universal transport by explicit
integrable changes of data; it introduces no derivative-decay or
uniqueness-of-Bessel-model premise.
\end{proof}
% END BNV21

% BEGIN BNV22
\paragraph{BNV22. Return to every original representation and its endomorphism integral.}
\label{claim:BNV22}
Suppose $\texttt{Bessel\allowbreak All}(\epsilon,\eta,u)$ holds in the analytic-gap
scope.  Then for every original $\sigma,w,\pi,z,c,\mu_W,\mu_N$ permitted
in (BNV.4), some $x_0,x\in G$ make the \emph{actual endomorphism-valued}
integral (BNV.4) nonzero.

\begin{proof}
First reconcile the independently existential unitary Cayley witnesses
in the original \texttt{Is\allowbreak DSHpk} and \texttt{Is\allowbreak Low\allowbreak Unit\allowbreak H} predicates.
If they are $P_0,P_1$, then $P_1P_0^{-1}$ is both Euclidean unitary
and $J_\epsilon$-unitary, hence belongs to $K_\epsilon$.
The resulting inner compact conjugacy preserves the entire lowest
$K_\epsilon$-isotypic subspace.  Thus the original $w$ belongs to the
lowest type for the same Cayley witness used for the representation.
The isomorphism theorem DC13 identifies the arbitrary supplied
$\texttt{Is\allowbreak DSHpk}$ representation with the constructed large-gap model,
after its allowed Cayley change.  It identifies its entire lowest
compact type with $T$.  The coefficient identity DC11, with the
embedding normalized isometrically, therefore gives a nonzero $v\in T$
such that
\[
 \langle w,\sigma(g)w\rangle
       =\langle v,\Psi(PgP^{-1})v\rangle\qquad(g\in G).
 \tag{BNV.40}
\]
Equivalently one may keep DC11's positive scalar; it multiplies all
following integrals by the same nonzero number.  No choice of the
original unit vector $w$ is being imposed.

By $\texttt{Bessel\allowbreak All}$ choose $\xi,x_0,x$ with
$A(v,v,\xi;x_0,x)\ne0$ for the original compact type, character, and
Haar pair.  Unitarity gives
\[
 \langle\sigma(x_0)w,\sigma(hn')\sigma(x)w\rangle
 =\langle w,\sigma(x_0^{-1}hn'x)w\rangle.
\]
Substitute (BNV.40).  BNV19 proves integrability of the operator
integrand, and evaluation $L\mapsto L\xi$ is a continuous linear map
on its finite-dimensional endomorphism space.  It commutes with the
Bochner integral, giving
\[
 \mathcal B_{\sigma,w}(x_0,x)\xi=A(v,v,\xi;x_0,x)\ne0.
\]
Thus the operator itself is nonzero.  This is precisely the original
conclusion, including its order of factors $hn'$, conjugated character,
two group translates, arbitrary lowest unit vector, and every original
Haar measure.  A nonzero abstract functional has not been substituted
for that conclusion.
\end{proof}
% END BNV22

\subsection{The remaining producer and the conditional original endpoint}

Here is the remaining source input, with its full scope.  The original
\texttt{Step\allowbreak Statement} says that, assuming the original
\texttt{Int\allowbreak SStatement}, for every $k,r$, every signature $\epsilon$ with
$p=k+r$, $q=r+1$, every $u$ and every real $h$ satisfying (BNV.12),
the predicate
\[
 \eta\longmapsto\texttt{Bessel\allowbreak NZ}(\epsilon,\eta,u)
 \tag{BNV.41}
\]
is step closed in exactly the sense (BNV.8).  In expanded form it has
two assertions: if $\eta$ and $\eta-e_s$ are Harish--Chandra and have
noncompact gaps at least $h$, if $u_j\ne\eta_s-1/2$ for every $j$, and
if the period for $\eta$ is nonzero, then the period for $\eta-e_s$ is
nonzero; and the analogous assertion for $\eta+e_s$ with exclusion
$u_j\ne\eta_s+1/2$.  The data witnessing the new existential period may
change.  This statement is a retained, unreconstructed source producer
in this section, not an external published theorem or a newly claimed
proof obligation closure.

The source proves the two directions in \texttt{Laplace\allowbreak Estimates/\allowbreak Step\allowbreak Down}
and \texttt{Step\allowbreak Up}.  Their substantive inputs concern standard and dual
tensor products, complementary Capelli minors, quotients with the shifted
Casimir and lowest weight, and nonzero period evaluation.  The factors
$\prod_j(u_j-\eta_s\pm1/2)$ explain the exclusions, but merely writing
these factors does not prove that the required quotient carries a
nonzero period.  That tensor-and-quotient argument is the remaining
analytic reconstruction task.  BNV1 removes its original subgroup
integrability antecedent; BNV18 removes its large-parameter base
antecedent; BNV6 supplies the entire finite path that will consume it.

% BEGIN BNV23
\paragraph{BNV23. Exact conditional closure of \texttt{Bessel\allowbreak NVGap} and the supplied-form setting.}
\label{claim:BNV23}
The original \texttt{Step\allowbreak Statement}, together with the inherited human
inputs named above, implies the whole original \texttt{Bessel\allowbreak NVGap}.
On the same supplied $\texttt{Forms0}$ scope as ASG, AMD, and BWP24, it
therefore implies the original \texttt{Concrete\allowbreak Setting\allowbreak Exists}, with
that supplied-form premise retained.

\begin{proof}
Take arbitrary data in the original target (BNV.4).  BNV7 supplies
$h$ satisfying the original tail gaps, $\texttt{Gap\allowbreak Lap}$, and
$\texttt{Step\allowbreak Gap}$.  BNV1 proves \texttt{Int\allowbreak SStatement}, so the retained
\texttt{Step\allowbreak Statement} makes (BNV.41) step closed.  By BNV18 choose
an integer $t\ge t_0$ for which
$\texttt{Bessel\allowbreak NZ}(\epsilon,\lambda^{(t)},u)$ holds.
BNV6 then gives $\texttt{Bessel\allowbreak NZ}(\epsilon,\lambda,u)$.
BNV21 gives $\texttt{Bessel\allowbreak All}(\epsilon,\lambda,u)$, and BNV22 gives
the nonzero endomorphism integral for the original arbitrary data.
This proves every quantifier of \texttt{Bessel\allowbreak NVGap} under the one
explicitly retained producer.

For the global setting, invoke the already proved BWP24.  Its input is
this very \texttt{Bessel\allowbreak NVGap}, not a branching substitute.  BWP24
constructs the local wave packet from the original coefficient and
nonvanishing integrals, transports it with the original frame and
$4\pi/g$ character, normalizes the global block before any later
\texttt{MBlk}, and joins ASG and AMD on the same data and exact measures.
It yields \texttt{Concrete\allowbreak Setting\allowbreak Exists} for each supplied block sequence
in their original \texttt{Forms0} scope.  No new choice of ambient forms,
no strengthening of that scope, and no orbital/fiber or Siegel--Weil
comparison follows from this invocation.
\end{proof}
% END BNV23

\paragraph{Resulting boundary.}
This section completely reconstructs the original subgroup integrability,
parameter walk, large-parameter Laplace base, universal-data transport,
and the return from a coefficient period to the prescribed endomorphism
integral.  The two one-unit propagation directions remain to be
reconstructed.  Thus (BNV.4), the supplied-form global setting, and all
subsequent orbital/fiber comparisons remain conditional at this stage.
No whole selected obligation, unrestricted RH theorem, or unrestricted
Dirichlet-GRH theorem is declared closed by this component.


\section{The analytic completion of a one-weight Bessel propagation}
\label{sec:bessel-weight-propagation-v50}

This section proves the original final reconstruction step
\texttt{Fin\allowbreak Statement}, the period-differentiation statement
\texttt{Per\allowbreak Statement}, and the generated-module Capelli eigenvalue statement
\texttt{Cap\allowbreak Val\allowbreak Statement}.  These are substantive producers used in both
\texttt{Step\allowbreak Down} and \texttt{Step\allowbreak Up}.  The full one-weight propagation is
still conditional on the complementary-minor evaluation and shifted
lowest-type quotient constructions.  We state that boundary explicitly
at the end.  None of the arguments below replaces the original integral
by an abstract branching assertion.

We retain BNV's notation $n=2r+k+1$, $G_H=U(H_{pk})$,
$G_\epsilon=U(J_\epsilon)$, $H_z$, $N$, and the actual character
$\psi_{z,c}$.  Put
\[
 \rho(g)=\|g\|_{\rm F},\qquad
 q_0=\min(p,q),\qquad
 a=\frac{h}{8\sqrt{2q_0}},\qquad \beta=a-1.
 \tag{BSP.1}
\]
The division in this notation is used only after positive rank has been
proved.  A Cayley matrix satisfies $P^*J_\epsilon P=H_{pk}$.
For a scalar function $F$ on matrices and a vector $\xi\in V_\pi$, the
original period is
\[
 \mathcal P_F(x)=
 \int_{H_z\times N}
 F\bigl(P(h'n'x)P^{-1}\bigr)\,
 \overline{\psi_{z,c}(n')}\,\pi(h')^{-1}\xi\,
 d\mu_W(h')d\mu_N(n').
 \tag{BSP.2}
\]
This is the source's \texttt{per}; the Haar measures retain their original
weak predicates, treated by AMD8.  Every statement concerning Bessel data
below allows every original $z,c,u,\pi,\mu_W,\mu_N,P$, with
$z$ a negative middle vector, all $c_a\ne0$, and the exact
$\texttt{Is\allowbreak HCType\allowbreak W}(k,r,z,u,\pi)$.

The inherited inputs are DC1--DC13, including the proved kernel Capelli
identity and degree \emph{orthogonality}, and BNV1.  No formula for the
formal degree, no Bessel multiplicity-one theorem, and no one-weight
propagation theorem is assumed.  Source evidence is pinned in the
accompanying manifest.

% BEGIN BSP1
\paragraph{BSP1. The whole original \texttt{Gap\allowbreak Need\allowbreak Statement}.}
\label{claim:BSP1}
For every $k,r\ge0$ and signature with $p=k+r$, $q=r+1$, the original
$\texttt{Step\allowbreak Gap}(\epsilon,r,h)$ implies
\[
 q_0\ge1,\qquad h>h_{\rm need}(\epsilon),\qquad
 \beta>4\max(2nq_0,2rn^2)>0.
 \tag{BSP.3}
\]
In particular it supplies DC's large-gap construction threshold and the
integrability inequalities $\beta/2>2nq_0$, $\beta/4>2rn^2$.

\begin{proof}
The two hypotheses are exactly
\[
 \max(3n/2,1+npq)<h-2,\qquad
 4\max(2nq_0,2rn^2)<h/(8\sqrt{2q_0})-1.
\]
If $q_0=0$, the second inequality, with the source's zero-division
convention, would say that a nonnegative number is less than $-1$.
Thus $q_0\ge1$.  Put $s=8\sqrt{2q_0}\ge8$ and
$X=\max(2nq_0,2rn^2)\ge2nq_0\ge2$.
The second hypothesis is $h>s(4X+1)$.
It supplies the third member of $h_{\rm need}$ in (BNV.11), since
$q-1=r$.  Moreover
\[
 s(4X+1)\ge64nq_0+s>16nq_0+s+2,
\]
which supplies the second member.  The first hypothesis supplies the
first member directly.  This proves $h>h_{\rm need}$ and (BSP.3).
Notice that $k=0$ was not removed from the source's scope: its impossible
rank-zero case is excluded by the hypotheses themselves.
\end{proof}
% END BSP1

% BEGIN BSP2
\paragraph{BSP2. The exact kernel datum and rigidity of its cyclic lowest type.}
\label{claim:BSP2}
Here the dimension $n$ and signature $\epsilon$ are arbitrary; the
signature restrictions of BSP1 are not imposed.
Let $\lambda$ be Harish--Chandra, with noncompact gaps at least $h$,
$q_0\ge1$, and $h>h_{\rm need}(\epsilon)$.
For every original compact type $T$ of weight
$\Lambda=\lambda+\rho_n-\rho_c$ and every $\mu=\lambda+\rho$, there
exist a complete Hilbert space $\mathcal H$, a unitary representation
$\sigma$, an invariant smooth derived domain $D$, matrix-unit operators
$E_{ab}$ on $D$, and a compact irreducible subspace $V\subset D$ of
highest weight $\Lambda$, together with an isomorphism $\iota:V\to T$,
such that
\[
 \begin{split}
 C_Gv&=(\|\lambda\|^2-\|\rho\|^2)v\quad(v\in V),\\
 \langle v',\sigma(g)v\rangle
   &=\langle\iota v',\Psi_{T,\mu}(g)\iota v\rangle
       \quad(g\in G_\epsilon,\ v,v'\in V).
 \end{split}
 \tag{BSP.4}
\]
The subspace $V$ is stable under the actual compact group, not merely its
Lie algebra.  This is the entire original \texttt{Psi\allowbreak Datum\allowbreak Statement}.

More generally, every smooth compact irreducible $V$ of this highest
weight and Casimir in a unitary representation generates a cyclic closed
Hilbert subspace unitarily equivalent to the constructed model.
Its lowest type corresponds to $T$, and its coefficients satisfy the
four-vector orthogonality relation with some $d>0$ for the same Haar
measure.  This latter assertion is used only at the displayed
large-gap scope.

\begin{proof}
For existence, take DC's constructed Hilbert model and its smooth-vector
domain.  Let $j:T\to\mathcal H$ be the lowest-type embedding of DC7.
Divide it by the square root of the positive scalar in DC11, so that
it is an isometry and its coefficients are exactly $\Psi_{T,\mu}$.
Let $V=j(T)$ and $\iota=j^{-1}$.  DC3--DC4 give the derived action and
Casimir, DC7 gives smoothness, and the compact equivariance gives actual
compact-group stability.  These are all fields of the original datum.

Here is the rigidity argument for the stated broader input; in
particular we do not assume that this input is already discrete series.
DC8 applies to any such unitary derived representation: the fixed
Casimir and gap kill the negative first noncompact tensor layer.
DC9 then puts every noninitial compact type in strictly higher
coercive PBW layers, and makes the original type occur once.
The energy estimate and convergent orbit Taylor series in DC10 apply to
the same finite initial compact type and scalar Casimir.  They show that
the closure of the generated Lie module is group invariant and retains
that unique compact type.  If a closed invariant subspace is projected
orthogonally onto, its projection commutes with the compact group.
Its restriction to the unique initial type is zero or identity by
Schur's lemma; cyclicity then makes the whole projection zero or
identity.  Thus the cyclic closure is irreducible.

Identify the initial type isometrically with the initial type of DC's
model and take its diagonal copy in the direct sum of the two unitary
representations.  It has the same weight and Casimir.  The preceding
PBW, closure, and irreducibility argument applies to this diagonal copy.
Its two coordinate projections are bounded nonzero intertwiners with
dense images onto the two cyclic closures.  Composing either projection
with its adjoint gives a positive scalar times the identity, by
irreducibility; the scalar is nonzero.  Hence each projection is a
nonzero scalar multiple of a unitary isomorphism.  The prescribed
isometry on the initial types makes the two scalars agree.  We obtain
a unitary equivalence identifying those initial types and intertwining
the actual group actions.
Finally DC12 transports its full coefficient orthogonality and positive
degree to this cyclic subspace.  This uses the proved Hilbert-space
orthogonality, not an unreconstructed degree polynomial.
\end{proof}
% END BSP2

% BEGIN BSP3
\paragraph{BSP3. The whole original \texttt{Cap\allowbreak Val\allowbreak Statement}.}
\label{claim:BSP3}
For arbitrary $n,\epsilon,\lambda,h,T,\mu$ satisfying BSP2's
Harish--Chandra, noncompact-gap, $h>h_{\rm need}$, $q_0\ge1$,
and weight hypotheses, let $(\mathcal H,\sigma,D,E,V,\iota)$ be
\emph{any} complete unitary derived datum with compact irreducible
type of weight $\lambda+\rho_n-\rho_c$, its stated Casimir,
actual compact-group stability, and the exact coefficient identity
(BSP.4).  The ambient representation need not be irreducible or be
the one constructed in BSP2.
Let $S=U(\mathfrak{gl}_n)V\subset D$ be the generated module, and put
$\theta_i=\lambda_i-(n-1)/2$.  Then for every $\omega\in\mathbb C$
and every $m\in S$,
\[
 \mathcal C_E(\omega)m
       =\prod_{i=0}^{n-1}(\omega+\theta_i)\,m.
 \tag{BSP.5}
\]
The Capelli operator uses the original ordered column determinant
\[
 \mathcal C_E(\omega)
 =\sum_{\tau\in S_n}\operatorname{sgn}(\tau)
   \prod_{j=0}^{n-1}
       \bigl(E_{\tau(j),j}+\delta_{\tau(j),j}(\omega-j)\bigr),
 \tag{BSP.6}
\]
with the rightmost operator acting first.  Thus in particular
$\mathcal C_E(-\theta_s)$ kills the entire generated module, as needed
by both one-weight steps.

\begin{proof}
The hypotheses include the exact coefficient identity (BSP.4);
$D$ is an invariant derived domain, so every finite derivative word
may be applied to its vectors.  Differentiating that identity on the
right and using DC4 gives, for $v,v'\in V$,
\[
 \langle v',\sigma(g)\mathcal C_E(\omega)v\rangle
   =p_\lambda(\omega)\langle v',\sigma(g)v\rangle,\qquad
 p_\lambda(\omega)=\prod_i(\omega+\theta_i).
 \tag{BSP.7}
\]
The shift is exact: substitute $Y=\omega-(n-1)/2$ in DC4's
$\mathcal C(Y+(n-1)/2)$ identity.

Set $u=\mathcal C_E(\omega)v-p_\lambda(\omega)v$.
Let $\mathcal K$ be the closed span of the group translates of $V$.
It contains every $E_{ab}m$ for $m\in D\cap\mathcal K$:
real infinitesimal actions are limits of difference quotients of group
orbits in the closed subspace, and complex matrix units are complex
linear combinations of real Lie algebra elements.  Induction shows
$S\subset\mathcal K$, so $u\in\mathcal K$.
Equation (BSP.7) and unitarity make $u$ orthogonal to every group
translate of $V$; hence $u=0$.  This proves (BSP.5) on $V$ without
assuming that the ambient representation is irreducible.

The ordered determinant (BSP.6) is central in the enveloping algebra.
This is the centrality calculation included in DC4: commuting a matrix
unit through the ordered column expansion pairs the off-diagonal
commutator terms by the corresponding transposition, with the
$(\omega-j)$ shifts canceling the diagonal commutators.  Consequently
$\mathcal C_E(\omega)$ commutes with every $E_{ab}$ in every
representation of the matrix-unit relations.  Therefore the kernel of
$\mathcal C_E(\omega)-p_\lambda(\omega)$ is invariant under all the
$E_{ab}$ and contains $V$.  By the definition of the generated module
it contains $S$, proving the full assertion.  Evaluation at
$\omega=-\theta_s$ makes the $s$th factor zero.
\end{proof}
% END BSP3

% BEGIN BSP4
\paragraph{BSP4. The original invariant standard and dual multipliers.}
\label{claim:BSP4}
For every middle vector $z$, without requiring negativity, set
$x_* = e_0$ if $r>0$, and $x_*=z$ if $r=0$.
For a column $y$ and a row $\alpha$ define
\[
 m_y(g)=x_*^*H_{pk}gy,\qquad
 m^\vee_\alpha(g)=\alpha g^{-1}x_*.
 \tag{BSP.8}
\]
These are exactly the original \texttt{c\allowbreak Mul} and \texttt{c\allowbreak Mul\allowbreak D}.
For every $s\in S_z=H_zN$ and $g\in G_H$ they satisfy
\[
 m_y(sg)=m_y(g),\qquad m^\vee_\alpha(sg)=m^\vee_\alpha(g).
 \tag{BSP.9}
\]
For the complexified right derivative $R_X$,
\[
 R_Xm_y=m_{Xy},\qquad
 R_Xm^\vee_\alpha=-m^\vee_{\alpha X}.
 \tag{BSP.10}
\]
For every $y,\alpha$ there is one constant $C$ such that both absolute
values in (BSP.8) are at most $C\sum_{i,j}|g_{ij}|$ on $G_H$.
These are all conclusions of the original \texttt{Mult\allowbreak Statement}.

\begin{proof}
If $r>0$, the defining isotropic flag makes every element of $N$ fix
its first vector $e_0$, and $H_z$ fixes the entire outer frame.
If $r=0$, $N=\{1\}$ and $H_z$ fixes $z$.  Thus $sx_*=x_*$ in either
case.  Since $s$ is $H_{pk}$-unitary, also
$x_*^*H_{pk}s=x_*^*H_{pk}$.  This proves both identities (BSP.9).
Differentiating $g\exp(tX)y$ gives $gXy$, while differentiating
$(g\exp(tX))^{-1}x_*=\exp(-tX)g^{-1}x_*$ gives
$-Xg^{-1}x_*$.  This proves (BSP.10) for real Lie directions and then
for their complexification; equivalently both expressions have local
holomorphic extensions on invertible matrices, so the same derivatives
are the source's ambient complex derivatives.
The first norm bound is a finite sum of fixed coefficients times entries
of $g$.  On $G_H$ one has $g^{-1}=H_{pk}g^*H_{pk}$, and $H_{pk}$ is
a signed permutation matrix.  Hence the sum of absolute entries of
$g^{-1}$ equals that of $g$, giving the second bound.  Take the larger
of the two constants.  The argument includes $z=0$ and all rank-zero
cases allowed in the original multiplier statement.
\end{proof}
% END BSP4

% BEGIN BSP5
\paragraph{BSP5. The whole original \texttt{Per\allowbreak Statement}.}
\label{claim:BSP5}
Assume the original signature identities and \texttt{Step\allowbreak Gap}.
If $F$ is continuous on $G_\epsilon$ and
$|F(g)|\le C\rho(g)^{-\beta}$ there, then (BSP.2) is absolutely
integrable for every original datum, every $x\in G_H$, and every $\xi$.
If $F$ is smooth in the source's local ambient sense and every
complexified right derivative word $E_{\mathbf w}F$ satisfies
\[
 |E_{\mathbf w}F(g)|\le C_{\mathbf w}\rho(g)^{-\beta},
 \tag{BSP.11}
\]
then differentiation commutes with the actual period:
\[
 \mathcal P_{E_{ab}F}(x)
    =\left.
       E_{ab}\bigl(y\longmapsto\mathcal P_F(P^{-1}yP)\bigr)
      \right|_{y=PxP^{-1}}.
 \tag{BSP.12}
\]
The operator on the right is the original \texttt{eab\allowbreak B}.

\begin{proof}
BSP1 gives $\beta>2rn^2$.  Conjugation by $P$ changes Frobenius norm
by fixed positive comparison constants.  BNV1 therefore makes
$\rho(P h'n'xP^{-1})^{-\beta}$ integrable.  The remaining character
and unitary factors in the period have norm $\|\xi\|$.
Continuity on the supported subgroups gives strong measurability,
and proves the first assertion.

For a real Lie direction $X$ in $G_\epsilon$, write
$x_t=xP^{-1}\exp(tX)P$.  The scalar integrand is then
$F(P h'n'xP^{-1}\exp(tX))$.
For $|t|\le t_0$ the inverse of $\exp(tX)$ has bounded Frobenius norm,
so uniformly in $h',n'$,
\[
 \rho(P h'n'xP^{-1}\exp(tX))
       \ge c\,\rho(P h'n'xP^{-1}),\qquad c>0.
\]
The derivative of the scalar integrand is
$(R_XF)(P h'n'xP^{-1}\exp(tX))$ and has the same integrable majorant,
by (BSP.11) and the finite expression of $X$ in the complex matrix units.
The fundamental theorem of calculus bounds its difference quotients
by that majorant.  Dominated convergence differentiates (BSP.2) under
the integral.  It also gives continuity of the resulting derivative.
Repeating this argument with every derivative word proves smoothness
in group coordinates and commutation with every real derivative.
Complex linearity yields (BSP.12) with exactly the source's matrix-unit
convention.

For completeness these conclusions have the local ambient meaning
required on the right side of (BSP.12).  Near any group point,
$G_\epsilon$ is a smooth embedded real submanifold of the invertible
matrices.  A local smooth retraction onto it extends any smooth function
on the group to an ambient smooth function.  Complexified tangential
derivatives of this extension are independent of the retraction.
Thus the calculated derivatives are the literal source derivatives,
not derivatives of an unrelated representative.  Haar support
eliminates all off-group values of the original functions.
\end{proof}
% END BSP5

% BEGIN BSP6
\paragraph{BSP6. The group power-integral threshold used by the final Fubini step.}
\label{claim:BSP6}
For any signature with $q_0\ge1$ and every $b>2nq_0$,
\[
             \int_{G_\epsilon}\rho(g)^{-b}\,dg<\infty.
 \tag{BSP.13}
\]
The Haar normalization is arbitrary.  On $G_\epsilon$ also
$\rho(g)\ge1$, $\rho(g^{-1})=\rho(g)$, and
$\rho(gh)\le\rho(g)\rho(h)$.

\begin{proof}
The last three assertions follow from Frobenius submultiplicativity,
$g^{-1}=J_\epsilon g^*J_\epsilon$, and $gg^{-1}=I$.
For the integral use the polar coordinates of DC6:
$g=\kappa e^X$ with $\kappa$ compact and
$X=X^*$, $J_\epsilon XJ_\epsilon=-X$.
Put $s=\|X\|_{\rm F}$.  Its nonzero eigenvalues are
$t_1,-t_1,\ldots,t_{q_0},-t_{q_0}$, allowing zeros, so
$\max_i|t_i|\ge s/\sqrt{2q_0}$ and
$\rho(e^X)\ge\exp(s/\sqrt{2q_0})$.
The polar Jacobian already calculated in DC6 is at most
$\exp(n\sqrt{2q_0}s)$.
Its remaining radial shell factor is polynomial in $s$, since the
noncompact tangent space is finite dimensional.  The resulting
integrand is bounded at infinity by a polynomial times
\[
 \exp\left[-\frac{b-2nq_0}{\sqrt{2q_0}}s\right],
\]
which is integrable.  Near zero the polar Jacobian is continuous.
The compact factor has finite Haar mass.  This proves (BSP.13),
including the exact threshold used in the source's
\texttt{fin\allowbreak B\_\allowbreak int\allowbreak J}.  A positive rescaling of Haar measure changes only
the finite constant.
\end{proof}
% END BSP6

% BEGIN BSP7
\paragraph{BSP7. Pointwise reproduction inside the actual right-regular space.}
\label{claim:BSP7}
Let $\mathcal K$ be a closed invariant subspace of $L^2(G_\epsilon)$
whose coefficients have the four-vector orthogonality identity with
degree $d>0$.  Suppose $F$ is continuous on $G_\epsilon$, is in $L^2$,
and its class $f$ belongs to $\mathcal K$.
For every $x\in G_\epsilon$,
\[
 \int_{G_\epsilon}F(y)
        \langle f,R(y^{-1}x)f\rangle\,dy
           =\frac{\|f\|_2^2}{d}F(x).
 \tag{BSP.14}
\]
This is an equality of continuous functions, not just of $L^2$ classes.

\begin{proof}
Haar measure on $G_\epsilon$ is invariant under right translation and
inversion.  Indeed the polar decomposition reduces the modulus of the
adjoint determinant to the compact factor, where it is one, and to
$\exp X$, where it is $\exp(\operatorname{tr}(\operatorname{ad}X))=1$:
the matrix-unit roots of the complexified Lie algebra occur in opposite
pairs, so the trace is zero.  Thus the group is unimodular.
Write $\gamma(g)=\langle f,R(g)f\rangle$.
Orthogonality makes $\gamma\in L^2$, so the integral in (BSP.14)
is absolutely convergent by Cauchy--Schwarz.  It is continuous in $x$
by strong continuity of translation in $L^2$.
Let $\varphi$ be any continuous compactly supported test function.
For fixed $x$, substitute $y=zt$ and apply Fubini:
\[
 \begin{split}
 &\int \overline{\varphi(z)}
       \left(\int F(y)\gamma(y^{-1}zx)\,dy\right)dz\\
 &\qquad=\int
       \langle\varphi,R(t)f\rangle\,
       \overline{\langle R(x)f,R(t)f\rangle}\,dt.
 \end{split}
 \tag{BSP.15}
\]
The double absolute integral is at most
$\|\varphi\|_1\|F\|_2\|\gamma\|_2$, which justifies this interchange.
If $\Pi$ is the orthogonal projection onto $\mathcal K$, replace
$\varphi$ by $\Pi\varphi$ in the first coefficient.  This is allowed
because all $R(t)f$ lie in $\mathcal K$.
Four-vector orthogonality makes the last line of (BSP.15)
\[
 \frac{\|f\|_2^2}{d}\langle\Pi\varphi,R(x)f\rangle
       =\frac{\|f\|_2^2}{d}
          \int\overline{\varphi(z)}F(zx)\,dz.
\]
Thus the difference of the two continuous functions in (BSP.14),
translated by $x$, pairs to zero with every compactly supported
continuous test function.  If it were nonzero somewhere, a sufficiently
small neighborhood and a nonnegative test function there would give
a nonzero pairing (multiply first by the conjugate phase at that
point).  Haar measure is positive on nonempty open sets.  This
contradiction proves the pointwise identity.
\end{proof}
% END BSP7

% BEGIN BSP8
\paragraph{BSP8. The full original \texttt{fin\allowbreak Proj} reconstruction.}
\label{claim:BSP8}
Assume the signature identities, \texttt{Step\allowbreak Gap}, and a Harish--Chandra
$\lambda'$ with noncompact gaps at least $h$.
Let $W_f$ be an original linear space of scalar matrix functions, and
$W'$ a compact irreducible subspace of the smooth vectors of
$L^2(G_\epsilon)$, with highest weight
$\Lambda'=\lambda'+\rho_n-\rho_c$ and Casimir $c_{\lambda'}$.
Assume actual compact-group stability, and the source's exact image
relation
\[
 v\in W'\ \Longleftrightarrow\
 \exists F\in W_f:\ v=[F]_{L^2}
 \quad\text{for every smooth }v.
 \tag{BSP.16}
\]
Every $F\in W_f$ is required to be continuous on the group and bounded
by $C_F\rho^{-\beta}$.
For every such $F$ whose class is a nonzero element of $W'$, there
exist an original compact type $T'$, $\mu'=\lambda'+\rho$, and
$v_1,v_2\in T'$ such that
\[
 F(x)=\int_{G_\epsilon}F(y)
      \langle v_2,\Psi_{T',\mu'}(y^{-1}x)v_1\rangle\,dy
       \qquad(x\in G_\epsilon).
 \tag{BSP.17}
\]

\begin{proof}
BSP1 supplies the large gap for BSP2.
The cyclic Hilbert space generated by $W'$ is therefore unitarily
equivalent to the constructed model, with the entire lowest type
matched.  In particular it has four-vector orthogonality for some
$d>0$, and there is $t\in T'$ with
\[
 \langle f,R(g)f\rangle
       =\langle t,\Psi_{T',\mu'}(g)t\rangle,\qquad f=[F]_{L^2}.
\]
Here the lowest-type embedding has been normalized isometrically,
as in BSP2.  The function $F$ is genuinely in $L^2$:
its displayed bound, BSP1, and BSP6 show this directly.
Consequently the formal \texttt{to\allowbreak L2} in (BSP.16) is its usual class,
and $f\ne0$.
Apply BSP7 and take
$v_2=t$, $v_1=(d/\|f\|_2^2)t$.
The scalar is nonzero and acts in the linear inner-product slot.
The result is (BSP.17), with no almost-everywhere ambiguity and with
the exact kernel required by the original \texttt{fin\allowbreak Proj}.
\end{proof}
% END BSP8

% BEGIN BSP9
\paragraph{BSP9. The full original final Fubini and nonvanishing extraction.}
\label{claim:BSP9}
Under the original signature, gap, Bessel-data, and kernel hypotheses,
suppose $F$ is continuous on $G_\epsilon$ and
$|F(y)|\le C\rho(y)^{-\beta}$.
For every $x\in G_H$ the following double integrand is Bochner integrable:
\[
 F(y)\,
 \langle v_2,\Psi_{T',\mu'}
        (y^{-1}P h'n'xP^{-1})v_1\rangle\,
 \overline{\psi(n')}\,\pi(h')^{-1}\xi.
 \tag{BSP.18}
\]
The measures are $dy\,d\mu_W(h')\,d\mu_N(n')$.
If in addition $F$ satisfies (BSP.17) and
$\mathcal P_F(x)\ne0$, then some $x_0\in G_H$ satisfies
\[
 A_{T',\mu',P}(v_1,v_2,\xi;x_0,x)\ne0
 \tag{BSP.19}
\]
for the original integral (BNV.5), with the same Bessel datum.
These are the original \texttt{fin\allowbreak Fub\_\allowbreak int'} and \texttt{fin\allowbreak Fub}
conclusions; BNV1 proves their \texttt{Int\allowbreak SStatement} input.

\begin{proof}
Put $g=P h'n'xP^{-1}\in G_\epsilon$ on the Haar support.
The Frobenius coefficient estimate proved in BNV19, or its proof
from DC6, gives
$|\langle v_2,\Psi(y^{-1}g)v_1\rangle|
       \le C'\rho(y^{-1}g)^{-a}$.
Since $\rho\ge1$ and $a=\beta+1$, this is at most
$C'\rho(y^{-1}g)^{-\beta}$.
Let $A=\rho(y)$, $B=\rho(y^{-1}g)$, and $G_0=\rho(g)$.
Then $A,B\ge1$ and $G_0\le AB$.
The following elementary majorant suffices:
\[
 A^{-\beta}B^{-\beta}
   \le G_0^{-\beta/4}A^{-\beta/2}.
 \tag{BSP.20}
\]
Indeed $G_0^{-\beta/4}\ge A^{-\beta/4}B^{-\beta/4}$,
and $A^{-\beta}B^{-\beta}\le
A^{-3\beta/4}B^{-\beta/4}$ because $A,B\ge1$.
This slightly simpler majorant replaces the source's sum of two
majorants without changing any required exponent.

The first factor on the right of (BSP.20) is integrable over
$H_z\times N$ by BNV1 and norm comparison under $P$:
$\beta/4>2rn^2$ by BSP1.
The second is integrable over $G_\epsilon$ by BSP6, since
$\beta/2>2nq_0$.  Their product is therefore integrable.
The character and representation factors have norm $\|\xi\|$;
the supported functions are continuous and strongly measurable.
This proves the full Bochner-integrability assertion (BSP.18).

Insert (BSP.17) into (BSP.2) and interchange the integrals using
that assertion.  With $x_0=P^{-1}yP$ one has exactly
\[
 y^{-1}P h'n'xP^{-1}
       =P(x_0^{-1}h'n'x)P^{-1}.
\]
Thus the inner integral after the interchange is
$A(v_1,v_2,\xi;P^{-1}yP,x)$.
If all these vector integrals were zero, the resulting expression
for $\mathcal P_F(x)$ would be zero, a contradiction.
Since $P^{-1}yP\in G_H$, this proves (BSP.19) with the original
two free translates and all original normalization factors.
\end{proof}
% END BSP9

% BEGIN BSP10
\paragraph{BSP10. The whole original \texttt{Fin\allowbreak Statement}.}
\label{claim:BSP10}
Take every datum and hypothesis in BSP8, together with every original
Bessel datum and a fixed $\xi\in V_\pi$.  Assume the source's
period-injectivity condition
\[
 \begin{gathered}
 F\in W_f,\quad
 \mathcal P_F(x)=0\ \text{for every }x\in G_H\\
 \Longrightarrow\quad F(g)=0\ \text{for every }g\in G_\epsilon.
 \end{gathered}
 \tag{BSP.21}
\]
Then $\texttt{Bessel\allowbreak NZ}(\epsilon,\lambda',u)$ holds.
This proves the full original \texttt{Fin\allowbreak Statement}, including its
$\texttt{Int\allowbreak SStatement}$ antecedent by BNV1, every original space
$W_f,W'$, every original Haar pair, and the exact period (BSP.2).

\begin{proof}
The nonzero irreducible compact type $W'$ contains a nonzero vector
$f$.  The exact image relation (BSP.16) supplies $F\in W_f$ representing
it.  If all periods of $F$ vanished, (BSP.21) would make its function
zero on the group and hence its $L^2$ class zero, contradicting $f\ne0$.
Choose $x$ with $\mathcal P_F(x)\ne0$.
BSP8 supplies $T',\mu',v_1,v_2$ and the pointwise reproduction formula;
BSP9 supplies $x_0$ with the nonzero original vector integral.
These, together with the unchanged $P,z,c,\pi,\mu_W,\mu_N,\xi,x$,
are every field in the existential definition of
$\texttt{Bessel\allowbreak NZ}(\epsilon,\lambda',u)$.
The construction has not assumed existence of such a period or of
a shifted discrete series beforehand: nonvanishing came from the
explicit injectivity hypothesis (BSP.21), and the shifted model came
from its stated lowest-type and Casimir data.
\end{proof}
% END BSP10

\subsection{What remains in each one-unit propagation}

The new complete producers are the original
\texttt{Gap\allowbreak Need\allowbreak Statement}, \texttt{Psi\allowbreak Datum\allowbreak Statement},
\texttt{Cap\allowbreak Val\allowbreak Statement}, \texttt{Mult\allowbreak Statement},
\texttt{Per\allowbreak Statement}, \texttt{Fin\allowbreak Statement}, and the intermediate
pointwise reconstruction and final Fubini theorems used to prove the last
one.  The source's two complete step proofs still require the following
substantive inputs.

First, the standard and dual tensor coefficient maps must be shown to be
smooth $L^2$ functions with the exact derivative bounds.  Their formulas
are fixed: with $x_1=P x_*$,
\[
 \begin{split}
 T_f(g)&=\sum_j (x_1^*J_\epsilon g e_j)
                  \langle w,\sigma(x_0^{-1}g)f_j\rangle,\\
 T_f^\vee(g)&=\sum_j (e_j^*g^{-1}x_1)
                  \langle w,\sigma(x_0^{-1}g)f_j\rangle.
 \end{split}
 \tag{BSP.22}
\]
Here $w$ lies in the original lowest compact type and every $f_j$ in
its generated module.  BSP4 verifies the exact scalar multipliers and
their derivatives; it does not by itself prove all derivative bounds
for these generated-module coefficients.

Second, the complementary Capelli minors must produce the relevant
Euler-operator eigenspaces in the standard and dual tensors.
BSP3 proves the required annihilation
$\mathcal C_E(-\theta_s)S=0$, but the noncommutative adjugate identities
turning it into those eigenspace vectors remain separate.
The original \texttt{Eval\allowbreak PStatement} and \texttt{Eval\allowbreak MStatement} then
assert a nonzero period of a specific such tensor image whenever
the old period is nonzero and
$u_j\ne\lambda_s+1/2$ or $u_j\ne\lambda_s-1/2$, respectively.
Their products of excluded factors alone are not a proof of this
evaluation assertion.

Third, \texttt{Quot\allowbreak PStatement} and \texttt{Quot\allowbreak MStatement} must produce
spaces $W_f,W'$ for the shifted weights $\lambda+e_s$ and $\lambda-e_s$,
with exactly (BSP.16), the shifted Casimir, actual compact stability,
the smooth decay properties, and the period-injectivity (BSP.21).
BSP5 supplies differentiation through their actual periods and BSP10
completes the last step once these spaces are constructed.  It does
not assert the quotient existence.

Consequently the whole original \texttt{Step\allowbreak Statement} remains open.
BNV23, the whole \texttt{Bessel\allowbreak NVGap}, and the supplied-form global setting
remain conditional on that source producer.  No selected global
obligation is closed by this section.


\section{The actual denominator table and its uniform global export}
\label{sec:denominator-export-v050}

\status{DNE1--DNE14 reconstruct the full implications
\texttt{Concl\allowbreak F}$\Rightarrow$\texttt{Adm\allowbreak F}$\Rightarrow$
\texttt{Link\allowbreak N}, including the actual local table, its coefficient-place
weights and the global prime sum.  The original common arithmetic
supplier of \texttt{Concl\allowbreak F} is retained.  Consequently
\texttt{GLOBAL\_\allowbreak DENOMINATOR} remains open.  A deliberate simplification
replaces the source's particular pointwise \texttt{Large} cutoff by a
uniform eventual cutoff; the original \texttt{Link\allowbreak N} conclusion and
its constant $2000(|C|+1)^3$ are unchanged.}

The notation in this chapter always denotes the objects of one original
$D:\texttt{D7}$, $L:\texttt{Size\allowbreak Data}_2(D)$ and
$P:\texttt{Parameter\allowbreak Choice}'(D,L)$.  No free module, table on a different
host, or replacement coefficient field is introduced.  Put
\[
 x=L.\mathrm{sc}.x,\quad \epsilon=L.\mathrm{sc}.\epsilon,
 \quad z=x^\epsilon,\quad b=e^z,\quad X'=e^{z^{1/3}},
 \quad b'=xb/X',\quad w''=b/(nX').                         \tag{DNE.1}
\]
Here $n=D.R.n$, $f=2n$, $m=D.R.m$, $k=D.R.k$,
$r=k+m-1$, $t=D.R.t$, $N=k+m(t+1)$ and $n_V=r+t$.
We have $m\ge2$, $m\le f\le k$, $r\le2k$, $t\ge r+5$,
and $tr$ is odd.  Write $S=L.\mathrm{Sc}$, $\eta=L.\mathrm{etaSz}$,
$\delta=L.\mathrm{de2}$ and
\[
\begin{split}
 \mathcal D&=nt^2(N+t)z^2+kn^2(S+t+z),\\
 p_*&=mb/\mathcal D,\qquad A=kz/(m\delta),\qquad R_0=n/z,\\
 H&=b+S+\eta+t+2,\qquad \eta_3=\eta+t+2,\\
 w_B&=\max\{w''+\eta,\ \sup_{a\in D.\mathrm{host}.TT}b_a+b'/x+\eta\},\\
 w_P&=w_B+(t+1)/2,\qquad G_P=w_P+S+\eta_3,\\
 e_p&=C\left(1+H/p+\frac{\log H}{\log p}\right),\qquad
 j_p=\max\{1,\lceil\log_{p-1}A\rceil\}.
\end{split}                                                       \tag{DNE.2}
\]
The source names are respectively \texttt{Den}, \texttt{l\allowbreak Star},
\texttt{l\allowbreak A}, \texttt{R0}, \texttt{h\allowbreak B}, \texttt{eta3},
\texttt{w\allowbreak B}, \texttt{w\allowbreak P}, \texttt{GP}, \texttt{e\allowbreak P} and \texttt{j\allowbreak A}.
The empty supremum in $w_B$ has the source's zero convention.
The host field is $E$, with $[E:\mathbb Q]=2f$; the coefficient
field of the rows is denoted $K$, and need not be $E$.
The two generator primes form $G=\{g_0,g_1\}$, the small primes
are the divisors of $Q_S$, and $q,d$ are the two distinct ramified
auxiliary primes.  The host embeds in
$\mathbb Q(\zeta_M)$ for $M=Q_Sqd$, and $4\mid Q_S$.

\subsection{Exact local data and weights}

% BEGIN DNE1
\paragraph{DNE1: place weights and the unchanged zero-row rank.}
\label{claim:DNE1}
For a coefficient place $w$ above a rational prime $p$, the actual
weight in $\eta\mathrm{Sum}$ is
\[
 a_w=\frac{\log \mathrm Nw}{[K:\mathbb Q]}e_w
     =\frac{e_w f_w}{[K:\mathbb Q]}\log p.                 \tag{DNE.3}
\]
Thus $a_w\ge0$, and the sum over any subset of places above $p$
is at most $\log p$.  Likewise, for the host places $\lambda\mid p$,
put $n_\lambda=e_\lambda f_\lambda$.  Then
\[
 \sum_{\lambda\mid p}n_\lambda=2f,\qquad
 \sum_{\lambda\mid p}(n_\lambda+1)\le4f.                 \tag{DNE.4}
\]
These are the number-field degree formula; every residue degree
and ramification index is positive.  In the source's finite-set
presentation one may use the corresponding inequalities, which
are sufficient below.

Fix a lift $j\in\{1,\ldots,m\}$ and write $d_j=L.dZ0(j)$.
The original \texttt{adm1}, \texttt{adm2}, \texttt{adm3} say,
away from $G$, that the local row loss is bounded by
\[
 \eta_{\mathrm{loc},i}\le
    \sum_{\pi\in Z_i^{\mathrm{good}}}
         m_{i,\pi}(v_{i,\pi}+v_p(J)),\qquad
 \sum_{i\in\operatorname{locs}(j,w)}
    \sum_{\pi\in Z_i^{\mathrm{good}}}m_{i,\pi}
 \le \operatorname{rk}_0(j,w)=d_j,                       \tag{DNE.5}
\]
where $J=D.\mathrm{what}_2.\mathrm{restRowsFiniteGroupCard}>0$.
The actual equality
$\eta_J(j,w)=\sum_{i\in\operatorname{locs}(j,w)}\eta_{\mathrm{loc},i}$
is part of the $D$ record.  Consequently, if every summand at $p$
satisfies $v_{i,\pi}+v_p(J)\le U_p$, with $U_p\ge0$, then
\[
 \eta\mathrm{At}(L,j,p)\le d_j\log p\,U_p.              \tag{DNE.6}
\]
Indeed multiply (DNE.5) by the nonnegative $U_p$, use the exact rank
identity, multiply by $a_w$, and sum using (DNE.3).  This also
handles an empty place set or $d_j=0$; no division by $d_j$ occurs.
% END DNE1

% BEGIN DNE2
\paragraph{DNE2: an elementary sufficient cyclotomic budget.}
\label{claim:DNE2}
For the actual cyclotomic character $\chi_p:G_E\to\mathbb Z_p^\times$
and the odd positive integer $a=tr$, set
\[
 \nu_p=\min_{\sigma\in G_E}v_p(\chi_p(\sigma)^a-1).
\]
This minimum is the source's \texttt{nu\allowbreak Of}, hence its real image is
\texttt{nu\allowbreak L}.  For every prime $p$,
\[
 \nu_p\le v_p(M)+v_p(a),\qquad
 \sum_{p\in T}\nu_p\log p\le\log M+\log a              \tag{DNE.7}
\]
for every finite prime set $T$.

\begin{proof}
An automorphism of a cyclotomic extension whose exponent is $1$
modulo $M$ fixes $E$, because $E$ is embedded in
$\mathbb Q(\zeta_M)$.  For any prescribed finite $p$-power precision,
the Chinese remainder theorem permits every $p$-adic unit congruent
to $1$ modulo $p^{v_p(M)}$ as such an exponent.  Extend the resulting
automorphism of the finite cyclotomic compositum to the algebraic
closure over $E$.

If $p$ is odd and $p\nmid M$, choose the unit $-1$; as $a$ is odd,
$(-1)^a-1=-2$ is a unit.  If $p\mid M$, put $h=v_p(M)$ and choose
$1+p^h$.  The binomial theorem gives
\[
 v_p((1+p^h)^a-1)=h+v_p(a)
\]
for odd $p$.  One direct proof first treats an exponent prime to
$p$, for which the linear term has strictly smaller valuation than
the others, and then repeatedly raises to the $p$th power; each
step increases the valuation by one.  For $p=2$ we have $h\ge2$
and $a$ odd, so the linear term again proves the same equality.
Use precision strictly greater than this displayed valuation.
Approximation of the unit by an actual cyclotomic character value
preserves it.  This proves the first inequality and also proves
that the minimum is finite.  The second follows from
$\sum_pv_p(u)\log p=\log u$ for positive integers $u$.
\end{proof}

The source proves the sharper budget $\log Q_S+\log(tr)$ by removing
$q,d$.  The sufficient budget (DNE.7) is used here only as an
intermediate estimate: $\log M\le z/72$ follows from the original
parameter inequality $M^{36}t^{150}\le b^{1/2}$.  The final
\texttt{Link\allowbreak N} bound below is exactly the original one.
% END DNE2

% BEGIN DNE3
\paragraph{DNE3: patching the actual local bounds, including the prime three.}
\label{claim:DNE3}
Suppose the original \texttt{Concl}(L,P,X,C) is supplied, where
$X$ is its actual \texttt{Hub\allowbreak Data}.  Its representations are denoted
$\rho_i$; its finite good-point sets, exceptional prime set and band
predicates are $Z_i^{\mathrm{good}}$, $\mathcal E$ and $B_{i,\lambda}$.
Put
\[
 L^X_{i,\lambda}=\max_{\pi\in Z_i^{\mathrm{good}}}
                       X.\mathrm{len0}_{i,\pi,\lambda},\qquad
 B_i^X=\max_{\pi\in Z_i^{\mathrm{good}}}
                       X.\mathrm{errorBoundHub}_{i,\pi},           \tag{DNE.8}
\]
with maximum zero for an empty set.  The same representation gives
$s_{i,\pi}=\operatorname{sPi}(\rho_i,\pi)$.
Take the one table $T_f$ provided by the \texttt{Star\allowbreak F} field, and set
$b_i^f=\sum_{\lambda\mid p_i}(T_f.\mathrm{len})_{i,\lambda}
             +T_f.\mathrm{BB}_i$.
Define the ordinary range
\[
 \mathcal R=\{p:p\notin\mathrm{small},\ p\ne3,q,d,
                         \ p\le p_*\ \hbox{or}\ p\in\mathcal E\}.
\]
The literal \texttt{tab} of the source is the following table,
with the branches taken in the displayed order:
\[
 \ell_{i,\lambda}=
 \begin{cases}
  0,&p_i\in\mathrm{small},\\
  b_i^f/\#\{\lambda\mid3\},&p_i=3\notin\mathrm{small},\\
  L^X_{i,\lambda},&p_i\in\mathcal R,\\
  0,&\text{otherwise},
 \end{cases}\qquad
 B_i=
 \begin{cases}
  b_i^f,&p_i\in\mathrm{small},\\
  \operatorname{midG}(L,C,p_i),&p_i=q\text{ or }d,\\
  B_i^X,&p_i\in\mathcal R,\\
  0,&\text{otherwise}.
 \end{cases}                                                   \tag{DNE.9}
\]
It retains $s$, $B_{i,\lambda}$ and $\mathcal E$, and sets
$\mathrm{gcost}(j,p)=Cx^C\eta$, $\mathrm{gs}(j,p)=p^{-2}$.
The division in (DNE.9) is legitimate: the number field $E$ has
finitely many and at least one prime above $3$.  For example $3$
is a nonunit in $\mathcal O_E$, since its norm is
$3^{[E:\mathbb Q]}$, so a maximal ideal contains it.
Thus summing its equal shares gives exactly $b_i^f$.

For every twisted local row away from $G$, (DNE.9) satisfies
\[
 v_{i,\pi}\le s_{i,\pi}+\sum_{\lambda\mid p_i}\ell_{i,\lambda}
       +V_{p_i}+v_{p_i}(h_E)+B_i+p_i^{-2}.               \tag{DNE.10}
\]
Here $V_p=D.\mathrm{what}_2.VV(p)$ and $h_E$ is the actual host
class number.  At a small prime or at $3$, this is exactly
\texttt{Star\allowbreak F}, by the equal-share identity.  At $q,d$, it is
\texttt{Star\allowbreak G}.  In $\mathcal R$, dominate each local length and
error by its maximum (DNE.8) and use \texttt{Star\allowbreak AD}.
In the remaining case $p>p_*$ and $p\notin\mathcal E$, so
\texttt{Star\allowbreak E} gives (DNE.10) with zero added length and error.
These alternatives exhaust the original prime cases.
% END DNE3

\subsection{All admissibility fields on this same table}

For precision, the numerical part of \texttt{Adm\allowbreak T} is recorded here.
The table has error $B_i\le B_0=Cfnt^2n_Vz^2$ away from
$\mathrm{small}\cup\{q,d\}\cup G$, and
$s_{i,\pi}+p_i^{-2}\le\nu_{p_i}+2p_i^{-2}$.
For the split and inert rows, put
\[
 F_{i,\lambda}(p)=\frac{w_P+\eta_3}{p}+S+z+
                          \eta_3\mathbf1_{B_{i,\lambda}},
 \quad R(p)=\frac H{(p-1)^{j_p+1}}+(G_P+1)(j_p+1)+z.     \tag{DNE.11}
\]
The four bounds, at the source's twisted rows away from $G$, are
\[
\begin{array}{ll}
\text{split},\ p\ge5,\ p\le p_*\text{ or }p\in\mathcal E:
 &2\sum_{\lambda\mid p}\ell_{i,\lambda}
       \le\sum_{\lambda\mid p}(2n_\lambda m+1)e_p,\\[2pt]
\text{inert},\ p<A,\ \operatorname{ord}_{\pm}(p)>4L_w:
 &\ell_{i,\lambda}\le C(n_\lambda+1)rR(p),\\[2pt]
\text{inert},\ p\ge A,\ p\le p_*\text{ or }p\in\mathcal E:
 &\ell_{i,\lambda}\le C(n_\lambda+1)r(H/p^2+F_{i,\lambda}(p))
                              +Cn_\lambda m e_p,\\[2pt]
\text{inert},\ p<A,\ \operatorname{ord}_{\pm}(p)\le4L_w:
 &\ell_{i,\lambda}\le C(n_\lambda+1)rH/p^2
       +C(n_\lambda+1)rF_{i,\lambda}(p)+Cn_\lambda m e_p.
\end{array}                                                     \tag{DNE.12}
\]
The order is in
$(\mathbb Z/2^{e+2}\mathbb Z)^\times/\{\pm1\}$.
The band implication is
\[
 B_{i,\lambda},\ \lambda\mid p\quad\Longrightarrow\quad
 \exists s\in[w'',w_B]:\ |ps-b|
   \le C\{p\eta_3+pS+b/p+s+pz\}.                       \tag{DNE.13}
\]
Above $p_*$ outside $G\cup\mathcal E$, both $\ell$ and $B_i$ vanish.
The exception bounds are
\[
 |\mathcal E|\le(|TT|+2)(S+3\eta+t+8b/p_*+9),\qquad
                    p\in\mathcal E\Longrightarrow p\le3b.       \tag{DNE.14}
\]
For $p\in\mathrm{small}\cup\{3\}$ the sum of the lengths and
error is at most $Cfm b\log x+Cx^Cz^2$; for $p=q,d$ it is at most
$\operatorname{midG}(L,C,p)\le2Cfm b$.
For $p\in G$ the direct bound is
\[
 \eta\mathrm{At}(L,j,p)\le
    d_j\log p\{Cx^C\eta+2p^{-2}+V_p+v_p(h_E)\}.        \tag{DNE.15}
\]
Finally \texttt{Parameters\allowbreak Polynomially\allowbreak Bounded} means exactly
\[
 \begin{gathered}
  P.CLw\le C,\quad P.CYL\le C,\quad k,n,t\le x^C,
  \quad x^{-\epsilon}\le\delta,\\
  \log h_E\le Cx^C,\qquad \log J\le Cx^Cz .
 \end{gathered}                                                   \tag{DNE.16}
\]
In particular a supplied admissible table forces $C>0$, since
$P.CLw>0$.

% BEGIN DNE4
\paragraph{DNE4: the full original table is admissible.}
\label{claim:DNE4}
The table (DNE.9) satisfies all seventeen fields of
\texttt{Adm\allowbreak T}$(L,\mathrm{tab},C,P)$.

\begin{proof}
The first three fields are \texttt{Concl.\allowbreak pass},
\texttt{Concl.\allowbreak mult} restricted to places outside $G$, and
\texttt{Concl.\allowbreak rank}.  DNE3 proves the fourth field (DNE.10).
For the fifth, at the ordinary range take the maximum in
\texttt{BBLe}; outside it $B_i=0$, including $p=3$ when it is not
small.  The required upper bound is nonnegative.  For the sixth,
\texttt{Kill\allowbreak At} gives $s_{i,\pi}\le\nu_p+p^{-2}$, and
$\mathrm{gs}+p^{-2}=2p^{-2}$.

For the four rows (DNE.12), in $\mathcal R$ the new length is
exactly the original maximum length supplied by \texttt{Row\allowbreak Num}.
Outside $\mathcal R$ it is zero.  The split row excludes $3$
by $p\ge5$; the inert rows exclude it by the original
\texttt{Concl.\allowbreak small} clause that $3$ is not inert.  Every
right-hand side in (DNE.12) is nonnegative, since $C\ge0$,
$H\ge1$, all ranks are nonnegative, and $p\ge2$.
Thus the zero alternative also satisfies all four bounds.
The band predicate was not changed, so (DNE.13) is unchanged.

The \texttt{Small} field puts $3,q,d$ and every small prime below
$p_*$.  Hence a prime above $p_*$ outside $\mathcal E\cup G$ lies
in the zero branch of (DNE.9).  This proves large-prime vanishing;
the exceptional set itself is unchanged, giving (DNE.14).
The small-prime bound is preserved by putting its whole cost in
$B_i$, or, at $3$ outside the small set, distributing it equally
over the primes above $3$.  At $q,d$, both primes are outside the
small set and unequal to $3$, and the length is zero while the
error is exactly $\operatorname{midG}$.  The retained
\texttt{Small} bound gives its advertised upper bound.
Finally (DNE.15) is \texttt{Star\allowbreak H} with the chosen
$\mathrm{gcost},\mathrm{gs}$, and (DNE.16) is
\texttt{Concl.\allowbreak par}.  This accounts for each field on the actual
table and actual host.
\end{proof}

As a separate elementary constructor, the source's
\texttt{Bound\allowbreak At\allowbreak Two\allowbreak Three} and \texttt{Star\allowbreak S} imply \texttt{Star\allowbreak F}:
take all length entries zero and $\mathrm{BB}_i=Cfm b\log x+Cx^Cz^2$.
At $2$ use the former hypothesis; at every other small prime use
the latter; at $3$ outside the small set use the former again.
This proves the source's full \texttt{star\allowbreak F\_\allowbreak of} implication
without dropping the $2$ or $3$ case.
% END DNE4

% BEGIN DNE5
\paragraph{DNE5: a scalar majorant and the original global row sum.}
\label{claim:DNE5}
For any admissible table define $I(p)\in\{0,1\}$ by the existence
of the $s$ in (DNE.13), and put
\[
 F(p)=(w_P+\eta_3)/p+S+z+\eta_3I(p).
\]
Define
\[
 \mathcal L(p)=
 \begin{cases}
 (2fm+f)e_p,&p\text{ split},\\
 4fCrR(p),&p<A,\ \operatorname{ord}_{\pm}(p)>4L_w,\\
 4fCr(H/p^2+F(p))+2fCm e_p,&\text{otherwise},
 \end{cases}                                                     \tag{DNE.17}
\]
and the ordered scalar table
\[
 \Lambda(p)=
 \begin{cases}
 Cx^C\eta,&p\in G,\\
 Cfm b\log x+Cx^Cz^2,&p\in\mathrm{small}\text{ or }p=3,\\
 2Cfm b,&p=q\text{ or }d,\\
 0,&p>p_*,\ p\notin\mathcal E,\\
 B_0+\mathcal L(p),&\text{otherwise}.
 \end{cases}                                                     \tag{DNE.18}
\]
These are the literal source \texttt{len\allowbreak B}, \texttt{lam\allowbreak G}; in
particular the occurrence of $C$ inside $e_p$ is retained.

For every lift $j$, assuming only summability of the original
$\mathrm{vvTerm}(D,p)$, one has
\[
\begin{split}
 \eta\mathrm{Sum}(L,j)\le d_j\biggl(&6+\log M+\log(tr)
                      +\log h_E+\log J\\
             &+\sum_{p\in T}\log p\,\Lambda(p)
                         +\sum_p\mathrm{vvTerm}(D,p)\biggr),
\end{split}                                                       \tag{DNE.19}
\]
where $T$ is exactly the image under residue characteristic of the
twisted places in $L.\mathrm{supp}$.

\begin{proof}
Sum (DNE.12) over $\lambda$, using (DNE.4).
In the split case, division by $2$ gives $(2fm+f)e_p$.
In the other cases, the bounds $\sum(n_\lambda+1)\le4f$
and $\sum n_\lambda\le2f$ give (DNE.17); the band implication
ensures $F_{i,\lambda}(p)\le F(p)$.  Together with the remaining
admissibility fields this proves
$\sum_{\lambda\mid p}\ell_{i,\lambda}+B_i\le\Lambda(p)$
outside $G$.

Use (DNE.10), the sixth field of admissibility, and DNE1, with
\[
 U_p=\nu_p+2p^{-2}+\Lambda(p)+V_p+v_p(h_E)+v_p(J).
\]
Every term here is nonnegative.  At $G$, (DNE.15) gives the same
bound directly, with harmless nonnegative extra terms.  Grouping
the original finite place sum by residue characteristic is an exact
finite-sum identity.  The cyclotomic sum is bounded by DNE2;
$\sum_{p\in T}2\log p/p^2\le4<6$ follows by comparison with the
integral of $\log u/u^2$ on $[2,\infty)$ and its first term.
Unique factorization gives
$\sum_{p\in T}v_p(h_E)\log p\le\log h_E$ and the same statement
for $J$.  Finally
\[
 \log p\,V_p=
  \log p\sum_{a=1}^{t-1}\mathrm{shaLen}_D(p,a)
       =\mathrm{vvTerm}(D,p).
\]
These are nonnegative terms.  Summability therefore permits bounding
their finite sum by the full series on this very host.  This proves
(DNE.19), retaining the exact $d_j$, $\mathrm{shaLen}$ and $\mathrm{vvTerm}$.
\end{proof}
% END DNE5

\subsection{Uniform estimates with a simpler eventual cutoff}

% BEGIN DNE6
\paragraph{DNE6: the uniform scale comparison.}
\label{claim:DNE6}
Fix $0<\epsilon<1/52$ and $C>0$.  All assertions below are uniform
over every original $D,L,P$ satisfying (DNE.16) with these fixed
$\epsilon,C$, once $x\ge x_0(\epsilon,C)$.
The original parameters give
\[
\begin{gathered}
 \eta\le e^{z^{1/4}},\quad S=4k\Delta\le4x^C e^{z^{1/4}/2},
 \quad t\le\eta,\quad 0<\delta\le1,\\
 b_a\le2b',\quad w_P+\eta_3\le9b',\quad G_P\le10b',\quad
 \eta_3\le4\eta,\quad r\eta_3\le4m w'',\\
 z\le A\le x^{C+2\epsilon},\quad
 \frac rA\le\frac{2m}{z},\quad
 p_*\le b/z,\quad \log p_*\le z,\quad A<p_*,\\
 o\le x^{\epsilon/100},\quad 2\le L_w\le C\log x,
 \quad n\le o2^e,\quad \frac km\le R_0^2,
 \quad R_0\ge\sqrt{x^{1/2}\log x}.
\end{gathered}                                                    \tag{DNE.20}
\]
Every finite polynomial with fixed coefficients in
$k,n,m,r,t,N,n_V,z,S,\eta,|TT|,\mathcal D,b/p_*$ is bounded by
\[
             Kx^a\exp(Bz^{1/4})                                \tag{DNE.21}
\]
for $K,a,B$ depending only on $C$ and that polynomial.  Such a
bound is $o(b)$, and its product with $b'$ is also $o(b)$.

\begin{proof}
The bounds on $\eta,\Delta$ are \texttt{logeta}, \texttt{log\allowbreak Del}.
The polynomial rank bounds and $|TT|\le2n x^{-\epsilon/8}\le2n$
give (DNE.21) by its displayed formula for $\mathcal D$ and
$b/p_*=\mathcal D/m$.
Since $z^{1/4}=x^{\epsilon/4}$ and
$z^{1/3}=x^{\epsilon/3}$, taking logarithms proves, for every
fixed $K,a,B$,
\[
 Kx^a e^{Bz^{1/4}-z}\longrightarrow0,
 \qquad Kx^{a+1}e^{Bz^{1/4}-z^{1/3}}\longrightarrow0.     \tag{DNE.22}
\]
This also proves $A<p_*$ uniformly: $A\mathcal D/(mb)$ tends to zero.

For the top weights, the \texttt{b\allowbreak T\_\allowbreak le} inequality and
\texttt{rho\_\allowbreak le} show $b_a\le2b'$.  If $L.\mathrm{LaG}<0$,
the correction is nonpositive; otherwise discard the nonnegative
second term of \texttt{rho\allowbreak P} and bound its first term by
$x^{-2}\le1$.  The definitions then give $w_B\le3b'+\eta$.
The inequalities \texttt{teta} and \texttt{k\allowbreak S\_\allowbreak le} give
$\eta,S\le b'$, while $t\le\eta$ and $\eta\ge1$ give
$\eta_3\le4\eta$.  Substitution proves the bounds on $w_P,G_P$.
Moreover \texttt{Neta} gives
$r\eta_3\le4N\eta\le4\delta m w''\le4m w''$.
The formulas in (DNE.2), $r\le2k$ and
$\delta\in[x^{-\epsilon},1]$ give the bounds on $A,r/A$.
Since $\mathcal D\ge kn^2z\ge mz$, $p_*\le b/z$.
The final line is \texttt{o\_\allowbreak tie}, \texttt{Lw\_\allowbreak le}, the original
frame identity $n\le o2^e$, and \texttt{radius\_\allowbreak ge\_\allowbreak sqrt\_\allowbreak ratio}
together with \texttt{k\_\allowbreak ge\_\allowbreak mx}.  This proves every assertion.
\end{proof}

Only finitely many comparisons of type (DNE.22) are used below.
Their thresholds can therefore be maximized before any local index,
coefficient place, prime set or member of the family is chosen.
% END DNE6

% BEGIN DNE7
\paragraph{DNE7: the small-order inert primes.}
\label{claim:DNE7}
Uniformly as in DNE6, every finite set $T$ of odd integers $p\ge3$
with $\operatorname{ord}_{\pm}(p)\le4L_w$ satisfies
\[
           r\sum_{p\in T}\frac{\log p}{p^2}\le4mz^2.     \tag{DNE.23}
\]
No distribution theorem for primes is needed.

\begin{proof}
The elementary dyadic calculation proved in LFS2 says that, with
$v=\lfloor\log_2(4L_w)\rfloor$ and
$M_0=2^{e+2-v}$, these integers are $\pm1\pmod{M_0}$.
Eventually $2^e/L_w\ge2R_0\ge R_0+1$, since
$n=R_0z\le o2^e$ and $2oL_w\le z$ by (DNE.20).
Since $R_0\ge3$ eventually, $2^e\ge6L_w>4L_w$,
so $v<e$.  Thus the residue calculation applies without truncation, and
$M_0\ge2^e/L_w\ge R_0+1$.

For $R\ge3$ and an integer $M_0\ge R+1$, a decreasing-function
comparison along each of the progressions $aM_0\pm1$ gives
\[
 \sum_{p\equiv\pm1\ (M_0),\ p\ge2}\frac{\log p}{p^2}
 \le\frac{2\log R}{R^2}
             +\frac{2(\log R+2)}{RM_0}.                 \tag{DNE.24}
\]
For completeness, in each progression the first term is at most
$\log R/R^2$, because its first admissible member is at least
$M_0-1\ge R$.  The remaining terms, separated by $M_0$, are
bounded by $M_0^{-1}\int_R^\infty\log u/u^2\,du$; the latter
is $(\log R+1)/(RM_0)$, even stronger than the stated term.
Apply (DNE.24) at $R=R_0$.  Since $r\le2mR_0^2$, its first
contribution after multiplication by $r$ is at most $4m\log R_0$.
Since $n=R_0z\le oM_0L_w$, the second is at most
$4moL_w(\log R_0+2)/z$.
Now $\log R_0\le C\log x$, $o\le z^{1/100}$ and
$L_w\le C\log x$; thus
$\log R_0+oL_w(\log R_0+2)/z\le z^2$ eventually.
This proves (DNE.23) with one threshold for every finite $T$.
\end{proof}
% END DNE7

% BEGIN DNE8
\paragraph{DNE8: common error and split-place contributions below $p_*$.}
\label{claim:DNE8}
For every finite prime set $T\subset[5,p_*]$, uniformly eventually,
\[
 \sum_{p\in T}\log p\,B_0\le Cfm bz^2,
 \qquad \sum_{p\in T}\log p\,e_p\le11Cbz^2.             \tag{DNE.25}
\]

\begin{proof}
For any finite integer set below $Y\ge1$,
$\sum\log u\le Y\log Y$, $\#T\le Y$, and
$\sum\log u/u\le(1+\log Y)^2$; the last assertion follows from
$\sum_{u\le Y}1/u\le1+\log Y$.
The first bound in (DNE.25) follows from
\[
 B_0p_*z=Cfnt^2n_Vz^2\frac{mb}{\mathcal D}z
      \le Cfm bz\le Cfm bz^2,
\]
using $n_V\le N+t$.  DNE6 gives $H\le2b$, $\log H\le2z$,
$p_*\le b/z$, and $\log p_*\le z$, eventually with $z\ge1$.
Consequently
\[
 \sum_{p\in T}\log p\,e_p
 =C\left(\sum\log p+H\sum\frac{\log p}{p}
                                     +\#T\log H\right)
 \le C(b+8bz^2+2b)\le11Cbz^2.
\]
The total coefficient of $e_p$ from (DNE.17) is bounded by
$(3+2C)fm$; this keeps both the split term and the $C$ already
present in the inert term.
\end{proof}
% END DNE8

% BEGIN DNE9
\paragraph{DNE9: inert primes below $A$ with large dyadic order.}
\label{claim:DNE9}
For every finite set $T\subset[5,A)$,
\[
        \sum_{p\in T}4fCr\log p\,R(p)\le40Cfm bz^2       \tag{DNE.26}
\]
uniformly eventually.  It is permissible here to include all
integers in the indicated interval, so this bounds the needed
inert primes of large order.

\begin{proof}
By definition $(p-1)^{j_p}\ge A$, and
$j_p+1\le3+\log A$ for $p\ge5$.  Therefore
\[
 R(p)\le\frac H{A(p-1)}+(G_P+1)(3+\log A)+z.
\]
The elementary bound
$\sum_{p\le A}\log p/(p-1)\le2(1+\log A)^2$ and $r/A\le2m/z$
show that the sum of the first term, multiplied by $r$, is at
most $8mb(1+\log A)^2/z\le8mb$ eventually.
For the other terms use $G_P\le10b'$, $\log A\le(C+2)\log x$,
$A\le x^{C+2\epsilon}$ and $r\le2x^C$.
Their sum times $r$ is bounded by a fixed polynomial in $x$
times $b'$, plus a fixed polynomial in $x$.  DNE6 makes this
at most $mb$ eventually, uniformly over $T$.
Multiplication by $4fC$ gives $36Cfm b$, hence (DNE.26).
\end{proof}
% END DNE9

% BEGIN DNE10
\paragraph{DNE10: the remaining inert terms $H/p^2$.}
\label{claim:DNE10}
For the inert primes in the third and fourth rows of (DNE.12),
their combined $H/p^2$ contribution satisfies
\[
 \sum 4fCrH\frac{\log p}{p^2}\le96Cfm bz^2              \tag{DNE.27}
\]
uniformly eventually.

\begin{proof}
For $p\ge A\ge2$, monotonicity and integration of $\log u/u^2$
give the convenient bound
$\sum_{p\ge A}\log p/p^2\le4(\log A+2)/A$.
With $r/A\le2m/z$, $H\le2b$ and
$\log A+2\le z$, this contributes at most $64Cfm b$.
For $p<A$ and order at most $4L_w$, DNE7 and $H\le2b$ give
$4fC(2b)(4mz^2)=32Cfm bz^2$.  Add these bounds and use $z\ge1$.
\end{proof}
% END DNE10

% BEGIN DNE11
\paragraph{DNE11: the band contribution and the full flat term.}
\label{claim:DNE11}
For every finite prime set $T\subset[5,p_*]$, uniformly eventually,
\[
 \sum_{p\in T}4fCr\log p\,F(p)
 \le\{248C+16C(2C+1)^2\}fm bz^2.                      \tag{DNE.28}
\]

\begin{proof}
The contribution of $(w_P+\eta_3)/p$ before multiplication by
$4fC$ is at most
$9rb'(1+\log p_*)^2\le36rb'z^2\le36mbz^2$.
The last inequality follows uniformly from
$r\le2x^C$ and $b'/b=xe^{-z^{1/3}}$.
The contribution of $S+z$ is bounded by
\[
 r(S+z)p_*z
   \le2kn^2(S+t+z)\frac{mb}{\mathcal D}z
   \le2mbz\le2mbz^2.                                  \tag{DNE.29}
\]

For the band, DNE6 permits one threshold such that
$C(\eta_3+S+z)\le w''/4$.  Suppose $I(p)=1$, witnessed by
$s\ge w''$, and $p\ge2C$.  Then the right side of (DNE.13)
is at most $ps/4+Cb/p+ps/2$.  Since $Cb/p\le b/2$,
\[
 ps-b\le3ps/4+b/2,\qquad ps\le6b,
                         \qquad p\le6b/w''.             \tag{DNE.30}
\]
There are at most $6b/w''$ such integers.  Their logarithms are
at most $z$ because $p\le p_*$.  Thus
$r\eta_3\sum\log p\le(4mw'')(6b/w'')z=24mbz$.
The primes $p<2C$ contribute at most
$4mw''(2C+1)^2\le4(2C+1)^2mb$, using
$\sum_{p<2C}\log p\le(2C+1)^2$.
Together the three contributions before $4fC$ are bounded by
$\{62+4(2C+1)^2\}mbz^2$, proving (DNE.28).
\end{proof}

This argument uses the actual band implication.  It does not
replace the band set by arbitrary primes or discard its dependence
on the local row and place.
% END DNE11

% BEGIN DNE12
\paragraph{DNE12: large exceptional primes and the three special classes.}
\label{claim:DNE12}
Set $X_0=fm bz^2$.  Uniformly eventually, the part of
$\sum\log p\,\Lambda(p)$ from the large exceptional primes is
at most $X_0$, the generator-prime part at most $X_0$, the
small-prime part at most $2X_0$, and the $q,d$ part at most
$4CX_0$.  These bounds apply to arbitrary finite subsets of each class.

\begin{proof}
For the exceptions $p>p_*$, (DNE.14) and DNE6 bound
$|\mathcal E|$ by an expression of form (DNE.21), while
$p\le3b$ gives $\log p\le2z$.  Also $p>p_*>A$ and
$\log p\ge z/2$ eventually.  Here $H/p$ is bounded by
$2\mathcal D/m$, and $\log H/\log p\le4$.
Thus $e_p$ and $B_0$ satisfy (DNE.21).  In the inert branch,
$H/p^2$ does too, and
\[
 F(p)\le9b'/p_*+S+z+\eta_3
       =9x\mathcal D/(mX')+S+z+\eta_3
\]
is again of that form, even with the crude bound $I(p)\le1$.
The high-order branch below $A$ cannot occur.  Therefore the
entire exceptional sum is of form (DNE.21), hence at most
$b\le X_0$ eventually.

There are at most two generator primes, and \texttt{g\_\allowbreak Y},
$P.CYL\le C$ give $\log g_i\le Cz+\log2$.  Their cost is
$Cx^C\eta$ each, so their total is again of form (DNE.21)
and eventually at most $b$.
The small primes divide $Q_S$, whence
\[
 \sum_{p\in\mathrm{small}\cup\{3\}}\log p
      \le\log Q_S+\log3\le2C\log x+\log3.
\]
Multiplication by $Cfm b\log x$ is at most $X_0$ eventually,
because $(\log x)^2=o(z^2)$.  Multiplication by $Cx^Cz^2$
is at most another $b$, giving $2X_0$.
Finally $\log q,\log d\le z$ follows from
$M^{36}t^{150}\le b^{1/2}$.  At most two terms of size
$2Cfm b\log p$ give $4Cfm bz\le4CX_0$.
\end{proof}
% END DNE12

% BEGIN DNE13
\paragraph{DNE13: the complete original denominator inequality.}
\label{claim:DNE13}
For fixed $\epsilon,C$ and every sufficiently large $x$ as above,
any original admissible table gives, for every original lift $j$,
\[
 \eta\mathrm{Sum}(L,j)
 \le d_j\left[2000(C+1)^3
       \{fm bz^2+(\log x)^2L.Nc\}
                       +\sum_p\mathrm{vvTerm}(D,p)\right]        \tag{DNE.31}
\]
whenever the last series is summable.  This is exactly
\texttt{Link\allowbreak N\_\allowbreak at}$(L,j,2000(C+1)^3)$.

\begin{proof}
Partition the ordered table (DNE.18) into its special classes,
ordinary primes below $p_*$, and exceptional primes above $p_*$.
One can also bound by the sum of the five overlapping nonnegative
majorants, so overlaps at special primes present no issue.
For ordinary primes use DNE8--DNE11.  The bound, in units of
$X_0=fm bz^2$, is
\[
 C+11C(3+2C)+40C+96C+248C+16C(2C+1)^2
                 =434C+86C^2+64C^3.                    \tag{DNE.32}
\]
DNE12 adds at most $4+4C$.
The non-prime part of (DNE.19) is at most $b\le X_0$
eventually: $\log M\le z/72$, $\log(tr)\le O_C(\log x)$,
$\log h_E\le Cx^C$ and $\log J\le Cx^Cz$.
Thus its inclusion changes the majorant to
$5+438C+86C^2+64C^3$, which is at most $2000(C+1)^3$ for $C>0$.
Substitute in (DNE.19).
The source quantity $L.Nc$ is nonnegative, by its definition
as a quotient of positive rank/scale factors and the nonnegative
thick-place count.  Adding $2000(C+1)^3(\log x)^2L.Nc$ therefore
gives exactly (DNE.31).  All multipliers use the same $d_j\ge0$.
\end{proof}

The small-prime Selmer finiteness guard is not removed: it enters
when \texttt{admitted} supplies the table under
\texttt{Sha\allowbreak Fin\allowbreak Odd}.  The summability guard is used at precisely
the nonnegative infinite-sum step in DNE5.
% END DNE13

% BEGIN DNE14
\paragraph{DNE14: uniform export and the original common-family consumer.}
\label{claim:DNE14}
For every fixed $\epsilon$ and real $C_a$ one may choose
\[
                  C_{\mathrm{out}}=2000(|C_a|+1)^3             \tag{DNE.33}
\]
before the family $\mathcal F:\texttt{Fam\allowbreak P}(\epsilon)$, so that
\texttt{Adm\allowbreak FC}$(\mathcal F,C_a)$ implies
\texttt{Family\allowbreak Thin\allowbreak Scales\allowbreak Bound}$(\mathcal F,C_{\mathrm{out}})$.
Consequently
\[
 \texttt{Concl\allowbreak F}(\mathcal F)\ \Longrightarrow\ 
 \texttt{Adm\allowbreak F}(\mathcal F)\ \Longrightarrow\ 
 \texttt{Link\allowbreak N}(\mathcal F).                              \tag{DNE.34}
\]

\begin{proof}
If a family exists, its scale record at its supplied member at
$x=0$ fixes $0<\epsilon<1/52$; the family parameter $x$ and
the scale coordinate are identified at all sufficiently late
indices by its original equality.  Take the cutoff in
\texttt{Adm\allowbreak FC} and maximize it with $\mathcal F.x1+1$ and the
uniform cutoff of DNE6--DNE13.
At a later member satisfying \texttt{Sha\allowbreak Fin\allowbreak Odd}, \texttt{admitted}
supplies an actual table with all fields of \texttt{Adm\allowbreak T}.
Its parameter bound forces $C_a\ge P.CLw>0$, so DNE13 applies
with $C=C_a$ and (DNE.33).  It is universal over the lift $j$,
and still has the original summability antecedent.
If $C_a\le0$, no such admissible table exists, so the same guarded
conclusion is vacuous; if there is no family the assertion is
also vacuous.  No choice of table is needed before the
\texttt{Sha\allowbreak Fin\allowbreak Odd} premise.

This proves the full original \texttt{Uniform\allowbreak Export\allowbreak Statement}:
$C_{\mathrm{out}}$ is independent of the family, and only the
final family threshold uses its initial cutoff.
The identity between the exported and original $\eta\mathrm{Sum}$
is definitional: both use the same twisted support, weights and
$\eta_J$.  This gives \texttt{Link\allowbreak N} with precisely its original
quantifier order.  Finally, \texttt{Concl\allowbreak F} supplies one $C_a$
and one threshold, then an actual $X$ under \texttt{Sha\allowbreak Fin\allowbreak Odd}
for each later member.  Apply DNE3--DNE4 to those very data;
this supplies \texttt{Adm\allowbreak FC} and proves (DNE.34).
\end{proof}

% END DNE14

\subsection{What remains upstream}

The construction of the actual common $X$ satisfying
\texttt{Concl} has not been assumed away.  In particular the
\texttt{Norm\allowbreak Choice\allowbreak Facts}, compatible Galois points and killing data,
the numerical row bounds on their actual lattices, and the original
two/three-prime and type-prime bounds must still be supplied on
the selected family.  DNE3--DNE14 now discharge the entire
table-patching, weighted aggregation and uniform export after those
inputs, and do so with the unchanged coefficient field and host.
They do not prove the closed selected-class supplier of these inputs,
or nonemptiness of the common construction class.

\paragraph{Source map.}
The exact pinned source ranges, full original contracts and file hashes
are recorded in \texttt{denominator\_\allowbreak export\_\allowbreak v50/\allowbreak provenance.\allowbreak json},
\texttt{contracts.\allowbreak json} and \texttt{source\_\allowbreak manifest.\allowbreak json}.
The main sources are \texttt{Bound\allowbreak At\allowbreak Two\allowbreak Three\allowbreak Family.\allowbreak lean},
\texttt{Admitted\allowbreak Statements\allowbreak Of\allowbreak Local\allowbreak Bounds\allowbreak Conclusion.\allowbreak lean}, its seven
table lemmas, \texttt{Twisted\allowbreak Places.\allowbreak lean},
\texttt{Mem\allowbreak Primes\allowbreak Over\allowbreak Of\allowbreak Mem.\allowbreak lean}, \texttt{NEq\allowbreak RZero\allowbreak Mul\allowbreak XPow.\allowbreak lean},
\texttt{Export\allowbreak Bound\allowbreak Of\allowbreak Admissible.\allowbreak lean},
\texttt{Theta\allowbreak Scalar\allowbreak Family\allowbreak Of\allowbreak Eigenvalue.\allowbreak lean} and
\texttt{Link\allowbreak NAlong\allowbreak Family\allowbreak Of\allowbreak Local\allowbreak Bounds.\allowbreak lean}.
The sole inherited mathematical argument used as a theorem is the
dyadic residue calculation LFS2; its range and hash are included.
All analytic comparisons used in the simplification are proved in
DNE6--DNE13.  No new external published theorem is invoked.

\endgroup
\end{document}
